% 本文件由 LaTeX 排版 Agent 根据有效 translation.json 自动生成。
% author=John Cassian (c. 360-c. 435); work_id=3507; title=会规

\chapter{\WorkTitle{会规}}

% structure: p0001 source=h1 target=omitted reason=page-title-already-rendered-by-parent

\Person{若望·卡西安}论隐修院会规及八种主要过犯之补救十二卷\par

第一卷\newline{}第二卷\newline{}第三卷\newline{}第四卷\newline{}第五卷\newline{}第七卷\newline{}第八卷\newline{}第九卷\newline{}第十卷\newline{}第十一卷\newline{}第十二卷\par

\SourceRecord{Translated by C.S. Gibson.；Nicene and Post-Nicene Fathers, Second Series；Vol. 11.；Edited by Philip Schaff and Henry Wace.；Buffalo, NY: Christian Literature Publishing Co.,；1894.；<http://www.newadvent.org/fathers/3507.htm>.}

% structure: p0001 source=h1 target=section reason=page-title

\section{\WorkTitle{隐修院制度（卷一）}}

\WorkTitle{若望·加西安的隐修院制度十二卷以及八种主要过失的补救}\par

% structure: p0003 source=h2 target=subsection reason=relative-heading-rank; base=h2

\subsection{第一章}

\Emph{论隐修士的腰带。}\par

既然我们将要讲述隐修院的习俗与规则，那么，赖天主的恩宠，我们如何能比从隐修士的实际服装开始更好呢？因为当我们把他们的外在之人摆在你们眼前之后，我们便能适时地阐述他们的内在生活。因此，隐修士作为基督的士兵，时刻准备战斗，应当常束着腰行走。因为圣经的权威也表明，在旧约中开创这种生活之始的那些人——我是指\Person{厄里亚}和\Person{厄里叟}——就是这样行走的；此外，我们知道新约的领袖和作者，即\Person{若翰}、\Person{伯多禄}和\Person{保禄}，以及同等的其他人，也都同样行走。其中首先提到的这位，即使在旧约中就展现了童贞生活的花朵和贞洁与节制的榜样，当他被主派遣去责备\Person{阿哈齐雅}、\Place{以色列}恶王的使者时——因为这位王因患病而躺在床，竟打算去求问\Place{厄刻龙}的神贝尔则步关于他的健康状况——于是这位先知遇见他们，说他不应从所躺的床上下来；而这位卧床的国王通过对其服装特征的描述认出了他。因为当使者回到他那里，带回了先知的信息，他问那遇见他们并说这些话的人是什么样子，穿着如何。他们说：\BibleQuote{一个身穿毛衣的人，腰束皮带。}国王一见这服饰，立刻知道那是天主的人，就说：\BibleQuote{这是提协贝人厄里亚。}\BibleRef{列下 1:1}也就是说，凭借腰带和那毛茸茸、不修边幅的身体外观，他毫不犹豫地认出了天主的人，因为当他居住在成千上万的\Place{以色列}人中时，这装饰总是附在他身上，仿佛印刻为他个人特有风格的特殊标志。关于\Person{若翰}，他作为新旧约之间某种神圣的边界来临，既是开端也是终结，我们通过福音作者的见证知道：\BibleQuote{这若翰身穿骆驼毛的衣服，腰束皮带。}\BibleRef{玛 3:4}当\Person{伯多禄}被\Person{黑落德}下在监里，第二天将要被提出处死时，天使站在他旁边，吩咐他说：\BibleQuote{束上带子，穿上鞋。}\BibleRef{宗 12:8}主的天使之所以吩咐他这样做，必定是因为看见他为了一夜休息，暂时解开了通常系在腰间的带子，让疲惫的四肢松开。同样，\Person{保禄}上\Place{耶路撒冷}去，即将被犹太人捆锁时，在\Place{凯撒勒雅}遇见先知\Person{阿加波}，他拿了保禄的腰带，捆上自己的手脚，以身体动作表明他要受的伤害，并说：\BibleQuote{耶路撒冷的犹太人将要这样捆绑这根腰带的主人，把他交在外邦人手里。}\BibleRef{宗 21:11}先知绝不会提出这点，也不会说\BibleQuote{这根腰带的主人}，除非保禄素来习惯把腰带系在腰间。\par

% structure: p0006 source=h2 target=subsection reason=relative-heading-rank; base=h2

\subsection{第二章}

\Emph{论隐修士的长袍。}\par

也让隐修士的衣袍只能恰好遮体，防止裸露之羞，并抵御寒冷之伤，而不是助长虚荣和骄傲的种子；因为同一宗徒告诉我们：\Quote{只要有食物和遮体之物，就当知足。}他说\Quote{遮体之物}，而非\Quote{衣服}，正如某些拉丁抄本错误记载的那样；即仅能遮盖身体之物，而非以华丽服饰取悦幻想之物；要平凡，使其在同行之人中不因颜色或款式的奇特而引人注目；且完全摆脱焦虑的操心，但也不因疏忽而沾染污渍。最后，它们应如此远离世俗时尚，以至于完全成为天主仆人的共同财产。因为凡被一个或少数天主仆人声称拥有、而非全体弟兄共同所有的，要么是多余要么是虚荣，因此应视为有害，并显出虚荣而非德行的外表。因此，我们所见到的任何模式，无论是古代奠定隐修生活基础的圣人，还是我们当代继承并保持其习俗的教父们所教导的，我们也应拒绝，视为多余和无用：因此他们全然不赞成粗毛布衣袍，因为它显眼而引人注目，且正因如此，不仅对灵魂无益，反而助长虚荣和骄傲，而且不便且不适合履行必要的工作——为这些工作隐修士应随时准备好并无障碍。但即使我们听说有些可敬之人曾穿着这种服装，我们也不应因此为隐修院制定规则，也不应推翻古代圣教父的法令，因为我们不认为少数人仗着拥有其他德行，即使在他们并非按公教规则所行之事上也应受责备。因为少数人的意见不应优先于或干预普遍的规则。我们应当毫不犹豫地忠诚和不加质询地服从的，不是那些少数人意志引入的习惯和规则，而是那些长久以来的古代和众多圣教父一致决议传给后代的。的确，约兰，以色列邪恶的君王，被敌军包围时撕裂衣服，据说里面穿着粗毛布衣，这件事不应作为我们日常生活的先例；\BibleRef{列下 6:30}或者尼尼微人为了减轻天主借先知向他们宣布的判决，身穿粗毛布；\BibleRef{纳 3:8}前者被显示秘密地穿在里面，因此除非外衣撕裂，否则无人能知道；后者忍受粗毛布覆盖，那时全城都为即将来临的毁灭哀悼，都穿着同样的衣服，无人能指摘为炫耀。因为若没有特别区别，大家都一样，便没有害处。\par

% structure: p0009 source=h2 target=subsection reason=relative-heading-rank; base=h2

\subsection{第三章。}

\Emph{论埃及人的风帽。}\par

此外，在埃及人的服饰中，有些方面与其说关乎身体的照料，不如说关乎品格的规范，以便借服装本身的特征保持纯朴和天真的遵守。因为他们昼夜常使用极小的风帽，垂到颈项和肩膀的末端，仅盖住头部，为使他们通过模仿孩童的实际服装，时常被提醒要保持孩童的纯朴和天真。这些人在基督里回归童年，时时刻刻用心和灵魂歌唱：\Quote{上主，我的心不骄矜，我的眼不高傲；我没有从事伟大及超越我的事。我若不一味谦卑，却高举我的灵魂，就如断奶的婴孩在他母亲怀中。}\par

% structure: p0012 source=h2 target=subsection reason=relative-heading-rank; base=h2

\subsection{第四章。}

\Emph{论埃及人的束腰外衣。}\par

他们还穿亚麻束腰外衣，几乎只到肘部，其余部分露出手臂，如此衣袖的截断暗示他们已斩断所有世俗的行为和事务，而亚麻服装教导他们已死于一切世俗的交往，并借此日日听见宗徒对他们说：\Quote{致死你们在地上的肢体；}他们的服装本身也在宣告：\Quote{因为你们已经死了，你们的生命与基督一同藏在天主内；}又说：\Quote{我活着，但不再是我，而是基督在我内活着。对我而言，世界已被钉在十字架上，我也被钉于世界。}\par

% structure: p0015 source=h2 target=subsection reason=relative-heading-rank; base=h2

\subsection{第五章。}

\Emph{论他们的绳索。}\par

他们还佩戴用羊毛线编织的双肩带，希腊人称为\OriginalTerm{肩带}{\GreekText{ἀνάλαβοι}}，而我们可以称之为束带或系绳，或更恰当地称为绳索。这些肩带从颈项顶部垂下，在喉咙两侧分开，环绕腋窝处的衣褶，并将它们从两侧收拢，以便他们可以将宽大的衣褶拉紧并贴身束起，如此束上臂膀，变得敏捷而准备好从事各种工作，竭力遵行宗徒的训诲：\Quote{这些手不仅供应了我，也供应了与我同在的人；} \Quote{我们没有白吃任何人的饭，而是昼夜辛苦劳作，免得成为你们任何人的负担。}又说：\Quote{若有人不愿工作，也不该吃饭。}\par

% structure: p0018 source=h2 target=subsection reason=relative-heading-rank; base=h2

\subsection{第六章。}

\Emph{论他们的披肩。}\par

接着他们用一条窄披肩覆盖颈项和肩膀，既追求服装的朴素，也注重便宜和经济；这在我们的语言和他们的语言中均称为\OriginalTerm{马福尔}{mafors}；如此他们便避免了斗篷和大衣的开销与炫耀。\par

% structure: p0021 source=h2 target=subsection reason=relative-heading-rank; base=h2

\subsection{第七章。}

\Emph{论绵羊皮和山羊皮。}\par

他们服装的最后一件是山羊皮衣，称为\OriginalTerm{默洛特}{melotes}或\OriginalTerm{佩拉}{pera}，以及一根手杖。他们携带此物，是为效仿那些在旧约中预示隐修生活的人，关于这些人，宗徒说：\BibleQuote{他们披着绵羊皮和山羊皮，忍受贫乏、忧苦和折磨；世界配不上他们；他们在旷野、山岭、山洞和地穴中飘流。}\BibleRef{希 11:37}而这件山羊皮衣的意义是：既已摧毁肉身情欲的一切放荡，他们应当持续处于最严格的品德节制中，且在他们身上不应再留存任何年少的放荡或热狂，或他们旧时的轻浮。\par

% structure: p0024 source=h2 target=subsection reason=relative-heading-rank; base=h2

\subsection{第八章。}

\Emph{论埃及人的手杖}\par

因为\Person{厄里叟}本人就是他们中的一位，教导说这些人习惯携带一根手杖；他对他的仆人\Person{革哈齐}说，当派他去复活那妇人的儿子时：\BibleQuote{拿我的手杖，跑去放在孩子的脸上，使他可以活过来。}\BibleRef{列下 4:29}先知绝不会把杖给他拿着，除非他惯于经常随身携带。而携带手杖在灵性上教导他们：绝不可赤手空拳地行走在如此之多的缺点吠犬和无形灵性邪恶的野兽中间（有福的\Person{达味}渴望脱离这些，遂说：\Quote{主，不要将信赖祢的灵魂交给野兽}），但当它们攻击他们时，他们应用十字架的记号击退它们，并将它们远远驱赶；当他们猛烈地对他们发怒时，他们应借不断思念主的苦难和效法祂克制生活的榜样来消灭它们。\par

% structure: p0027 source=h2 target=subsection reason=relative-heading-rank; base=h2

\subsection{第九章。}

\Emph{论他们的鞋}。\par

他们拒绝鞋子，因为福音的命令禁止；但如果身体虚弱，或冬天的晨寒，或中午的酷热迫使，他们仅用凉鞋保护双脚，解释说：使用凉鞋和主的许可暗示，如果我们仍在此世，不能完全摆脱对肉身的忧虑和焦虑，也不能完全从中解脱，我们至少要以尽可能少的忙乱和纠缠来供应身体的需要。至于我们灵魂的双脚——它们应当为我们灵性的赛跑做好准备，并时刻准备宣讲福音的平安（我们用这双脚追随基督香膏的芬芳，\Person{达味}说：\Quote{我渴切奔跑}，\Person{耶肋米亚}说：\Quote{但我跟随祢，并不困乏}），我们不应让它们被此世致命的忧虑缠住，使我们的思想充满那些不关供应自然需要，而关不必要的和有害的享乐之事。如果按照宗徒的劝告，我们\BibleQuote{不为肉身作准备，去放纵私欲}\BibleRef{罗 13:14}，我们就能这样做到。但尽管他们合法地使用这些凉鞋，如主的命令所允许，他们却绝不容它们留在脚上，当他们走近举行或领受神圣奥迹时，因为他们认为应当按字面遵守对\Person{梅瑟}和农的儿子\Person{若苏厄}所说的话：\Quote{解开你的鞋带，因为你站的地方是圣地。}\par

% structure: p0030 source=h2 target=subsection reason=relative-heading-rank; base=h2

\subsection{第十章。}

\Emph{论根据气候特征或当地习俗可能允许对惯例所作的修改。}\par

说这么多，是为了显得我们没有遗漏埃及人服装的任何一件。但我们只需保留那些地方环境和当地习俗所允许的。因为冬季的严寒不容我们满足于拖鞋、短衣或一件单衣；而戴小帽或披羊皮会引起嘲笑，而非令观者受益。因此，我们认为只应引入我们上面描述的那些，且适合我们修会谦卑的品格和气候的性质，使我们服装的主要特点不是新颖的款式（那可能得罪世俗之人），而是其可敬的简朴。\par

% structure: p0033 source=h2 target=subsection reason=relative-heading-rank; base=h2

\subsection{第十一章。}

\Emph{论灵性腰带及其奥秘意义。}\par

因此，基督的士兵穿上这些服装，首先应当知道，他因腰间所系的\OriginalTerm{腰带}{girdle}而受到保护，这腰带不仅使他在心里准备好为隐修院的一切工作和事务服务，而且还使他始终行动不受衣饰的阻碍。因为他对服从和工作的热诚越真诚，他在心灵的纯洁、圣德的进步和对神圣事物的认识上就越热心。其次，他应当意识到，实际佩戴腰带这件事本身并非小奥秘，它指示着对他的要求。因为用死皮束腰并绑住腰部，意味着他背负着对这些肢体的克制，这些肢体中含有邪淫和好色的种子，始终知道福音的命令，即\BibleQuote{你们腰间要束上带}\BibleRef{路 12:35}，借着宗徒的解释适用于他，即：\BibleQuote{所以，你们要致死属于地上的肢体，即邪淫、污秽、邪情、恶欲}\BibleRef{哥 3:5}。因此我们在圣经中发现，只有那些体内肉欲的种子已被消灭的人，才被束上腰带，他们竭力高唱圣\Person{达味}的这句话：\BibleQuote{因为我好像成了霜中的皮囊}，因为当有罪的肉身在内心被摧毁时，他们能借圣神的威力扩张外在之人的死皮。因此他特意加上\BibleQuote{在霜中}，因为他们不仅满足于内心的克制，而且连外在之人的活动和本能的冲动，也被来自外部的节制的寒霜所冻结，只要他们如同宗徒所说，不再让罪恶在他们可死的身体中掌权，也不再穿戴反抗圣神的肉身。\par

\SourceRecord{Translated by C.S. Gibson.；Nicene and Post-Nicene Fathers, Second Series；Vol. 11.；Edited by Philip Schaff and Henry Wace.；Buffalo, NY: Christian Literature Publishing Co.,；1894.；<http://www.newadvent.org/fathers/350701.htm>.}

% structure: p0001 source=h1 target=section reason=page-title

\section{\WorkTitle{修院规条（卷二）}}

% structure: p0002 source=h2 target=subsection reason=relative-heading-rank; base=h2

\subsection{第一章}

\Emph{论夜间祈祷与圣咏的正典制度}\par

因此，用我们所说的这一双重腰带束腰后，基督的士兵应接着学习东方圣教父们早已安排的夜间祈祷与圣咏的正典制度。但关于它们的性质，以及我们如何能如宗徒所指示的那样\Quote{不断祈祷}（\BibleRef{得前 5:17}），我们将在适当的时候，当主允许我们开始讲述\WorkTitle{长老会议录}时，予以论述。\par

% structure: p0005 source=h2 target=subsection reason=relative-heading-rank; base=h2

\subsection{第二章}

\Emph{论各地区规定咏唱圣咏数量之差异}\par

因为我们发现，许多人在不同地区，按照自己内心的想象（诚然，如宗徒所说：\Quote{对天主有热心，但不按真知识}（\BibleRef{罗 10:2}）），在这件事上为自己制定了不同的规则和安排。有些人规定每夜应诵念二十或三十篇圣咏，并通过对唱咏唱的音乐以及增添一些变调来延长它们。还有人甚至试图超过这个数目。有些人用十八篇。这样，我们发现不同地方有不同的规则，而我们所见到的制度和规定几乎与我们访问过的隐修院和单人小室的数目一样多。也有人认为，在白天的祈祷时辰，即第三时辰、第六时辰和第九时辰，圣咏和祈祷的数目应与向主献上礼仪的时辰数目完全对应。有人认为日间每项礼仪应指定六篇圣咏。因此，我认为最好陈述教父们最古老的制度，它至今仍被整个埃及的天主仆人们遵守，以便你们新建的隐修院在基督里的未成熟幼年期，能受到最早教父们最古老制度的教导。\par

% structure: p0008 source=h2 target=subsection reason=relative-heading-rank; base=h2

\subsection{第三章}

\Emph{论全埃及遵守统一规则及选立管理弟兄者}\par

因此，在全埃及和泰拜德，隐修院并非按每个弃绝俗世之人的臆想所建立，而是藉由教父的传承及其传统延续至今，或者以此方式建立并延续；我们在这些隐修院中注意到，在晚间聚会和夜间守夜时，遵守着一套规定的祈祷制度。因为若一个人尚未不仅舍弃自己的一切财产，而且还学会了自己不是自己的主人、对自己的行为没有主权这一事实，他是不允许主持弟兄的聚会，甚至不允许管理自己的。因为弃绝俗世之人，无论他拥有多少财产或财富，必须寻求集体隐修院的共同居所，以免他以自己已舍弃或带入隐修院的任何东西而自夸。他还必须服从众人，以学会他必须如主所说（\BibleRef{玛 18:3}）再变成小孩子，不因自己的年龄和如今算为虚度而在世俗中浪费的年岁而自恃，并且作为初学者，由于他所知自己在基督的服役中所服务的新学徒身份，他应当毫不迟疑地甚至服从比他年轻的人。此外，他必须习惯于劳动和辛劳，以便按照宗徒的命令（\BibleRef{得前 4:11}）亲手准备每日的饮食，供自己使用或供应陌生人的需要；并且他也可能忘记过去生活的骄傲和奢侈，通过辛苦的劳动获得内心谦逊。因此，没有人被选来管理一群弟兄，除非那要掌权的人已通过服从学会了该命令那些将要服从他的人什么，并从长老们的规则中发现了该教导他晚辈什么。因为他们说，善于管理或善于被管理需要一个有智慧的人，他们称这是圣神最大的恩赐和恩宠，因为除了那先前已在一切美德规则中受过训练的人，没有人能向服从他的人颁布有益的训诫；也没有人能服从一位长老，除非他充满了对天主的爱并完善了谦逊之德。因此，我们看到其他地区有各种不同的规则和规定在使用，因为我们常常甚至没有学习长老们的制度就胆敢管理隐修院，并把自己指定为院长，而尚未按应有的方式自称门徒，并且更愿意要求遵守自己的发明，而不是保留我们前人久经考验的教导。但是，当我们本意要解释应遵守的最佳祈祷制度时，由于急切地谈论教父的规条，我们通过一次仓促的离题，提前讲述了我们原要留到适当地方才说的内容。所以现在让我们回到当前的论题。\par

% structure: p0011 source=h2 target=subsection reason=relative-heading-rank; base=h2

\subsection{第四章}

\Emph{论全埃及和泰拜德如何规定圣咏为十二篇}\par

因此，正如我们所说，在全埃及和泰拜德，无论是在晚祷还是在夜间礼仪中，圣咏的数目固定为十二篇，结束时有两篇读经，一篇来自旧约，一篇来自新约。这一安排，早在很久以前就已经确定，并且在如此多的时代中，至今在所有那些地区的隐修院中延续不断，因为据说它不是出于人的发明，而是借着天使的职务从天上传给教父们的。\par

% structure: p0014 source=h2 target=subsection reason=relative-heading-rank; base=h2

\subsection{第五章}

\Emph{论圣咏定为十二篇的事实如何来自天使的教导}\par

因为在信仰初期，只有少数最优秀的人被称为隐修士；他们秉承真福的记忆者 \Person{马尔谷}——首位以主教身份管理亚历山太教会——从这位福音作者接受了那种生活方式，不仅保存了那些伟大的特征（我们在 \BibleBook{宗徒}{大事录} 中读到初期教会和信徒群众以此闻名：\BibleQuote{信徒群众都同心合意，没有人说自己的财物中有哪一样是自己的；相反，一切都是公用的。凡拥有田产房屋的，都卖了，把变卖的价钱拿来，放在宗徒脚前，然后照每人的需要分配。}），\BibleRef{宗 4:32} 而且还添加了其他更崇高的特征。因为他们退居城外更僻静的地方，过着如此严格的斋戒生活，以致即使异教之人也对那崇高生命的生活惊叹不已。因为他们专心于研读圣经、祈祷和日夜手工劳动，如此热诚，以致对食物没有欲望或念头——除非到第二或第三天，身体的饥饿提醒他们；他们进食和饮水并非出于喜好，而是生活所必需；即使如此，他们也只在日落之后才进食，以便将白昼的时光用于灵修默想，将身体的照料留给夜晚，并且还从事比这些更高尚的事。关于这些事，凡是从未从熟悉此类事物的人那里听到者，都可以从教会历史中学到。因此，当初期教会的完美仍未被破坏，仍由其追随者和继承者新鲜地保存着，并且少数人的热诚信仰尚未因分散到多数人中而变得冷淡之时，可敬的教父们以警醒的关怀为后来者作出安排，聚在一起讨论应如何为全体弟兄的日常祈祷制定计划；以便将一份毫无争议和分裂的虔诚与和平遗产传给后继者，因为他们担心在日课方面，那些在同一敬拜中结合的人之间可能产生差异或争论，并且有时会在后来者中产生错误、嫉妒或分裂的有毒根苗。由于每人按自己的热诚（而忽略了别人的软弱），认为那按自己的信仰和力量看来很容易的事就应该这样规定，很少考虑通常对大多数弟兄（其中必定有很大比例的软弱者）是否可能；他们以不同程度按各自能力竭力确定圣咏的极大数量，有的定为五十首，有的定为六十首，有的还不满足这数量，认为实际应该超过——就在他们对教规这虔诚讨论中发生了如此神圣的意见分歧，以至于晚祷的时间在神圣决议之前就到了。当他们正要举行日常礼仪和祈祷时，有一个人站起来在中间向主咏唱圣咏。他们全都坐着（正如埃及仍保留的习惯那样），心思专注地听着歌咏者的言辞，当他唱了十一首圣咏，其间穿插着祈祷，一句接一句平均地诵出，他以阿肋路亚的应答结束了第十二首；然后他突然从众人眼前消失，同时结束了他们的讨论和礼仪。\par

% structure: p0017 source=h2 target=subsection reason=relative-heading-rank; base=h2

\subsection{第六章。}

\Emph{论十二次祈祷的习俗。}\par

于是，可敬的教父们领悟到，藉着天使的指引，上智安排为弟兄们的集会定下了普遍规则，因此决定在晚祷和夜祷中都保留这一数目。他们又在这两篇读经之外加上了两篇——一篇选自 \WorkTitle{旧约}，另一篇选自 \WorkTitle{新约}——但只是作为附加的、出于他们自己决定的，仅给那些愿意并渴望通过恒常研习而拥有饱含圣经智慧的心志者。但在星期六和星期日，他们则从 \WorkTitle{新约} 中诵读两篇：一篇选自 \WorkTitle{书信} 或 \BibleBook{宗徒}{大事录}，一篇选自 \WorkTitle{福音}。那些负责诵读和记忆圣经的人，从复活节到圣神降临节，也如此行。\par

% structure: p0020 source=h2 target=subsection reason=relative-heading-rank; base=h2

\subsection{第七章。}

\Emph{论他们的祈祷方法。}\par

这些上述的祈祷，他们开始和结束的方式是：当圣咏结束时，他们不会立刻匆忙跪下，不像我们这里的有些人，在圣咏尚未完全结束之前，就急忙俯伏在地祈祷，急于尽快结束祈祷仪式。因为虽然我们选择了超过古人所定的时限，补足剩余圣咏的数目，但我们急于结束祈祷仪式，考虑的是疲惫身体的休息，而非从祈祷中寻求益处和裨益。然而，在他们那里并非如此：他们在屈膝之前先祈祷片刻，站立时花大部分时间祈祷。然后，在极短的时间里，他们俯伏于地，仿佛只是朝拜天主的仁慈，然后尽快起身，再次伸开双手直立——正如他们先前站着祈祷一样——心思专注地停留在祈祷中。因为，他们说，如果你长时间俯伏在地，更容易受到攻击，不仅是游荡的思绪，还有瞌睡。但愿我们也不凭经验和日常实践知道这一点——我们这些俯伏在地时，常常希望这种姿势延长一段时间，与其说是为了祈祷，不如说是为了休息。但是，当那位要\Quote{集祷}的人从地上起身时，所有人都立刻站起，这样就没有人敢在\Emph{他}俯伏之前屈膝，也没有人在\Emph{他}起身之后延迟，以免被认为他独立献上了自己的祈祷，而非跟随主领者直到结束。\par

% structure: p0023 source=h2 target=subsection reason=relative-heading-rank; base=h2

\subsection{第八章}

\Emph{论圣咏之后的祈祷。}\par

我们在本国所观察到的这种做法——即当一人唱完圣咏时，所有的人都站起来一起大声唱\Quote{光荣归于圣父、圣子及圣神}——我们从未在东方任何地方听说过，但在那里，当圣咏结束时，所有人都保持静默，接下来由歌咏者献上祈祷。但只有这颂扬天主圣三的赞美诗才通常用来结束全部圣咏。\par

% structure: p0026 source=h2 target=subsection reason=relative-heading-rank; base=h2

\subsection{第九章}

\Emph{论祈祷的特征，其更完整的论述将留待\WorkTitle{长老会议}。}\par

既然这部\WorkTitle{院规}引导我们进入会规祈祷的体系，其更完整的论述将留待\WorkTitle{长老会议}（届时我们将更详细地讨论他们的祈祷特征以及如何持续），但我认为，就地方和叙述所允许的，趁着机会，目前简要地提及几点也是好的，以便通过同时描绘外在之人的动态，并像打下祈祷的基础，然后当我们论及内在之人时，能以较少的劳力建立起完整的祈祷大厦；尤其是要提供这一点：如果生命的终结追上我们，使我们无法完成我们（\OriginalTerm{若天主愿意}{D.V.}）想要恰当撰写的叙述，至少在这部著作中为你留下如此必要之事的开始，对你而言，由于你渴望的热忱，一切都显得延迟；这样，如果我们被赐予再多几年寿命，至少可以为你勾勒出他们祈祷的一些轮廓，使特别是在修院中生活的人对此有所了解；同时也提供那些可能只读到这本书而无法获得另一本书的人，能发现他们对祈祷性质有所了解；并且正如他们在外在之人的服装和衣着上受到教导，同样，他们在呈献属神祭献时应该如何行为也不至于无知。因为，虽然我们现在依靠主的帮助安排撰写的这些书主要涉及外在之人和\OriginalTerm{修院}{Cœnobia}的习俗，但那些书则更关注内在之人的培育、心的成全以及\OriginalTerm{独修隐修士}{Anchorites}的生活和教导。\par

% structure: p0029 source=h2 target=subsection reason=relative-heading-rank; base=h2

\subsection{第十章}

\Emph{论埃及人献集祷经时的静默与简洁。}\par

当他们聚集举行上述礼仪（他们称之为\Emph{集会}）时，大家都非常静默，以至于虽然聚集了如此众多的弟兄，你也不会觉得有人在场，除了那站在中间起立咏唱圣咏的人；尤其是在献上祈祷时更是如此，因为那时没有吐痰、清喉咙、咳嗽声、没有睡意打哈欠、张口打呵欠，也没有发出可能分散旁边站着的人注意的呻吟或叹息。除了司铎结束祈祷的声音之外，听不到任何声音，或许有一种声音是因心不在焉而逃逸，无意中使心灵惊异，这心灵被一种无法控制、无法抑制的心神热忱所燃烧；而那时，炽热的心灵无法保持的，通过一种不可言喻的呻吟，努力从胸膛最内室冲出。但如果任何人被心灵的冷漠感染，大声祈祷或发出我们提到的任何声音，或被一阵哈欠所控制，他们宣称此人犯有双重过错。\par

他有可责之处，首先，就其自身的祈祷而言，因为他以漫不经心的方式将它献给天主；其次，因为他以无礼的喧哗打扰了另一个人的思绪，否则那人或许本能以更大的注意力祈祷。因此，他们的规则是：祈祷应以快速的方式结束，免得我们在拖延期间，多余的唾液或痰涎干扰祈祷的结束。所以，当祈祷仍在热忱中时，要尽快从仇敌的口中夺走它，因为仇敌虽然一直敌视我们，但当他看到我们急于向天主献上祈祷以对抗他的攻击时，他比任何时候都更为敌视；他通过激发游荡的思绪和各种体液，竭尽全力分散我们专注于祈祷的心思，并以此试图使它变冷，尽管开始时是热忱的。因此，他们认为最好的方式是祈祷要简短且频繁地献上：一方面，通过如此频繁地向主祈祷，我们能持续地与他结合；另一方面，当魔鬼在暗中等候我们时，我们可以通过祈祷的简洁简短来避开他意图伤害我们的箭矢，尤其是在我们祈祷的时候。\par

% structure: p0033 source=h2 target=subsection reason=relative-heading-rank; base=h2

\subsection{第十一章。}

\Emph{论埃及人诵咏圣咏的方式。}\par

因此，他们甚至不试图连续不断地完整诵完礼仪中所唱的圣咏。而是按照节数将它们分成两段或三段，分别重复，中间穿插祈祷。因为他们不关心节的数量，而关心心思的理解；全力以赴追求这一点：\BibleQuote{我要以心神歌唱，也要以理智歌唱。}\BibleRef{格前 14:15}因此，他们认为，以理解和思考歌唱十节，比困惑地诵完整篇圣咏更好。而这种情况有时是由于诵咏者的仓促造成的，当他想着剩余圣咏的性质和数量时，他不努力使意思对听众清晰，而是匆忙结束礼仪。最后，如果有较年轻的隐修士，或因精神的热忱，或因尚未受教，超出了应唱圣咏的恰当限度，则诵咏圣咏者会被长者从他坐着的座位上拍手制止，并使他们全体起立祈祷。这样，他们尽可能小心，不让任何倦怠在长时间坐着诵咏圣咏时混入，因为这不仅会使诵咏者本人失去理解的果实，而且会使那些因他持续太久而感到疲惫的人遭受损失。他们还极其注意这一点：即任何圣咏都不应以\Quote{阿肋路亚}的应答方式诵念，除非那些在标题中标注了\Quote{阿肋路亚}的圣咏。但他们将上述十二篇圣咏这样分配：如果有两位弟兄，每人唱六篇；如果有三位，则每人四篇；如果有四位，则每人三篇。他们在集会中从不唱少于这个数目，因此，无论聚集了多大的会众，在礼仪中唱圣咏的弟兄不得超过四位。\par

% structure: p0036 source=h2 target=subsection reason=relative-heading-rank; base=h2

\subsection{第十二章。}

\Emph{论在礼仪中一人诵咏圣咏时其余人坐下，以及其后他们在斗室中延长夜课直至黎明的热忱。}\par

这十二篇圣咏的常规系统，我们已提及，他们通过身体的休息使其更轻松；当他们按惯例举行这些礼仪时，除了一人站在中间诵咏圣咏外，其余人都坐在极低的座位上，以最大的心思注意力跟随诵咏者的声音。因为他们因日夜禁食和工作而疲惫不堪，除非有某种这样的帮助，否则他们不可能站着完成这个数目。因为他们不容许任何时间无工作地流逝，因为他们不仅竭尽全力做那些白天可以完成的手工活，还焦虑地思考那些连黑夜也无法阻止的工作，他们认为，他们越热忱地致力于工作和劳动，便越能通过心灵的纯洁更深入地洞察灵性默观的课题。因此，他们认为，由天主安排的常规祈祷的适度允许是为了给那些信德非常热忱的人留有空间，使他们永不疲倦的德行之流得以倾注，同时也不至于因数量过多而使他们疲惫脆弱的身体产生厌倦。因此，当常规祈祷的礼仪恰当完成后，每个人回到自己的斗室（他独自居住，或只允许与一位因工作或学习门徒训练而结合在一起的人同住，或许还有性格相似者成为同伴），并再次以更大的热忱献上同样的祈祷服务，作为他们特殊的私人祭献；此后，没有人再让自己休息和睡眠，直到白昼的亮光降临，白天的劳作接替夜晚的劳作和默想。\par

% structure: p0039 source=h2 target=subsection reason=relative-heading-rank; base=h2

\subsection{第十三章。}

\Emph{为何夜课之后不允许他们去睡觉的理由。}\par

他们坚持这些劳苦，有两个原因，此外还基于这一考虑：他们认为，当他们勤勉地操劳时，是在向天主献上自己双手成果的祭献。而且，如果我们追求成全，我们也应以同样的勤勉遵守这一点。首先，免得我们嫉妒的仇敌——他嫉妒我们的纯洁，常对我们设计谋，并不断敌视我们——借着梦中的某种幻象，玷污那已通过夜间的圣咏和祈祷获得的纯洁；因为在我们为疏忽和无知作了赔补，并在告解中以深长的叹息恳求了赦免之后，他焦急地试图，若发现有些时间给予休息，就来玷污我们；尤其当他看到我们借着祈祷的纯洁，最热切地趋向天主时，就竭力推翻和削弱我们对天主的信赖；因此，有时当他整夜未能伤害某些人时，便在那短短的一小时内尽其所能羞辱他们。其次，因为即使没有出现如此可怕的魔鬼幻象，单是睡眠的间隙也会使那本该很快醒来的隐修士产生懈怠；并且，在心灵中引起迟钝的麻木，使他一整天的精力迟钝，削弱那敏锐的觉察力，耗尽那心灵的力量——这力量本可使我们全天更警醒地防范仇敌的一切陷阱，更为刚强。因此，在日课之外，还加上了这些私人守夜，他们以更大的谨慎躬行此事，既为不失去通过圣咏和祈祷获得的纯洁，也为借着夜间的默想，预先确保更热切的谨慎，以在白天勤勉地守护我们。\par

% structure: p0042 source=h2 target=subsection reason=relative-heading-rank; base=h2

\subsection{第十四章。}

\\Emph{论他们在静修室中如何同样致力于手工劳作和祈祷。}\par

因此，他们在祈祷之外又加上劳作，免得懒惰者被睡意偷袭。因为他们几乎没有任何闲暇时间，所以对属灵默想也没有限制。因为同时操练身体和灵魂的德行，他们用那有益于内在之人的来平衡外在所当行的；以\\Emph{劳作}的重量，如同坚固不动的锚，稳住心灵易滑动的活动和思绪的变迁；藉此，心灵的变化无常和游荡被固定在静修室的界限内，如在完全安全的港湾中闭锁；于是，只专注于属灵默想和思想的警醒，不仅禁止警觉的心灵轻易同意任何邪恶的暗示，而且使其免于任何无用和空闲的念头。因此，很难说哪一方依赖另一方——我是说，他们是为着属灵默想而进行不间断的体力劳作，还是为着持续不断的劳作而获得如此卓越的属灵进步和知识的明光。\par

% structure: p0045 source=h2 target=subsection reason=relative-heading-rank; base=h2

\subsection{第十五章。}

\\Emph{论明智的规则：每人祈祷结束后必须回到自己的静修室；以及若有违反者将受谴责。}\par

因此，当圣咏唱完，日常的集会如我们上面所说解散后，他们中无人敢稍作停留或与另一人闲谈；也不敢在整天内离开静修室，或放弃通常在静修室中进行的工作，除非他们偶然被召出去执行某项必要任务——他们出外完成这些任务，以便在他们中间绝无任何闲谈。但每个人都以这样的方式完成分派给他的工作：通过默诵某篇圣咏或一段圣经经文，不让口和心有任何机会或时间从事危险的图谋、邪恶的意图，甚至闲谈，因为口和心不停地专注于属灵默想。因为他们特别严格遵守这一规则：无人，尤其是较年轻者，可能被发现在片刻间与另一人同行，或互相握手。但是，如果他们发现有人违犯此规则的纪律而实施了这些禁止之事中的任何一件，他们便判定那人犯有严重过错，是顽固不化、不服从规则；并且他们不会免去那人参与谋划和邪恶设计的嫌疑。而且，除非那人在全体弟兄聚集时通过公开的补赎赎罪，否则他的过错不被允许参与弟兄们的祈祷。\par

% structure: p0048 source=h2 target=subsection reason=relative-heading-rank; base=h2

\subsection{第十六章。}

\\Emph{论无人被允许与暂停祈祷的人一同祈祷。}\par

再者，若有人因所犯的过错被暂停祈祷，在他在众人面前俯伏做补赎，且院长当着所有兄弟的面公开赐予他宽恕和赦免之前，无人有权与他一同祈祷。因为藉此类计划，他们如此与那人断绝祈祷的共融，原因在于：他们相信被暂停祈祷的人，正如宗徒所说，是\\BibleQuote{交付与撒殚：}；若有人出于不合时宜的情感，在未被长者接纳之前就胆敢与他一同祈祷，便使自己分担他的丧亡，并自愿把自己交付给撒殚——那人原是为了纠正罪过而被交托给撒殚的。在此他犯下更严重的过错，因为既在谈话中又在祈祷中与他联合共融，便给他提供更大的傲慢理由，只会鼓励和加剧犯错者的顽固。因为给予他只有害处的安慰，会使他的心更硬，不让他为被暂停祈祷的过错而自卑；并由此使他对长者的谴责视为无关紧要，对赔补和赦免怀有欺诈的念头。\par

% structure: p0051 source=h2 target=subsection reason=relative-heading-rank; base=h2

\subsection{第十七章。}

\\Emph{论那唤醒他们祈祷的人应如何按时召唤他们。}\par

然而，被托付召集宗教集会并负责礼仪服侍的人，不应随心所欲、或按自己夜间醒来的时间、或随自己睡眠或失眠的偶然倾向，擅自不规律地唤醒弟兄们进行每日守夜祈祷。尽管每日的习惯可能迫使他在惯常的时刻醒来，但他仍应频繁而谨慎地根据星辰的轨迹确定正确的祈祷时刻，召唤他们前来祈祷，以免在两方面有所疏忽：要么因贪睡而错过夜间适当的时刻，要么因急于就寝、渴望睡眠而提前召唤，从而被人认为只顾自己的休息，而非关注属灵的职分和他人的安宁。\par

% structure: p0054 source=h2 target=subsection reason=relative-heading-rank; base=h2

\subsection{第十八章。}

\Emph{他们如何从星期六晚上至主日晚上不跪拜。}\par

这一点我们也应当知道：从星期六晚上（即主日前夕）到次日晚上，以及在埃及人中，从复活节到圣神降临期，他们从不跪拜；在这些时期，他们也不遵守禁食的规则。其原因，若主允许，将在\WorkTitle{长老会谈}的适当地方予以解释。目前，我们只打算简要地概述这些原因，以免本书超出应有的篇幅，使读者感到厌倦或负担。\par

\SourceRecord{Translated by C.S. Gibson.；Nicene and Post-Nicene Fathers, Second Series；Vol. 11.；Edited by Philip Schaff and Henry Wace.；Buffalo, NY: Christian Literature Publishing Co.,；1894.；<http://www.newadvent.org/fathers/350702.htm>.}

% structure: p0001 source=h1 target=section reason=page-title

\section{\WorkTitle{圣规 第三卷}}

% structure: p0002 source=h2 target=subsection reason=relative-heading-rank; base=h2

\subsection{第一章}

\Emph{论叙利亚地区所遵守的第三、第六、第九时辰的祈祷}。\par

在埃及全境所遵守的夜间祈祷与圣咏制度，据我拙见，已赖天主的助佑，尽我绵力阐述完毕。如今必须论及\OriginalTerm{第三时辰}{Tierce}、\OriginalTerm{第六时辰}{Sext}及\OriginalTerm{第九时辰}{None}的祈祷，如我们在序言中所说，此祈祷乃按巴勒斯坦及美索不达米亚隐修院的规矩，并须借这些地方的习俗，对埃及纪律之完美与不可仿效的严格加以调节。\par

% structure: p0005 source=h2 target=subsection reason=relative-heading-rank; base=h2

\subsection{第二章}

\Emph{埃及人如何终日持续祈祷和圣咏，并加上劳作，不拘时辰。}\par

因为在埃及人中，那些我们按不同时辰、相隔一段时间、由召集人提醒才奉行于主的祈祷，却借着加上劳作，且出于他们自愿，在整个白天持续不断地举行。原来在他们的小室里，手工劳作从不间断，而默想圣咏及圣经其他部分也从未完全停止。并且他们在每一时刻都将祈祷与祈求混入劳作中，整天都从事那些我们在固定时辰所举行的祈祷。因此，除了\OriginalTerm{晚祷}{Vespers}与\OriginalTerm{夜祷}{Nocturns}之外，他们白天没有公共祈祷，除非在星期六和星期日，那时他们在第三时辰为了领受圣体而聚集。因为持续不断的奉献，胜过间隔时间的奉献；作为自由奉献，也比出于规则强制的义务更蒙悦纳；为此达味有些得意地欢呼说：\BibleQuote{我必自由地向你献祭；}又说：\BibleQuote{上主，愿我口中的自由奉献蒙你悦纳。}\par

% structure: p0008 source=h2 target=subsection reason=relative-heading-rank; base=h2

\subsection{第三章}

\Emph{整个东方如何以每时辰仅三篇圣咏及祈祷来完成第三、第六、第九时辰的祈祷；以及为何这些属灵祈祷更特别地归于这些时辰。}\par

所以，在巴勒斯坦、美索不达米亚和整个东方的隐修院里，上述时辰的日课每日以三篇圣咏结束，以便在指定的时间不断向天主献上祈祷，同时，由于神业以适当节制完成，必要的工作职责也不致受到任何妨碍：因为我们知道，达尼尔先知也在这三个时辰，每日在开着窗户的房间里向天主倾心祈祷。\BibleRef{达 6:10} 这些时辰被特别指定用于宗教敬礼，并非没有充分的理由，因为在这些时辰，那完成了诸项应许并总括我们救恩的事得以实现。因为我们可以证明，在第三时辰，那自古由先知所预许的圣神，首先降临在聚集祈祷的宗徒们身上。因为当不信的犹太人民因他们因圣神倾注而说舌音感到惊奇时，嘲笑说他们充满了新酒，于是伯多禄站在他们中间说：\BibleQuote{以色列人及一切住在耶路撒冷的人！你们要知道，要留意我的话。这些人并不像你们所想象的那样醉，因为现在只是白天第三时辰；但这正是藉岳厄尔先知所说的：在末日，天主说，我要将我的神倾注在一切有血肉的人身上，你们的儿子和女儿要说预言，你们的青年要见异象，你们的老年人要做异梦。在那些日子里，我甚至要将我的神倾注在仆人和婢女身上，他们要说预言。} \BibleRef{宗 2:14} 这一切都在第三时辰应验了，当那先前由先知们宣告的圣神，在那时辰来临并停留在宗徒们身上。但在第六时辰，那无玷的祭献、我们的主和救主，被献给圣父，且为拯救全世界登上十字架，为人类的罪作了补赎，剥夺了率领者和掌权者，公开地将他们掳走；我们所有本应死亡、被无法偿还的债据所束缚的人，祂解救了，将债据从中间取走，钉在祂的十字架上作为战利品。同一时辰，伯多禄神魂超拔时，天主启示了他：外邦人的蒙召，由从天上降下的福音器皿所表明，以及里面所有活物的洁净，当时有声音向他说：\BibleQuote{起来，伯多禄，宰了吃！} 这由四角从天上降下的器皿，显然不是别的，正是福音。因为虽然它按四福音作者的叙事被分开，似乎有\Quote{四角}（或开端），但福音的躯体却只有一个；它包含同一个基督的诞生和神性，以及奇迹和受难。它也巧妙地说，不是\Quote{细麻布}，而是\Quote{\Emph{仿佛}是细麻布}。因为细麻布象征死亡。既然我主的死亡与受难并非出于人性规律，而是出于祂自己的自由意志，所以说\Quote{仿佛是细麻布}。因为祂虽按肉体死了，按灵却未死，因为\BibleQuote{祂的灵魂没有被留在阴府，祂的肉身也没有见到腐朽。} 祂又说：\BibleQuote{没有人夺去我的生命，是我自己舍弃的。我有权舍弃，也有权再取回。} \BibleRef{若 10:18} 因此，在这从天上降下的福音器皿中，即由圣神所记载的，所有先前在法律之外、被视为不洁的万民，现在因信仰的信德而聚集，为获救恩而远离偶像崇拜，成为有益健康的食物，并被带到伯多禄面前，藉主的声音而得洁净。但在第九时辰，祂深入阴府，以祂辉煌的光辉熄灭了地狱那难以言表的黑暗，冲破其铜门，打碎铁闩，将祂带往天上的那被囚的圣者队伍，他们原被关押在严酷地狱的黑暗中，祂取走了火剑，藉祂虔诚的信德宣告，使原住民重返乐园。同一时辰，百夫长科尔乃略，以惯常的虔诚继续祈祷，通过天使与他的交谈，得知他的祈祷和施舍在主前蒙记念，并在第九时辰，外邦人蒙召的奥秘清楚地显示给他，这奥秘曾在第六时辰于神魂超拔中启示给伯多禄。在《宗徒大事录》的另一处，关于同一时间，我们读到：\BibleQuote{伯多禄和若望在祈祷的时辰，即第九时辰，上圣殿去。} \BibleRef{宗 3:1} 由这些记载清楚证明，这些时辰被圣洁的宗徒们奉献于宗教敬礼，并非没有充分理由，我们同样应当遵守；除非我们被某种规则强迫，至少在规定时间履行这些虔诚职务，否则我们或因懒惰，或因遗忘，或因沉浸于事务，终日不祈祷。但关于晚祭，又当如何说呢？因为在旧约中，这些已由梅瑟法律规定要不断地奉献？因为早晨的全燔祭和晚上的祭献，在圣殿中日复一日地奉献，尽管是用象征的祭品，我们可以从达味所唱的诗歌中得知：\BibleQuote{愿我的祈祷在祢面前如馨香上腾，愿我双手的举起如同晚祭。} 在此，我们可按更高意义理解为那真正的晚祭，即我们的主救主在晚餐时于傍晚赐给宗徒们的，祂建立了教会的神圣奥迹，以及那真正的晚祭，即祂自己在次日、在时代的终结，举起双手为拯救全世界献给圣父的；祂在十字架上伸开双手，确可恰当地称为\Quote{高举}。因为当我们都躺在阴府时，祂将我们提升到天上，正如祂自己应许的话说：\BibleQuote{当我从地上被举起来时，我要吸引众人归向我。} \BibleRef{若 12:32} 但关于晨祷，那每日惯常歌唱的也教导我们：\BibleQuote{天主，我的天主，我向祢在破晓时守望；} 以及\BibleQuote{我将在清晨默想祢；} 以及\BibleQuote{我在天破晓前呼求；} 又说\BibleQuote{我的双眼在祢面前胜过清晨，为使我默想祢的话语。} 在这些时辰，福音中的家主也雇工人进他的葡萄园。因为如此描述：他清晨雇了工人，这时间表示晨祷时辰；然后第三时辰；第六时辰；之后第九时辰；最后第十一时辰\BibleRef{玛 20:1}，这表示点灯时分。\par

% structure: p0011 source=h2 target=subsection reason=relative-heading-rank; base=h2

\subsection{第四章}

\Emph{晨祷并非由古老传统所制定，而是在我们这个时代为特定理由而开始。}\par

但你们必须知道，现今在西方国家非常普遍奉行的\OriginalTerm{晨祷}{Mattin}，是在我们这个时代、也是在我们自己的隐修院中被定为会规时辰的。在这隐修院里，我们的主耶稣基督由童贞女诞生，并屈尊作为人经历婴孩的成长，又藉祂的恩宠扶持了我们尚在宗教薰陶中的婴孩时期，以奶哺育。因为直到那时，我们发现，在高卢的隐修院中，当夜祷的圣咏和祈祷结束后不久通常举行的晨祷时辰，连同每日的夜祷一同完成后，剩余的时间由我们的长老们分配给身体的休息。但有些人相当漫不经心地滥用这种恩惠，过分延长睡眠时间，因为他们不受任何礼仪要求的约束，直到第三时辰才离开他们的房间或起床；而且，除了丧失劳动外，他们在白天本应专注于某些职责时（尤其是在那些夜晚从黄昏至清晨守夜造成异常沉重的疲倦的日子里），却因睡眠过度而昏昏欲睡。于是，一些热心且对此漫不经心深感不安的弟兄向长老们投诉。长老们经过长时间讨论和慎重考虑后决定：在日出之前，当他们可以无害地准备阅读或从事体力劳动时，应给疲乏的身体休息时间；此后，应召集所有人遵守此礼仪，从床上起来，背诵三篇圣咏和祈祷（按照古时定为第三和第六时辰礼仪的顺序，以表示对天主圣三的宣认），从而同时以统一的安排结束睡眠并开始工作。此形式虽看似出于偶然，且在近期因上述理由而制定，但它显然按字面实现了有福的达味所指的数字（虽然也可从灵性上理解）：\BibleQuote{我每日七次赞美祢，因为祢的判决公正。}因为增加了此礼仪，我们确实每日七次举行这些灵性集会，并显示在其中七次歌颂天主。最后，虽此形式从东方开始，极有益地传播至这些地区，但在东方一些古老隐修院中（它们不愿容忍对父辈古老规矩的丝毫违反），似乎从未被引入。\par

% structure: p0014 source=h2 target=subsection reason=relative-heading-rank; base=h2

\subsection{第五章}

\Emph{他们不应在晨祷后再次回到床上。}\par

但此地区有些人不知此礼仪为何制定和引入，在完成晨祷后又回到床上，反而陷入那种长老们为防范而设立此礼仪的习惯中。因为他们急于在那一时刻结束晨祷，以便给那些倾向于漠不关心和不够谨慎的人提供再次上床的机会。这是绝不应做的（如我们在前册描述埃及人的礼仪时更详细地指出的），以免我们自然情欲的力量被激起，玷污我们藉黎明前的谦卑告解和祈祷所获得的纯洁，或以免仇敌的幻象污染我们，甚至纯洁自然的休息干扰我们灵性的热忱，使我们整天懒散怠惰，因睡眠引起的迟钝而冰冷。为避免此点，埃及人，尤其是他们习惯于在鸡鸣前甚至固定时间起床，在举行会规礼仪后，继续守夜直至天亮，这样，当晨光降临他们时，会发现他们保持着灵性的热忱，并使他们整天更加谨慎和热忱，因他们已通过夜间守夜和灵性默想为对抗魔鬼的日常争斗做好准备并得到加强。\par

% structure: p0017 source=h2 target=subsection reason=relative-heading-rank; base=h2

\subsection{第六章}

\Emph{在制定晨祷时，长老们并未改变古代圣咏系统。}\par

但我们也应知道，即决定添加此晨祷礼仪的长老们并未改变古代圣咏的安排；而是此礼仪在他们的夜间集会中始终按以前相同的顺序举行。因为在此地区用于晨祷礼仪的赞美诗，即在鸡鸣后黎明前结束夜祷时所惯用的，他们仍以同样方式咏唱；即咏148，开始为\BibleQuote{请从天上赞美上主}，以及其余续篇；但第五十篇圣咏、第六十二篇和第八十九篇，我们知道被分配给了这个新礼仪。最后，在全意大利直到今天，当晨祷赞美诗结束后，所有教堂都咏唱第五十篇圣咏，我确信这只能源于此来源。\par

% structure: p0020 source=h2 target=subsection reason=relative-heading-rank; base=h2

\subsection{第七章}

\Emph{凡在第一篇圣咏结束前未到每日祈祷者，不得进入祈祷所；但在夜祷中，迟到至第二篇圣咏结束时可被宽恕。}\par

但在第三时辰、第六时辰或第九时辰，若有人在圣咏开始和结束前尚未前来祈祷，便不敢再进入祈祷所，也不敢加入唱圣咏者的行列；他站在外面，等待会众散去，当众人出来时，他俯伏在地做补赎，并为自己的疏忽和迟到获得赦免，他知道除非藉着真谦卑立即为自己的过失做补赎，否则无法赎清懈怠之过，也绝不能被允许参加三小时后举行的日课。但在夜间集会中，迟到至\Emph{第二}篇圣咏是被允许的，前提是在圣咏结束、弟兄们俯首祈祷之前，他急忙加入会众的行列；但若他迟延超过允许的时限一点点，就必定会受到我们之前提到的同样责备和补赎。\par

% structure: p0023 source=h2 target=subsection reason=relative-heading-rank; base=h2

\subsection{第八章}

\Emph{论安息日前夕举行的夜课；其长度及遵守方式}\par

然而在冬季，当夜晚较长时，隐修院的长老们将每周安息日开始前夕举行的夜课安排至第四次鸡鸣，理由如下：经过整夜守望后，他们可在剩余约两小时的时间里让身体休息，以免整日因困倦而萎靡不振，并以这短暂的休息代替整夜安眠。我们也应当极其谨慎地遵守这一点，即满足于夜课结束后至黎明——即至晨祷圣咏——所允许的睡眠，之后将整个白天用于工作和必要职责，以免因夜课带来的疲惫和虚弱，被迫在白天补回夜间削减的睡眠，从而被认为不是减少身体的休息，而是改变了休息和夜间睡眠的时间。因为我们软弱的肉身不可能完全被剥夺整夜休息，却仍能在次日保持精力充沛而不致精神困倦、心灵沉重；若夜课结束后未能享受片刻小睡，反而会受妨碍而非帮助。因此，若如我们所建议，在黎明前至少争取一小时的睡眠，我们就能保全整夜祈祷所花费的全部夜课时间，既给予本性当得的份额，也无须在白天补回夜间削减的睡眠。因为人若不以理性方式只减少一部分，而试图完全拒绝，并（坦率地说）急于削减的不是多余的部分而是必要的部分，就必将把一切让给这肉身。因此，若夜课不合理地、无度地延长至黎明，就必须以更大的代价来补偿。于是他们将夜课分为三段经课，藉此多样性分散劳苦，并通过某种愉快的放松缓解身体的疲惫。因为他们站着轮流唱完三篇圣咏后，便坐在地上或极低的座位上，由一人轮流咏唱三篇圣咏，其余人应答，每篇圣咏指派一位弟兄，依次轮流；他们还在静坐中添加三篇读经。如此，通过减少身体的劳累，他们能以更集中的心神守夜课。\par

% structure: p0026 source=h2 target=subsection reason=relative-heading-rank; base=h2

\subsection{第九章}

\Emph{为何在安息日黎明举行夜课，以及为何在整个东方在安息日免于斋戒}\par

在整个东方，自从宗徒宣讲、基督徒信仰与宗教奠基之时起，便已规定这些夜课应在安息日黎明举行，原因如下：当我们的主和救主在一周第六日被钉十字架后，门徒们因祂受苦的新近而悲痛，整夜守候，不让眼睛休息或睡觉。因此，自那时起，这一夜便设立了夜课礼仪，至今仍在整个东方以同样方式遵守。所以，在夜课劳苦之后，宗徒们同样规定在安息日免于斋戒，这在东方所有教会中并非没有理由地施行，符合训道篇的这句话（尽管它有另一种神秘意义，但用在此处也并非不当），它吩咐我们要将同样的服务分给两日——即第七日和第八日——相等，正如经上所说：\BibleQuote{当将一份给这七个，也给这八个。}\BibleRef{训 11:2}因为不应将这免斋理解为那些显然摆脱一切犹太迷信的人参与犹太节日，而是如我们所说，有助于疲乏身体的休息；身体连续五天每周斋戒，全年如此，若不以至少两天的间隔恢复，就很容易疲惫衰竭。\par

% structure: p0029 source=h2 target=subsection reason=relative-heading-rank; base=h2

\subsection{第十章}

\Emph{如何在城中导致他们在安息日斋戒}\par

但是，在西方某些地区，尤其是这座城市里，有些人不明白此宽免的原因，认为在安息日绝对不应允许免除斋戒，因为他们说宗徒\Person{伯多禄}正是在这一日，在遇见\Person{西满}之前禁食了。由此清楚可见，他这样做并非遵循教会法规则，而是因着即将面临的斗争的必需。同样，在那里，为同一目的，\Person{伯多禄}似乎给门徒们规定的不是公斋，而是特斋；如果他知道斋戒是通过教会法规则规定的，他一定不会这样做；正如如果斗争的机会落在主日，他肯定也会准备在主日规定斋戒：但是由此并不能制定出普遍性的教会法斋戒规则，因为导致此事的并非普遍的惯例，而是迫于必需，这迫使他在一个单独的场合遵守斋戒。\par

% structure: p0032 source=h2 target=subsection reason=relative-heading-rank; base=h2

\subsection{第十一章。}

\Emph{论主日礼仪与其他日子惯常礼仪不同之处。}\par

但是，我们也应当知道，在主日，午餐前只庆祝一台日间祈祷，出于对礼仪本身和主的圣体的敬意，他们使用更隆重且更长的圣咏、祈祷和读经的礼仪，并因此认为\OriginalTerm{第三时辰}{Tierce}和\OriginalTerm{第六时辰}{Sext}都被包含其中。由此导致，由于增加了读经，他们的祈祷量并未减少，然而却有所区别，为尊敬主的复活，似乎对弟兄们施予了在其他时间的宽免；这似乎减轻了整周的遵守，并且由于插入的区别，使主日更隆重地被期待为一个庆日，由于期待来临的一周的斋戒也较少感到。因为任何疲惫总是以更大的平静承受，若插入一些变化或交替工作，无厌倦地承担的劳作也会变得容易。\par

% structure: p0035 source=h2 target=subsection reason=relative-heading-rank; base=h2

\subsection{第十二章。}

\Emph{论当为弟兄们提供晚餐时，召集用餐时不像午餐通常那样咏唱圣咏的日子。}\par

最后，也在那些日子——即星期六和主日以及庆日——在这些日子，通常为弟兄们提供午餐和晚餐，晚上不咏唱圣咏，无论是他们进入晚餐时还是结束后，不像在他们的常规午餐和斋戒日的惯例小斋中那样，惯常的圣咏通常在其前后咏唱。他们只是做一个简短的祈祷，然后进入晚餐，又在结束时，单独以祈祷结束；因为这一餐在隐修士中是特殊的：并非所有的人都必须参加，而只是为前来探望弟兄们的旅客，以及那些因身体虚弱或自身意愿而受邀的人。\par

\SourceRecord{Translated by C.S. Gibson.；Nicene and Post-Nicene Fathers, Second Series；Vol. 11.；Edited by Philip Schaff and Henry Wace.；Buffalo, NY: Christian Literature Publishing Co.,；1894.；<http://www.newadvent.org/fathers/350703.htm>.}

% structure: p0001 source=h1 target=section reason=page-title

\section{\WorkTitle{圣规（第四卷）}}

% structure: p0002 source=h2 target=subsection reason=relative-heading-rank; base=h2

\subsection{第一章}

\Emph{关于那些弃绝此世之人的培训，以及在塔本纳和埃及的隐修士中，被接纳入隐修院者所受教育的方式。}\par

从每日应在隐修院中遵守的圣咏与祈祷的规范体系，我们按照叙述的适当顺序，转而讨论弃绝此世之人的培训；首先，尽我们所能，用简短记述概括那些渴望转向主的人被接纳入隐修院的条件；并加入一些来自埃及人的规则，以及来自塔本纳隐修士的规则，该隐修院位于底比斯，其人数之众多、纪律之严格均超过其他所有隐修院，因为其中有五千多位弟兄同属一位院长管辖；而所有隐修士对一位长老始终如一的服从，是我们中间无人能暂时对他人做到或要求的。\par

% structure: p0005 source=h2 target=subsection reason=relative-heading-rank; base=h2

\subsection{第二章}

\Emph{在他们中，人们甚至在极老的年纪仍留在隐修院的方式。}\par

我认为，在一切之前，我们应当论及他们那持久的坚忍、谦卑和服从——它如何持续如此长久，又由怎样的体系形成，使他们得以留在隐修院直至年老背弯；因为它如此巨大，以致我们无法记起有谁加入我们的隐修院并连续遵守它一整年；这样，当我们看见他们弃绝此世的开始，便会明白，从这样的起点出发，他们如何达到如此完美的境界。\par

% structure: p0008 source=h2 target=subsection reason=relative-heading-rank; base=h2

\subsection{第三章}

\Emph{考察将被接纳入隐修院之人的考验。}\par

因此，一个寻求被接纳进入隐修院纪律的人，在他在门外躺卧十天或更久，证明了他的坚忍和渴望，以及谦卑和忍耐之前，永不被接纳。当他俯伏在所有经过的弟兄脚前，并被所有人刻意拒绝和轻视，仿佛他进入隐修院并非为了宗教而是出于被迫；当他受到许多侮辱和冒犯，并以实际证据证明了他的坚定，通过忍受羞辱显明他在试探中将会如何；当借此确认了他灵魂的热忱而被接纳后，他们便极其仔细地查问，他是否还沾染着哪怕一枚来自先前财产的硬币。因为他们知道，如果他知晓即使极小的一笔钱被隐藏，他便无法长久留在隐修院的纪律之下，也永不能学会谦卑和服从的德行，也不会满足于隐修院的贫穷和艰苦生活；反而，一旦因某种缘由兴起骚动，他便会立刻像从投石器中射出的石头那样从隐修院飞走，因信赖那笔钱财而受驱动。\par

% structure: p0011 source=h2 target=subsection reason=relative-heading-rank; base=h2

\subsection{第四章}

\Emph{为什么被接纳入隐修院的人不允许带任何东西进去的理由。}\par

基于这些理由，他们不同意收下他的钱财，甚至为了隐修院的利益也不行：首先，以免他因这奉献而心生傲慢，不屑于与较贫穷的弟兄平等；其次，以免他因这骄傲而不能屈就基督的谦卑，以致无法承受隐修院的纪律而离开，之后当他冷静下来，会怀着恶意想要取回他在弃绝初期、属灵热忱炽热时奉献的东西——这会给隐修院带来损失。他们经常从许多实例中得知，这条规则应始终遵守。因为在一些不够谨慎的隐修院中，有些毫无保留被接纳的人，后来竟以最亵渎的方式要求归还他们已经奉献并用于天主工作的财物。\par

% structure: p0014 source=h2 target=subsection reason=relative-heading-rank; base=h2

\subsection{第五章}

\Emph{为什么那些舍弃世界的人，在被接纳入隐修院时，必须脱下自己的衣服，由院长给他们穿上别样的衣服。}\par

\Emph{因此}，每位入修者被剥夺先前一切所有物，甚至不再允许保留身上的衣服；而是在弟兄们的议会上，他被带到中间，脱下自己的衣服，由院长亲手穿上隐修院的服装，借此他不仅知道被剥夺了一切旧物，而且知道放下了所有世俗的骄傲，降到基督的贫穷与匮乏中，从此不再依靠世俗手段获得的财富，也不依靠从不信主时的状态中保留的任何东西来维持生活，而是要从隐修院的圣洁基金中领取他服务的给养；并且，既然他知道自己从此将由该处得到衣着和食物，且一无所有，他也能学会不为明天忧虑，按照福音的教导，并且不以与穷人——即弟兄们的团体——平等为耻，基督不以被列在其中并称他们为弟兄为耻，反而要因被置于与自己的仆役同伙而自豪。\par

% structure: p0017 source=h2 target=subsection reason=relative-heading-rank; base=h2

\subsection{第六章}

\Emph{为什么院长保管弃绝者加入隐修院时所穿的衣服。}\par

但那些他脱下的衣服，交由总管保管，直到藉由各种不同的诱惑和考验，他们能认出他进步、生活与忍耐的美德。若他们见他随着时间的推移能持之以恒，并保持起初的热忱，便将衣服施舍给穷人。但若发现他有任何抱怨的过失，或甚至最小的不顺服，便从他身上剥下隐修院的衣服，重新给他穿上已被没收的旧衣，将他赶走。因为他不应穿着所领受的衣服离开，也不允许任何人一旦在修院规矩上冷淡下来后，再穿着这些衣服。因此，也没有公开出去的机会，除非他像逃跑的奴隶一样，趁黑夜最浓密时逃脱，或被判定不配此一会规与身份，脱去隐修院的衣服，当着众弟兄的面，蒙羞受辱地被驱逐出去。\par

% structure: p0020 source=h2 target=subsection reason=relative-heading-rank; base=h2

\subsection{第七章}

为何被接受入隐修院的人，不得立即与弟兄团体共处，而要先托付给客房\par

任何人经过我们所说的那种坚持而被接纳后，脱下自己的衣服，穿上隐修院的衣服时，他不得立即与弟兄团体共处，而是交给一位长老，这位长老独自住在离隐修院入口不远处，负责照料陌生人和客人，并尽心竭力地善待他们。他在那里服务满一年，没有任何抱怨，并向陌生人提供了服务的证据，以此接受了谦逊与忍耐的初步训练，并在长期的实践中得到认可；当他要从这里被接纳进入弟兄团体时，他再交给另一位长老，后者管理十位初学弟子，这些人由院长托付给他，他按照我们在\BibleBook{出谷}{纪}中读到梅瑟所安排的制度，教导和管理他们。\par

% structure: p0023 source=h2 target=subsection reason=relative-heading-rank; base=h2

\subsection{第八章}

初学弟子首先操练哪些实践，以便能精通克服自己的一切意愿\par

他的焦虑和教导的主要部分——通过它，交托给他的初学弟子能够适时攀登至完美的高峰——将是首先教导他克服自己的意愿；他特意精心训练他这一点，并故意安排一些他知道会违背他喜好的命令；因为他们通过许多实例教导说，一位隐修士，尤其是年轻的隐修士，除非首先通过服从学会克制自己的意愿，否则无法控制情欲的冲动。因此他们认定，一个尚未学会克服自己意愿的人，不可能灭除愤怒、愠怒或邪淫的精神；也无法保持内心的真正谦逊，或与弟兄们持久的合一与稳定的和谐；更无法在隐修院中长久居留。\par

% structure: p0026 source=h2 target=subsection reason=relative-heading-rank; base=h2

\subsection{第九章}

为何命令初学弟子不得向长老隐瞒任何思想\par

藉由这些实践，他们便急于给那些正在接受训练的人——如同以字母表和最初音节来导向成全——打下烙印和教导，因为他们能由此清楚看出这些人是否建立在虚假和想象的基础之上，还是建立在真正的谦逊之上。为了易于达到这一点，他们接下来教导这些初学弟子，不要因虚假的羞耻而隐藏心中任何刺人的思想，而是每当这些思想一出现，就立刻向长老坦白；在判断这些思想时，不要相信自己的任何辨别力，而是相信长老的判断所认为和宣判为善或恶的。这样，我们狡猾的仇敌便无法以任何方式诱骗年轻而缺乏经验的隐修士，或利用他的无知胜过他，或藉任何诡计欺骗那受长老辨别力而非自己辨别力保护的人，也不能说服他向长老隐瞒那些像火箭般射入他心中的建议；因为魔鬼虽狡猾，除非他藉着骄傲或羞耻引诱初学弟子隐藏思想，否则无法败坏或毁灭他。他们定下一条普遍而明确的证据：若我们对向长老坦白感到羞耻，那思想便是来自魔鬼。\par

% structure: p0029 source=h2 target=subsection reason=relative-heading-rank; base=h2

\subsection{第十章}

初学弟子甚至在那些属于共同必需之事上的服从是何等彻底\par

其次，他们以如此严格的服从遵守规矩，以至于初学弟子在没有上级的知晓和许可下，不仅不敢离开自己的小室，而且甚至不敢擅自满足共同的自然需要。因此，他们迅速且不加任何讨论地执行上级所命令的一切，仿佛这些命令来自天上的天主；以至于有时，当命令他们做不可能的事时，他们怀着如此的信心和虔敬去承担，以致全力以赴、毫不迟疑地努力完成和执行；出于对长老的尊敬，他们甚至不考虑命令是否可能。关于他们的服从，我目前暂且不多谈，因为我们打算在适当的地方稍后谈，并附上实例——若藉着你们的祈祷，主保守我们安然度过。我们现在转向其他规矩，省略所有那些无法在这个国家的隐修院中施加或遵守的规矩，正如我们在序言中所承诺的；例如，他们从不穿羊毛衣服，只穿棉衣，且是单层的；每位上级在看见他所照管的十位隐修士衣服变脏时，便给他们换发。\par

% structure: p0032 source=h2 target=subsection reason=relative-heading-rank; base=h2

\subsection{第十一章}

他们认为哪种食物更美味\par

我也略过那种艰难而崇高的自制，他们认为最大的奢侈莫过于将那名为\Quote{白屈菜}的植物，用盐腌制并浸水后摆在桌上供弟兄们享用；还有许多类似的事，在本地的气候和我虚弱的体质都无法容许。我只追述那些无论肉体软弱或地域环境都无法阻碍的事，只要心灵的软弱或精神的冷淡不将它们消除。\par

% structure: p0035 source=h2 target=subsection reason=relative-heading-rank; base=h2

\subsection{第十二章}

他们如何一听到有人敲门的声音就立即放下一切工作，急切地立刻应答。\par

因此，他们坐在自己的单间中，同样专心致志于工作和默想，当听到有人敲门叩击各人单间的声音，召唤他们去祈祷或某项工作时，每个人都急切地冲出单间，以至于正在练习写字的人，即便刚刚开始写一个字母，也不敢完成它，而是在叩门声传入耳际的那一刻以最快的速度跑出去，甚至不等写完已开始的那个字母；他们不力求缩减和节省自己的劳力，反而以最热切的诚恳和热情去追求服从的德行，他们认为服从不仅优于手工劳作、阅读、静默和独居，而且超越所有德行，以至于他们认为一切事情都应让位于服从，并且甘愿承受任何不便，只要能表明他们丝毫没有忽略这一德行。\par

% structure: p0038 source=h2 target=subsection reason=relative-heading-rank; base=h2

\subsection{第十三章}

任何人声称任何东西，哪怕微不足道，是他自己的，被视为何等错误。\par

在其他操练中，我想甚至无需提及这一德行，即：不准任何人拥有自己的箱子或篮子作为私人财产，也不得拥有任何可以用自己印章封存的东西，因为我们很清楚，他们被剥夺得如此彻底，除了衬衫、斗篷、鞋子、羊皮和草席外一无所有；因为即使在那些允许某些宽限和放松的其他隐修院里，我们看到这条规则仍被严格恪守，以至于没有人敢在言语上声称任何东西是自己的：如果一位隐修士口中漏出诸如\Quote{我的书}、\Quote{我的写字板}、\Quote{我的笔}、\Quote{我的外套}或\Quote{我的鞋子}之类的词，便是大错；若因疏忽或无知偶尔脱口而出此类话，他必须通过适当的补赎来赎罪。\par

% structure: p0041 source=h2 target=subsection reason=relative-heading-rank; base=h2

\subsection{第十四章}

即使每人通过劳动积累了大量金钱，仍无人敢超出所规定的适当限度。\par

尽管他们每人每天靠工作劳动给隐修院带来如此丰厚的收入，以致他不仅足以满足自己适度的需求，还能丰裕地供给许多人的需要，但他毫不骄傲，也不因自己的辛劳和由此带来的巨大收益而自夸，除了两个在当地售价几乎不到三便士的饼子外，无人认为自己有权获得更多。在他们中间，没有任何东西（我说这话感到羞愧，并且衷心希望在我们自己的隐修院里无人知晓）被任何人——我不敢说行为上，甚至思想中——视为自己的私产。虽然他相信隐修院的整个粮仓都是他的财富，并作为万有的主人将全部关心和精力投入其中，然而，为了保持他已获得并努力完整保持到最后的那个卓越的匮乏与贫穷的状态，他视自己为这一切的外人和异乡人，以陌生人及寄居者的身份行于世上，并将自己视为隐修院的学徒和仆人，而不是幻想自己是任何事物的主人和主宰。\par

% structure: p0044 source=h2 target=subsection reason=relative-heading-rank; base=h2

\subsection{第十五章}

论我们中间过度的占有欲。\par

对此，我们这些可怜虫还能说什么呢？我们虽住在修院，在院长的治理和照管之下，却随身携带自己的钥匙，践踏因我们修道誓愿而应有的羞耻感，竟然毫不羞耻地公然在手指上戴着印章戒指来封存我们囤积之物；在我们当中，不仅箱子和篮子，甚至柜子和橱柜都不足以装下我们收集或弃绝世俗时保留的东西；有时我们为琐碎和微不足道之物（但我们却视为己有）如此动怒，以至于若有人胆敢碰触其中任何一件，我们就对他满腔愤怒，无法控制内心的怒火，不让它表露在嘴唇和身体的激动上。但是，且放过我们的过失，闭口不谈那些连提及都羞耻的事，正如这句经文所说：\Quote{我的口不谈论人的作为。}让我们按照已开始的叙述方式，继续论述他们中间所实践的那些德行，我们应当以全部热忱追求这些德行；让我们简要而迅速地记下实际的规则和体系，以便随后在详细记述长老们的某些言行（我们打算仔细保存以供回忆）时，能以最有力的证据支持我们在论著中所阐述的内容，并通过生活中的事例和实例进一步证实我们所说的一切。\par

% structure: p0047 source=h2 target=subsection reason=relative-heading-rank; base=h2

\subsection{第十六章}

论各种斥责的规则。\par

若有人不慎打碎了一只陶罐（他们称之为\Quote{\OriginalTerm{陶罐}{baucalis}}），他只能藉公开补赎来补偿其疏忽；当全体弟兄聚集举行礼仪时，他必须躺在地上求赦罪，直到祈祷礼仪结束；并将在院长的命令下被吩咐从地上起来时获得赦免。同样，若有人应召去工作或参加常规礼仪时来得太迟，或在咏唱圣咏时稍有迟疑，也必须做同样的补赎。同样，若他回答得不必要、粗鲁或无礼，若他在执行所托付的服务时疏忽大意，若他稍作抱怨，若他偏好阅读甚于工作或服从，从而缓慢履行指定的职责，若礼仪结束后他不立即赶回自己的独室，若他与别人一起停留片刻，哪怕极短，若他拉住别人的手，若他敢于与不与他同住一室的人讨论任何微小的事情，若他与一个被剥夺祈祷资格的人一起祈祷，若他在未经其长者允许的情况下见到任何世上的亲友并与他们交谈，若他未经院长许可试图从任何人那里收信或回信。精神上的惩戒竟如此严厉，并涉及这类事情和类似的过失。但至于其他事情，在我们中间若随意犯下也被视为应受责备的，即：公开争吵、明显的轻蔑、傲慢的顶撞、随意无节制地离开修院、与妇女熟悉、愤怒、争吵、嫉妒、争论、将某物据为己有、贪恋的感染、欲望及获取并不需要的（其他弟兄所没有的）东西、两餐之间私下进食，以及诸如此类——它们不是通过我们所说的那种精神上的惩戒，而是通过鞭打来处置；或通过驱逐来弥补。\par

% structure: p0050 source=h2 target=subsection reason=relative-heading-rank; base=h2

\subsection{第十七章}

\Emph{关于那些引入在弟兄们用餐时在\OriginalTerm{修院}{Cœnobia}中诵读圣课之计划的人，以及关于埃及人中所遵守的严格缄默。}\par

但据我们所知，弟兄们用餐时在修院中诵读圣课的计划并非源于埃及制度，而是源于卡帕多细亚制度。无疑，他们建立此计划与其说是为了属灵操练，不如说是为了制止餐桌上常常出现的无益闲谈，尤其是争论；因为他们看到，除此之外无法阻止这些事。因为在埃及人，尤其是\OriginalTerm{塔本纳}{Tabenna}的隐修士中，所有人遵守如此严格的缄默，以至于当一大群弟兄一起坐下来用餐时，除监院外，无人敢低声交谈；监院若看到需要添菜或撤盘，则以手势而非言语示意。用餐时，他们如此严格地遵守这一缄默规则，以致将头巾拉下遮住眼睑（以防目光游移、好奇张望），除餐桌、桌上所放及所取的食物外，什么也看不见；因此无人注意到别人在吃什么。\par

% structure: p0053 source=h2 target=subsection reason=relative-heading-rank; base=h2

\subsection{第十八章}

\Emph{任何人在公共餐桌之外取用任何饮食如何违反规则。}\par

在常规的共同用餐之间，他们特别谨慎，不准任何人擅自以任何食物满足口腹之欲：因此，当他们随意穿过花园或果园时，树上诱人悬挂的果实不仅擦过他们经过的胸膛，还落在地上任人践踏，而且（由于果实已成熟可摘）很容易引诱看见的人满足食欲，并因眼前的机遇和果实之多，甚至激起最严厉、最节制的人渴望之；然而，他们认为，除了公开放在桌上供全体共同用餐、由管事供应并经弟兄服务共享的食物外，不仅品尝一个果子是错的，连用手碰一下也是错的。\par

% structure: p0056 source=h2 target=subsection reason=relative-heading-rank; base=h2

\subsection{第十九章}

\Emph{在\Place{巴勒斯坦}和\Place{美索不达米亚}各地，弟兄们如何每日履行礼仪。}\par

为了不使我们似乎遗漏了\OriginalTerm{隐修院}{Cœnobium}的任何规章，我认为应简要提及，在其他地区，弟兄们也每日轮流承担服务。因为在整个美索不达米亚、巴勒斯坦、卡帕多细亚以及整个东方，弟兄们每周轮流承担某些职责，服务人数根据隐修院中全体隐修士的数量而定。他们以奴仆对最严厉强大的主人也未曾有的热忱和谦卑，急切地履行这些职责；以至于不满足于仅按会规提供的服务，竟在夜间因热忱而起，取代专门负责的弟兄；并秘密地抢在他们前面，试图完成别人本应承担的工作。每一位承担当周值班者，须服务至主日（星期日）晚餐，此后其整周职责即告完成。因此，当全体弟兄聚集咏唱圣咏（按习惯在就寝前咏唱）时，值班期满者轮流为众人洗脚，真诚地以此为自己整周工作祈求祝福的赏报，即全体弟兄的共同祈祷伴随他们，以完成基督的命令，即那为他们的无知和因人性软弱所犯之罪代祷的祈祷，并将他们虔诚的完全服务像丰厚祭品一样荐于天主。于是，在星期一早课赞美诗后，他们将用于服务的器皿和用具交给接替者；后者极其谨慎地接收和保管，唯恐任何一件损坏或丢失，因为他们相信，即使是最小的器皿，也须作为圣物交代，不仅向现任管库，而且向主交代，若有任何因疏忽而损坏者。这种纪律的限度，及其遵守中的忠信与审慎，可由我将举的一个实例看出。因为我们急于满足您透过要求详尽说明一切并在本论著中重复您已知晓之事所表现的热忱，同时也担心超出简洁的界限。\par

% structure: p0059 source=h2 target=subsection reason=relative-heading-rank; base=h2

\subsection{第二十章。}

\Emph{关于管家发现的三颗扁豆。}\par

在某位弟兄当周期间，管家经过时看见地上有三颗扁豆，是从当周值班修士手中滑落的。当时他正匆忙准备煮食，连同洗豆的水一起洒落。管家立即就此事请示院长；院长判定该修士为偷窃者且对圣物疏忽，遂暂停其祈祷。其疏忽之过，只在以公开补赎赎罪后才得赦免。因为他们相信，不仅他们自己不属自己，而且他们所拥有的一切皆奉献于主。因此，任何物品一旦进入隐修院，他们便认为应以最大敬意视之为圣物。即使对待被认为微不足道或视为平常与琐碎之物，他们也极其忠信地照料安排；例如，若将它们移至更佳位置，或给水瓶装水，或给人从中饮水，或从祈祷所或小室拂去一点灰尘，他们都深信必从主获得赏报。\par

% structure: p0062 source=h2 target=subsection reason=relative-heading-rank; base=h2

\subsection{第二十一章。}

\Emph{论某些弟兄的自发服务。}\par

我们听说有些弟兄，他们当周时木材极度匮乏，不足以准备弟兄们的常规食物；院长依权责下令，在更多木材运来之前，他们应以干粮为足。尽管众人皆同意，无人可指望熟食；但这些人，仿佛因未照惯例在轮值次序中为弟兄准备食物而被剥夺了劳苦与服务的果实与赏报——竟将如此不必要的劳苦与照顾加诸自身：在那些干燥贫瘠、除果树枝外无法获得木材的地区（因为那里不像我们这里有野生灌木），他们穿越广阔荒漠，跋涉向死海延伸的荒野，在衣襟和衣褶中收集风吹来的稀疏残梗与荆棘，如此凭自愿服务为弟兄们准备全部常规食物，使日常供应毫无减少。他们以如此忠信履行对弟兄的职责，尽管木材匮乏和院长的命令给予他们公平的借口，但为顾及自身的利益与赏报，仍不愿利用这自由。\par

% structure: p0065 source=h2 target=subsection reason=relative-heading-rank; base=h2

\subsection{第二十二章。}

\Emph{埃及人为弟兄们每日服务所指定的制度。}\par

这些事是按照我们先前所说整个东方的制度讲述的，我们也说它在我们本国自当遵守。但在埃及人当中，他们主要关心工作，并无互相轮换的每周服务，以免因职务之需而妨碍遵守工作规则。他们从最受认可的弟兄中选出一人，负责储藏室和厨房，且只要体力与年岁允许，永久掌管此职。因为他不受巨大体力劳动的劳累，因他们准备食物或烹煮不花费太多心思，多使用干粮与生食；在他们当中，每月切下的韭葱叶、山芥、食盐、橄榄、称为沙丁鱼的小咸鱼，构成最大的美味。\par

% structure: p0068 source=h2 target=subsection reason=relative-heading-rank; base=h2

\subsection{第二十三章。}

\Emph{院长\Person{若望}的服从，藉此他得以被高举甚至达到预言的恩宠。}\par

既然本书关於训练一位弃绝此世者，使他藉由开始真正的谦卑和完美的服从，也能攀升至其他诸德的高峰；我以为最好仅作为样本，按我们应许的，记下一些长老们在其中出类拔萃的事迹，从许多事例中选取少数；若有人渴望追求更高的境界，他们不仅由此得到朝向完美生活的激励，还能获得所立志的模范。因此，为使本书尽可能简短，我们将从全部教父中举出两三位；首先是院长\Person{若望}，他住在\Place{特拜德}的一个城镇\Place{吕孔}附近；因其令人钦佩的服从，被高举甚至达到预言的恩宠，且名满天下，甚至在这世界的君王中也因他的功德而闻名。因为，如我们所说，他住在\Place{特拜德}最偏远的地区，但\Person{德奥多西}皇帝若不先受他的言辞和答复鼓励，都不敢向最强大的暴君宣战；皇帝信赖这些，如同从天而降，在看似无望的战役中取胜于敌人。\par

% structure: p0071 source=h2 target=subsection reason=relative-heading-rank; base=h2

\subsection{\Emph{第二十四章}}

\Emph{院长\Person{若望}遵其长上之命，长期浇灌一根枯枝，彷彿它能生长。}\par

这位有福的\Person{若望}，从青年直到成熟的高龄，一直顺服他的长上，只要他仍活在此世；他如此谦卑地执行命令，以致长上本人也对他的服从感到完全惊讶。长上欲确证这德行是来自真诚的信德和纯朴的心，还是装出来的、勉强的、只在发令者面前才显示；于是他常给他许多多余、几乎不必要、甚至不可能的命令。我将从中选出三件，向愿知道他的心志与顺服何等完美的人展示。因为老人从柴堆里取出一根先前已砍好、预备生火的木棍，因无机会煮饭，已因时日而不仅乾燥，而且发霉。他在他眼前把木棍插进土里，命他取水每日浇灌两次，好藉此每日浇灌，使它生根复活如先前的树，伸展枝桠，既悦目，又为在盛夏炎热中坐于其下者提供荫凉。这青年以惯常的恭敬领受了命令，从未考虑其不可能，日复一日执行，几乎每天需提水步行两哩，从未停止浇灌木棍；整整一年，没有身体疾病、节日服务、必要的事务（这些都合理可免除他执行命令），最后也没有寒冬的严酷能妨碍他遵行此令。老人暗中观察了几日，见他以纯朴的心心甘情愿地持守命令，仿佛来自天上，面色不变，亦不考量其合理性——赞许他谦卑中不虚假的服从，同时怜悯他因热心忠诚而持续一年的辛劳——便来到枯枝前，说：\Quote{\Person{若望}，这棵树生根了没有？}他回答说不知道。于是老人彷佛探求实情，试探它是否已扎根，便当他的面轻轻摇动泥土拔起木棍，扔在一旁，命他今后可停止浇灌。\par

% structure: p0074 source=h2 target=subsection reason=relative-heading-rank; base=h2

\subsection{\Emph{第二十五章}}

\Emph{院长\Person{若望}遵其长上之命，丢弃唯一一瓶油。}\par

这样，青年经过此类练习，日益增进服从之德，并以谦卑的恩宠更加显耀。当他服从的美名传遍所有隐修院时，一些弟兄为考验他，或更为使他受益，前来见那长者，惊叹于他们所闻的服从。于是长者突然叫他，说：\Quote{上去，拿这瓶油（这是沙漠中唯一的一瓶，为数不多的优质液体，供他们自己和访客使用），把它从窗户扔出去。}他闻命立刻飞步上楼，将油瓶从窗户扔下地面摔碎，丝毫未曾思考命令的愚蠢、每日所需、身体的软弱、贫穷，以及这可怜沙漠的艰苦与困难；在这样的沙漠中，即使有钱，一旦失去，也无法购得或替代同样品质的油。\par

% structure: p0077 source=h2 target=subsection reason=relative-heading-rank; base=h2

\subsection{\Emph{第二十六章}}

\Emph{院长\Person{若望}如何顺从长上，试图滚动一块众人都无法移动的巨石。}\par

又有一次，当另一些人渴望从服从的榜样中得到益处时，长者叫他，说：\Quote{\Person{若望}，跑去，尽快把那块石头滚到这儿来。}他立刻时而用颈，时而用全身，倾尽全力试图滚动一块巨石，那是一大群人也无法移动的；以致不仅衣服被汗水湿透，连石头也被他的颈项润湿了；在此事上他也从未衡量命令与行为的可能性，出于对长上的尊敬和事奉的纯朴真诚，他坚信长者不会命令他做任何虚空无理之事。\par

% structure: p0080 source=h2 target=subsection reason=relative-heading-rank; base=h2

\subsection{\Emph{第二十七章}}

\Emph{院长\Person{帕特穆齐乌斯}的谦卑与服从，他毫不犹豫地遵长上之命，将自己的幼子投入河中。}\par

以上我略举数事论及院长\Person{若望}，尚足为证；今我将述院长\Person{帕特穆修}一件可纪念的事迹。原来他切愿弃绝世俗，长久躺在隐修院门前，直到他凭借顽强的坚持，违反隐修院的一切规定，说服他们收留他和他的幼子——那孩子大约八岁。及至最终获准入院，他们立刻不仅被交给不同的上司照管，而且还被安置住在分隔的静室中，以免父亲因不断看见幼子而想起：在他已抛弃并舍弃的一切财产和世俗财富中，至少他的儿子仍留给他；并且，既然他已被教导不再是一个富人，他或许也会忘记自己是一个父亲的事实。为了更彻底地考验他是否会将他对自己骨肉的情感与爱看得比服从和基督徒的克苦（所有弃绝世俗者，因爱基督之故，都当优先选择克苦）更重，那孩子故意被忽视，穿着破衣而非合宜的衣服；并被泥土弄得如此污秽丑怪，以致每当父亲看见他时，只会感到厌恶而非喜悦。此外，他还遭受来自众人的拳打脚踢，父亲经常亲眼看见这些无缘无故加在他无辜孩子身上的殴打，以至于他每次看见孩子的面颊，都沾满了泪水的污痕。尽管孩子日复一日在他眼前这样受虐待，但出于对基督的爱和服从的德行，父亲的心依然坚定不动摇。因为他已不再视他为自己的儿子，因为他已将他与自己一同奉献给了基督；他也不关心他当前的伤害，反而欢欣，因为他看到这些伤害不是没有益处地忍受；他看轻儿子的眼泪，却挂虑自己的谦逊和成全。当隐修院院长见到他心思的坚毅与不可动摇的坚决，为彻底证明他志向的恒常，有一天看见孩子哭泣，便装作恼恨他，命令父亲将他扔进河里。于是，他仿佛这是主所命令的，立刻以最快的速度抱起孩子，抱在怀中带到河岸要将他扔进去。随即，凭着他信德与服从的热忱，这事本来就要在行动上实现，若非有几位弟兄预先受命严密监视河岸，并在孩子被扔下时，设法从河床中将他夺回，阻止了那条命令——这命令其实已因父亲的服从与奉献而在动机上完成——在行动和结果上落实。\par

% structure: p0083 source=h2 target=subsection reason=relative-heading-rank; base=h2

\subsection{第二十八章。}

\Emph{关于\Person{帕特穆修}，如何启示给院长他行了\Person{亚巴郎}的事迹；以及当同一院长去世后，\Person{帕特穆修}如何继任管理隐修院。}\par

这个人的信德与奉献如此中悦天主，以致立即得到了神圣见证的认可。因为立即启示给院长，他因这一服从行为，效法了圣祖\Person{亚巴郎}的事迹。不久之后，同一隐修院院长离世归主，他将他置于众弟兄之上，并指定他为继承人和隐修院院长。\par

% structure: p0086 source=h2 target=subsection reason=relative-heading-rank; base=h2

\subsection{第二十九章。}

\Emph{论一位弟兄的服从，他奉院长之命，在公开场合携带十只篮子并零卖它们。}\par

我们也不愿缄默一位我们所认识的弟兄，他按这世界的阶层属于高门，因为他出身于一位伯爵父亲，极其富有，并曾受过自由教育的良好教养。此人离开父母逃往隐修院后，为证明其性情的谦逊与信德的热忱，立即被院长命令以肩负起十只篮子（其实并无必要公开出售），并沿街叫卖：附加此条件，为使他更久留于此工，即：若有人偶然想一次全部购买，他不得允许，而应分别卖给购者。他极其热诚地执行了这事，践踏一切羞愧与慌乱，出于对基督的爱，并为祂圣名的缘故，他将篮子搁在肩上，按定价零卖，将银钱带回隐修院；丝毫不因如此卑贱不寻常之任务的奇特性而烦恼，也不介意此事的屈辱、他出生的高贵和出售的丢脸，因为他一心追求通过服从的恩宠获得基督的谦逊，那才是真正的尊贵。\par

% structure: p0089 source=h2 target=subsection reason=relative-heading-rank; base=h2

\subsection{第三十章。}

\Emph{论院长\Person{皮努菲乌}的谦逊，他离开一所他担任司铎主持的极著名隐修院，出于对服从之爱的追求，前往一所遥远的隐修院，好能获收留为初学生。}\par

书籍篇幅的限制迫使我们收尾；然而，在诸善德中居于首位的服从之德，却不允许我们将那些以之超群出众者的行径完全略过。因此，将两者巧妙结合（即兼顾简洁以及恳切者的愿望与益处），我们只添加一个谦逊的榜样，它既非由新手，而是由一位已臻完善且身为院长的人所展现，当读者阅读时，不仅可教导较年轻的，也能激励年长者臻于谦逊的完美之德。我们曾目睹院长\Person{品纽弗}，他曾是\Place{埃及}一处巨大隐修院的司铎，该隐修院距\Place{巴内弗西斯}城不远，因他的一生、年岁或司铎职之故，众人皆以敬意与尊崇待他。当他察觉正因如此，他无法实践自己以全部热诚所渴望的谦逊，也没有机会操练他所切望的服从之德时，他便秘密逃离隐修院，独自退往\Place{特拜德}最偏远的地区，在那里脱下隐修士的会衣，换上世俗服装，这样便投奔\Place{塔本纳}隐修院——他深知那是所有隐修院中最严格的一所，并以为或因地点的遥远而不会被人认出，或可因隐修院的规模与弟兄的人数而轻易隐藏在那里。他在入口处停留了很久，如同恳求者般跪在弟兄膝前，以最恳切的祈祷寻求获准入院。当他最终因被视为一个终生在世度日、年老力衰，且在老年时请求进入隐修院（因他们称他寻求此并非出于宗教，而是受饥饿与匮乏所迫）的老者而被勉强接纳后，他们派他照管和管理花园，因他看似年长且不特别适合某项工作。他在另一位较为年轻的弟兄手下从事此务，那位弟兄将他当作受托付者留在身边；他如此顺服于他，并以如此服从的态度培植所渴望的谦逊之德，以至于他每日不仅极其勤勉地完成与花园照管和管理相关的一切，而且还执行那些被他人视为艰辛、卑贱且令人不快的所有职责。他也夜间起身，暗中做许多事，无人观看也无人知晓，黑暗将他隐藏，无人能查明事主。当他藏身那里三年，并被弟兄们走遍埃及上下寻找之后，终于被一位来自埃及地区的人看见，但因其衣着低微和所从事的卑微职责几乎无法认出。他正弯腰锄地为蔬菜松土，并用肩膀背负粪肥放在它们的根部。看见他的弟兄迟疑许久才认出他，但最终走近，仔细留意其面容及语调，立即扑倒在他脚前。起初，所有看见这事的人都极其惊讶，为何他对一个被他们视为最卑微者（犹如新手且刚刚弃世者）如此行；但此后，当他们立刻听到他报出姓名——那个甚至在他们中间也早已藉声誉闻名者时，就更加震惊。众弟兄因先前的无知而请求他原谅，因他们长久以来将他与年幼者和孩童归为一类，便将他带回他自己的隐修院，他虽泪流满面且不情愿，因魔鬼的嫉妒骗走了他可敬的生活方式和那久寻后欣喜发现的谦逊，也因他未能在所得的那服从状态中结束此生。于是他们极其谨慎地看守他，以防他再以同样方式溜走并逃脱他们。\par

% structure: p0092 source=h2 target=subsection reason=relative-heading-rank; base=h2

\subsection{第三十一章}

\Emph{院长\Person{品纽弗}被带回他的隐修院后，在那里停留了片刻，便又逃往\Place{叙利亚巴勒斯坦}地区。}\par

他在那里停了片刻，又被对谦逊的渴望与向往所抓住，趁夜寂静逃出，这次他不再寻邻近地区，而是前往陌生、遥远且相隔广阔的地区。他搭船设法到了\Place{巴勒斯坦}，相信若前往那些从未听闻他名字的地方，会更安全地隐藏。他到了那里，即刻寻访我们的隐修院，它离我们主曾由童贞女诞生的山洞不远。虽然他在这里隐藏了一段时间，却如同（用我们主的话说）\BibleQuote{建在山上的城}\BibleRef{玛 5:14}，不能长久隐藏。不久，一些从\Place{埃及}前来圣地祈祷的弟兄认出他，并以最热切的祈祷将他召回他自己的隐修院。\par

% structure: p0095 source=h2 target=subsection reason=relative-heading-rank; base=h2

\subsection{第三十二章}

\Emph{同一院长\Person{品纽弗}在我们面前接纳一位弟兄进入隐修院时给他的训导。}\par

这位年长者后来我们在埃及努力寻访，因我们曾在自己的隐修院与他亲近；我打算在这次著作中插入一段劝勉，那是他在我们面前给一位被他接纳进入隐修院的弟兄的，因我认为这或许有益。他说：\Quote{你知道，在门口躺卧多日之后，你今天将被接纳。首先，你应当知道设置困难的原因。因为在这条你渴望进入的路上，你若明白其方法，并据此适当地接近基督的服役，这或许对你大有裨益。}\par

% structure: p0098 source=h2 target=subsection reason=relative-heading-rank; base=h2

\subsection{第三十三章}

\Emph{如何正如按照教父规条勤劳的隐修士应得大赏报，懒惰的隐修士也当受惩罚；因此不应轻易接纳任何人进入隐修院。}\par

因为正如按照这一制度的规则忠信事奉天主、紧依天主的人，将来有无穷的荣耀被应许；同样，那些粗心冷淡地履行、未能向祂呈上与他们所宣认或人们以为他们所是相称的圣德果实的人，有最严厉的惩罚为他们存留。因为圣经说：\Quote{人许愿不还，不如不许愿；}又说：\Quote{懒惰行主工的可咒骂。}因此，我们长久拒绝你，并非好像我们不从心里渴望确保你和所有人的得救，也不是好像我们不关心去迎接甚至远处那些渴望归向基督的人；而是恐怕我们轻率接纳你，使我们在天主面前犯轻浮之罪，并使你陷入更重的惩罚，假如你被我们过于轻易地接纳，却没有认识到这一修道的责任，后来竟成了逃兵或冷淡者。因此，你首先应当明白弃绝世界的实际原因，当你明白了这一点，就可以从那原因更清楚地被告知你该做什么。\par

% structure: p0101 source=h2 target=subsection reason=relative-heading-rank; base=h2

\subsection{第三十四章。}

\Emph{论我们的弃绝不外乎是死罸与基督被钉之像。}\par

\Emph{弃绝}不外乎是十字架与死罸的明证。因此你必须知道，今天你已死于这世界及其作为和欲望，并且，如宗徒所说，你已被钉于这世界，这世界也被钉于你。\BibleRef{迦 6:14}因此你要思想十字架的要求，你今后应在今生生活于这十字架的标记之下；因为\Emph{你}不再生活，而是那为你被钉的\Emph{祂}在你内生活。\BibleRef{迦 2:20}我们必须以那种形式和样式在今生度过我们的时间，就是祂为我们被钉在十字架上的样式，以便（如达味所说）用敬畏主刺透我们的肉身，使我们的一切意愿和欲望不服从于我们自己的私欲，而是被钉于祂的死罸。因为这样我们才能实现主的命令，祂说：\Quote{不背起自己的十字架跟随我的，不配是我的。}\BibleRef{玛 10:38}但或许你会说：人怎么能不断背起自己的十字架？或者一个活着的人怎么能被钉死？请简短地听这是如何的。\par

% structure: p0104 source=h2 target=subsection reason=relative-heading-rank; base=h2

\subsection{第三十五章。}

\Emph{论敬畏主如何是我们的十字架。}\par

敬畏主就是我们的十字架。正如被钉的人不能再随意向任何方向移动或转动肢体，我们也当将我们的意愿和欲望——不依照现在愉悦和快乐的事物，而依照主的律法约束我们之处——钉住。又如被钉在十字架木头上的人不再顾及现在的事情，也不想自己的喜好，不为明天的焦虑和挂虑所困，不为任何拥有的欲望所扰，不为任何骄傲、争斗或竞争所激动，不因当前的伤害而悲伤，不记念过去的，并且还在肉身中呼吸时就认为自己已死于一切属地之事，把心灵的思想提前送到他确信不久就要去的地方；同样，当我们被敬畏主所钉时，也当确实死于这一切，即不仅死于肉身的恶习，也死于一切属地之事，将我们心灵的眼目定在那里，即我们时刻希望不久就要经过的地方。因为这样我们就能使一切欲望和肉体的情感死去。\par

% structure: p0107 source=h2 target=subsection reason=relative-heading-rank; base=h2

\subsection{第三十六章。}

\Emph{论我们的弃绝若再次被所弃绝之事缠住，便毫无用处。}\par

因此，你要谨慎，免得任何时候你再拿起你所弃绝和抛下的任何东西，违背主的命令，从福音工作的田地返回，并发现自己重新穿上你已脱下的外衣；\BibleRef{玛 24:18}也不要沉回到这世界卑下和属地的情欲与欲望中，并违背基督的话，从完美的棍杖下来，胆敢再拿起你已弃绝和抛下的任何东西。要谨慎，不要记念你的亲属或你以前的情感，也不要被这世界的挂虑和焦虑召回去，（如我主所说）\Quote{手扶着犁向后看，以致被认为不配进天国。}\BibleRef{路 9:62}要谨慎，免得任何时候你开始涉猎圣咏和这个生活的知识，就渐渐自高自大，想要复兴那骄傲——就是你在开始时凭信德的热忱和极度的谦卑所践踏的；并因此（如宗徒所说）\Quote{重建你所拆毁的，使自己成为后退者。}\BibleRef{迦 2:18}反而应当注意，要持续到底，在你于天主和祂的天使面前所宣认的那种赤裸状态中。同样，在谦卑和忍耐中——你在门前坚持了十天，流了许多眼泪恳求被接纳进入隐修院——你不仅应当继续，还要增加并前进。因为当你应当从初学和开端被带领，前进到成全时，却开始从这些退回到更坏的事情，那就太糟了。因为开始的不是得救的，而是坚持到终点的才能得救。\BibleRef{玛 24:13}\par

% structure: p0110 source=h2 target=subsection reason=relative-heading-rank; base=h2

\subsection{第三十七章。}

\Emph{论魔鬼如何总是伺机等我们的终结，以及我们应如何时刻注视他的头。}\par

因为狡猾的古蛇常常\Quote{窥伺我们的脚跟}，即埋伏在末了，企图在我们生命终了时绊倒我们。因此，仅仅有好的开始并热切迈出第一步，以满腔热忱弃绝世界，是无用的，除非相应的结局也与之相配并完成它，并且你们如今在祂面前所宣认的基督的谦卑与贫穷，如同最初被确保的那样，由你们持守到底，直至生命终了。为使你们能成功做到这点，你们要经常\Quote{注视它的头}，即思想的初次萌动，立即将它们呈报给长上。因为如果你们不耻于向长上揭露其中任何一点，你们就将学习\Quote{踏碎}它危险的开端。\par

% structure: p0113 source=h2 target=subsection reason=relative-heading-rank; base=h2

\subsection{第三十八章}

\Emph{论\OriginalTerm{弃绝者}{renunciant}对诱惑的准备，以及少数值得效法的人。}\par

\Emph{因此}，如圣经所说：\Quote{你去事奉主时，当存敬畏主的心，预备你的心}\BibleRef{德 2:1}，不是为了安息、疏忽或快乐，而是为了诱惑和磨难。因为\Quote{我们必须经历许多苦难，才能进入天主的国。}因为\Quote{门是窄的，路是狭的，通往生命，找到的人很少。}因此，你要认定自己属于那少数被拣选的人；不要因多数人的冷淡榜样而冷却；却要像少数人那样生活，好能与少数人一起在天主的国里获得一席之地；因为\Quote{被召的多，但被选的少。}并且\Quote{你们这小群，父乐意把国赐给你们。}因此，你应当认识到，一个宣认了\OriginalTerm{成全}{perfection}的人却追求不完全，这不是小罪。你要通过以下步骤和方式，达到这个成全的境界。\par

% structure: p0116 source=h2 target=subsection reason=relative-heading-rank; base=h2

\subsection{第三十九章}

\Emph{论我们如何攀登\OriginalTerm{成全}{perfection}之路，并由此从\OriginalTerm{敬畏天主}{fear of God}上升至爱德。}\par

\Quote{开端}我们得救并确保它的是，如我所说，\Quote{敬畏主}。\BibleRef{箴 9:10} 因为藉着它，那些在成全道路上受训的人，得以开始悔改、净化邪恶并在美德中得保安全。当这进入人心时，它产生对一切事物的轻视，并生起忘却亲属和对世界本身的恐惧。但因轻视损失一切所有，便获得了\OriginalTerm{谦卑}{humility}。谦卑由这些迹象证实：首先，如果一个人克制了自己的一切欲望；第二，如果他毫不隐瞒自己的行为乃至思想，不向长上呈报；第三，如果他不信赖自己的意见，而完全信赖长上的判断，并热切而乐意地听从长上的指引；第四，如果再一切事上持守服从、温和和恒久的耐心；第五，如果他不仅不伤害任何人，而且对别人加给他的伤害也不感到恼怒或烦恼；第六，如果他不做也不尝试任何未经大众规则或长上榜样所敦促的事；第七，如果他满足于最低可能的位置，并视自己为拙劣的工人，对所吩咐的一切觉得自己不配；第八，如果他不只在嘴上承认自己比众人低下，而且在内心深处真正相信；第九，如果他管束自己的舌头，不过多言语；第十，如果他不易激动或过于好发笑。因为通过这些以及类似的迹象，能认出真正的谦卑。一旦这谦卑真正获得，它立刻便引领你藉着更高的一步，达到那没有恐惧的\OriginalTerm{爱德}{love}；\BibleRef{若一 4:18} 藉着这爱德，你开始毫不费力，并仿佛自然地，持守所有从前你因害怕惩罚而遵守的一切；不再出于对惩罚的顾虑或恐惧，而是出于对美善本身的爱和对德行的喜悦。\par

% structure: p0119 source=h2 target=subsection reason=relative-heading-rank; base=h2

\subsection{第四十章}

\Emph{论\OriginalTerm{隐修士}{monk}不应从许多榜样，而应从一个或极少数榜样寻求\OriginalTerm{成全}{perfection}的典范。}\par

为使你们更容易达到这一点，应只从极少数，甚至只从一个或两个榜样，而不是从过多榜样，寻求那值得模仿的、住在修院团体中的成全生活的典范。因为除了经过考验、炼净和纯化的生活只能在少数人身上找到之外，还有这一点益处，即：一个人通过某\Emph{一}个榜样的影响，能更彻底地受教并被塑造，达到他所追求的成全，即\OriginalTerm{修院团体}{Cœnobite life}生活的成全。\par

% structure: p0122 source=h2 target=subsection reason=relative-heading-rank; base=h2

\subsection{第四十一章}

\Emph{论住在\OriginalTerm{修院团体}{Cœnobium}中应显现怎样的\OriginalTerm{软弱}{infirmities}。}\par

为使你能达到这一切，并持续服从这属灵的规则，你必须在会众中遵守这三件事：即如圣咏作者所说：\Quote{我如同聋子，不听见；如同哑巴，不开口；我成了一个不听的人，口中没有反驳。} 因此，你也应当像聋子、哑巴和瞎子一样行走——放下你被正确选为完美模范的那一位的默观——你应像瞎子一样，看不见任何你发现无益的事物，也不受那些做这些事的人的权威或风尚的影响，而屈服于更坏的事和你以前所谴责的事。如果你听见有人不服从、不顺从、诋毁别人，或做了任何与你所学不同的事，你不应犯错并被这样的榜样误导去效仿他；反而，\Quote{如同聋子}，仿佛从未听过一样，一概放过。如果有人侮辱你或别人，或者做了错事，你要坚定不移，就报复的回答而言，要缄默\Quote{如同哑巴}，常在你心中吟诵圣咏作者的这句诗：\Quote{我说：我要谨遵我的道路，免得我以口舌犯罪。当罪人站在我面前时，我给我的口设了守卫。我成了哑巴，卑微，并沉默不语，连善事也不说。} 但最重要的是培养这第四件事，它装饰并使上述三件事优美：即你要如同宗徒所教导的，\BibleRef{格前 3:18}，在这世界成为愚笨，好使你成为有智慧的，在一切命令你的事上不行自己的辨别和判断，而总是在一切纯朴和信德中表示服从，认为唯有天主的法律或你的长上的决定所宣告的才是圣洁、有益和智慧的。因为建立在这样的教导体系上，你可以永远持续在这纪律之下，不会因任何仇敌的诱惑或诡计而从隐修院堕落。\par

% structure: p0125 source=h2 target=subsection reason=relative-heading-rank; base=h2

\subsection{第四十二章。}

\Emph{隐修士不应从别人的德行中寻求自己得忍耐的祝福，而应从自己的含忍中得来。}\par

因此，你不应从别人的德行中为自己寻求忍耐，以为只有在不受任何事物激怒时才能获得它（因为防止此事发生并不在你的能力之内）；而应将它视为你自己的谦卑和含忍的后果，这\Emph{确实}取决于你自己的意志。\par

% structure: p0128 source=h2 target=subsection reason=relative-heading-rank; base=h2

\subsection{第四十三章。}

\Emph{隐修士如何攀登至完美的说明之综述。}\par

为使以上这些以稍长篇幅阐述的事情更易铭刻在你心中，并牢牢粘附于你的思想，我将作一总结，使你因其简短扼要而能熟记所有这些变化。那么，请听几句简短的话，你如何能毫不费力地攀登至完美的高峰。依照圣经，我们救恩的\Quote{开端}和\Quote{智慧}是\Quote{敬畏上主}。从敬畏上主产生有益的痛悔。从心灵的痛悔生出弃绝，即赤贫和轻视一切财物。从赤贫生出谦卑；从谦卑生出克治欲望。通过克治欲望，一切过犯被根除而消灭。通过驱除过犯，美德勃发并增长。通过美德的萌芽，获得内心的纯洁。通过内心的纯洁，获得宗徒之爱的完美。\par

\SourceRecord{Translated by C.S. Gibson.；Nicene and Post-Nicene Fathers, Second Series；Vol. 11.；Edited by Philip Schaff and Henry Wace.；Buffalo, NY: Christian Literature Publishing Co.,；1894.；<http://www.newadvent.org/fathers/350704.htm>.}

% structure: p0001 source=h1 target=section reason=page-title

\section{\WorkTitle{修院规则（第五卷）}}

% structure: p0002 source=h2 target=subsection reason=relative-heading-rank; base=h2

\subsection{第一章}

\Emph{从隐修士的规则到对抗八种主要恶性的过渡。}\par

我们这第五卷书，如今赖天主的助佑，即将诞生。因为在已完成的论隐修院习俗的四卷书之后，我们现今在天主藉你们的祈祷所赐的力量下，打算开始对抗八种主要恶性的战斗，即：第一，贪饕或口腹之乐；第二，邪淫；第三，悭吝，即贪婪，或更恰当的说法是爱财；第四，忿怒；第五，忧愁；第六，\Quote{懈怠}，即心中的沉重或厌倦；第七，\OriginalTerm{虚荣}{\GreekText{κενοδοξία}}，即愚妄或虚浮的荣耀；第八，骄傲。在开始这项艰巨任务时，至福的教宗\Person{卡斯托尔}，我们比以往更需要你们的祈祷，好使我们首先能够值得探究这些恶性的本质，无论多么琐细、隐秘或晦涩；其次充分清楚地解释它们的原因；第三恰当地提出它们的疗法和补救。\par

% structure: p0005 source=h2 target=subsection reason=relative-heading-rank; base=h2

\subsection{第二章}

\Emph{这些恶性的缘由虽存在于每个人身上，却被每个人忽略；我们需要主的帮助来显明它们。}\par

这些情欲的缘由，一旦藉长老的教导揭露出来，便为众人所认识，但在它们被揭示以前，尽管我们都被它们征服，它们存在于每个人身上，却无人知晓。但我们相信，若因你们的祈祷，那藉\Person{依撒意亚}所宣告的主的话也能应验在我们身上——\BibleQuote{我必在你前面行，打压地上的强豪，打破铜门，砍断铁闩，将隐藏的宝贝和隐密的宝藏赐给你}\BibleRef{依 45:2}——好使主的话也走在我们前面，首先打压我们本地的强豪，即这些我们渴望战胜的邪恶情欲，它们在我们的必死肉身中自称主宰并以最可怕的暴政统治；并使它们屈服于我们的探究和阐述，从而打破我们无知的铜门，砍断将我们隔绝于真知的罪恶铁闩，引领我们进入自身隐秘的深处，并光照我们，按宗徒的话，\BibleQuote{揭露暗中的隐情，显明心中的意念}\BibleRef{格前 4:5}，好使我们以纯洁的心眼透视罪恶的污秽黑暗，能揭露它们并将它们拖到光中；并能向那些已摆脱它们或仍被其束缚的人解释其缘由和本质，从而如先知所说，经过那可怕地燃烧我们心灵的罪恶之火，能立即无误地经过那熄灭它们的德行之水，并（仿佛）以属灵的疗法得以滋润，堪当以纯洁的心灵被引至完美的慰藉。\par

% structure: p0008 source=h2 target=subsection reason=relative-heading-rank; base=h2

\subsection{第三章}

\Emph{我们的第一场战斗必须对付贪饕的恶神，即口腹之乐。}\par

因此，我们必须首先进行对抗贪饕的战斗，即我们所说的口腹之乐；首先，既然我们要谈论禁食制度和食物的质量，我们必须再次回到埃及人的传统和习俗，因为众所周知，他们包含着更高级的节制训练和完美的分辨方法。\par

% structure: p0011 source=h2 target=subsection reason=relative-heading-rank; base=h2

\subsection{第四章}

\Emph{\Person{安当}院长的证言，教导每一种德行都应从特别擅长此德的人那里寻求。}\par

因为这是真福安当的一句古老而卓越的名言：一位隐修士在隐修生活计划中努力达到更高级成全的高峰，并学习了明辨的考虑，能够凭自己的判断站稳，并抵达独修生活的顶峰时，他绝不该向某一个人——即使他多么杰出——寻求所有种类的德行。因为一个人以知识之花为装饰，另一个人更坚强地以明辨的方法武装，另一个以忍耐的尊荣建立，另一个在谦逊的德行上卓越，另一个在节制上，另一个以纯朴的恩宠装饰。此人在宽宏大量上超越众人，彼人在怜悯上，另一人在守夜上，另一人在沉默上，另一人在工作的热忱上。因此，渴望采集属灵蜂蜜的隐修士，应当像最谨慎的蜜蜂，从那些特别拥有美德的人那里吸取德行，并殷勤地储存在自己胸中的器皿里：他不应调查任何人缺乏什么，而只应关注并收集他所拥有的任何德行。因为如果我们想从某一个人那里获得所有德行，我们将非常困难，或许永远找不到适合我们模仿的榜样。因为虽然我们尚未看到基督如宗徒所说成为\BibleQuote{万有中的万有}\BibleRef{格前 15:28}，但以这种方式我们可以在所有人中一点点地找到祂。因为论到祂说：\BibleQuote{祂由天主成了你们的智慧、正义、圣化和救赎}\BibleRef{格前 1:30}。因此，既然在一人身上找到智慧，在另一人身上找到正义，在另一人身上找到圣化，在另一人身上找到良善，在另一人身上找到贞洁，在另一人身上找到谦逊，在另一人身上找到忍耐，那么基督现今就被分肢分节地分散在众圣人之中。但当众人进入信德与德行的合一时，祂就形成\BibleQuote{圆满的人}\BibleRef{弗 4:13}，完成祂身体的圆满，在祂所有肢体的关节与特性中。直到那时，当天主成为\Quote{万有中的万有}时，目前天主可以以我们所说的方式，通过个别的德行\Quote{在万有内}，虽然尚未通过它们的圆满成为\Quote{万有中的万有}。因为虽然我们的宗教只有一个终极和目的，但有不同方式我们接近天主，正如在\WorkTitle{长老会议}中将更充分显示。因此，我们必须特别从那些我们看见借着圣神的恩宠更丰沛地流出这些德行的人那里寻求明辨和节制的榜样；并不是说任何人能独自获得那些分散在多人身上的事物，而是为了在我们所能达到的那些善性中，我们向那些特别获得它们的人前进。\par

% structure: p0014 source=h2 target=subsection reason=relative-heading-rank; base=h2

\subsection{第五章}

同一守斋的规则不能为所有人遵守。\par

因此，关于守斋的方式，不易遵守统一的规则，因为并非人人都有同样的力量；也不像其他德行那样仅靠心志的坚定就能获得。因此，由于它不仅取决于心灵的坚定，还涉及身体的可能性，我们得到了关于它的这一解释，这解释已传授给我们，即：饮食的时间、方式和质量应根据身体状况、年龄和性别的差异而有所不同；但就心灵的节制和精神的德行而言，对所有人都有同一的约束规则。因为不可能每个人都能延长他的斋戒一周，或通过两三天的不进食来推迟进餐。许多人因疾病、尤其是年老而衰弱，甚至无法忍受从日出到日落的禁食而不受苦。潮湿豆类的病弱饮食并不适合每个人；新鲜蔬菜的节俭饮食也不适合所有人；干面包的少量餐食并非人人都能同样接受。一个人对两磅食物不感到满足，另一个人一磅或六盎司的餐食就太多了；但所有人节制的同一目标和目的就是：没有人因贪饕而超过自己食欲的分量。因为不仅是食物的质量，而且数量也会钝化心智的敏锐，当灵魂和肉体都过饱时，就会点燃有害和炽热的恶习诱因。\par

% structure: p0017 source=h2 target=subsection reason=relative-heading-rank; base=h2

\subsection{第六章}

心智不仅被酒灌醉。\par

肚子充满各种食物时，就产生放荡的种子；心智被食物的重量窒息时，无法保持思想的引导和管治。因为不仅是酒的醉意惯于使心智昏醉，而且所有食物的过量也使它软弱不定，剥夺其纯洁清晰默观的一切能力。索多玛倾覆和放荡的原因并不是因酒而醉，而是食物的饱足。请听主通过先知责备耶路撒冷：\BibleQuote{因为你的妹妹索多玛怎样犯罪？岂不是她吃饱喝足？}\BibleRef{则 16:49}因为由于食物的饱足，他们被不可控制的肉体情欲所燃烧，被天主的审判用天上的火和硫磺焚毁。但如果单单食物的过量就通过饱食的恶习将他们推入如此坠落的罪中，那么对于那些身强力壮、竟敢以无节制的放纵享用肉和酒的人，我们该作何想？他们不仅取用身体软弱所需，而且取用心智的急切欲望所暗示的。\par

% structure: p0020 source=h2 target=subsection reason=relative-heading-rank; base=h2

\subsection{第七章}

身体的软弱如何不妨碍心灵的纯洁。\par

身體的軟弱並不是心靈純潔的障礙，只要所取的食物僅是身體軟弱所需，而非貪求享樂。較之那些在允許的必需品上適度取用的人，更容易找到完全戒除肥甘厚味的人；較之那些以身體軟弱為由取食卻仍能保持適中限度的人，更容易找到因熱愛節制而拒絕一切的人。因為身體的軟弱在自我節制上有其榮耀：雖然食物被允許給衰弱的身體，但人卻在需要時仍放棄滋補，只取用嚴格的明智判斷為生活所需所決定的適量食物，而非貪求的食慾所渴望的。越是精緻的食物，只要適度取用，既有利於身體健康，也無損於貞潔的純潔。因為由此獲得的力量，都消耗在辛勞和照顧的浪費之中。所以，正如任何生活狀態都不能被剝奪克己的德行，同樣，完美的榮冠也不會被拒絕給任何人。\par

% structure: p0023 source=h2 target=subsection reason=relative-heading-rank; base=h2

\subsection{第八章。}

\Emph{關於應以何種方式取用食物以達到完美節制的目標。}\par

因此，教父們有一句非常真實且卓越的言論：齋戒與克己的正確方法在於適中與身體克制的尺度；並且這是所有人完美德行的目標，即：雖然我們仍被迫渴望食物，但應當對所取用的食物實行自我節制，因為我們不得不為了維持身體的需要而進食。因為即使一個人身體軟弱，他也能達到完全的德行，並且與那些完全強壯健康的人相等，只要他以堅定的心神克制那些並非由於肉體軟弱而產生的慾望和貪情。因為宗徒說：\Quote{不要為肉體安排，去滿足私欲。}\BibleRef{羅 13:14}他並非完全禁止照顧肉體，而是說不要為肉體的慾望和貪情安排。他剪除了對肉體的奢侈喜愛，並未排除生活所必需的管束：他這樣做，是為避免因縱容肉體而陷入慾望的危險糾纏；後者則是為防止因我們的過錯而使身體受損，無法履行屬靈和必要的職責。\par

% structure: p0026 source=h2 target=subsection reason=relative-heading-rank; base=h2

\subsection{第九章。}

\Emph{關於應採取的克苦尺度以及齋戒的補救。}\par

因此，克己的完美並不單憑時間的長短來衡量，也不僅憑食物的品質，而首先取決於良心的判斷。因為每個人都應根據自己身體鬥爭的需要，為自己規定如此簡樸的飲食。教會規定的齋戒確實有價值，並且必須遵守。但是，如果齋戒之後不是有節制的進食，人就不能達到完美的目標。因為長期齋戒的克己——隨後卻是飽腹——帶來的是一時的疲倦，而非純潔與貞潔。心智的完美確實取決於肚腹的節制。凡不滿足於始終維持均衡適度的飲食者，便沒有持久的純潔與貞潔。齋戒雖然嚴厲，但若隨後有不必要的鬆懈，就會變得無用，並立即導致貪饕的惡習。每日以適度合理的進食，勝過偶爾的嚴厲長期齋戒。過度的齋戒不僅會削弱心智的堅定，還會因身體的極度疲倦而削弱祈禱的力量。\par

% structure: p0029 source=h2 target=subsection reason=relative-heading-rank; base=h2

\subsection{第十章。}

\Emph{僅憑禁食不足以保持身體和心靈的純潔。}\par

為了保持心智和身體於完美狀態，僅憑禁食是不夠的：除非同時結合心智的其他德行。因此，必須先通過服從的德行學習謙卑，以及艱苦的勞作和身體的疲憊。不僅必須避免擁有錢財，而且必須徹底根除對錢財的慾望。因為僅僅不擁有錢財——這對許多人來說是出於無奈——還不夠：而且我們若偶然得到錢財，甚至不應有擁有的\Emph{意願}。憤怒的狂亂應受控制；沮喪的低落表情應被克服；虛榮應被藐視；驕傲的傲慢應被踐踏；心智的飄忽遊蕩的思想應藉對天主的持續紀念而受到約束。我們心中的遊移不定應被帶回對天主的默觀，每當奸詐的敵人試圖將心智從這默觀中擄走，潛入內心深處時。\par

% structure: p0032 source=h2 target=subsection reason=relative-heading-rank; base=h2

\subsection{第十一章。}

\Emph{身體的私慾只有通過徹底根除惡習才能消除。}\par

因为身体炽烈的冲动除非先将其他主要恶习的刺激彻底根除，否则绝不可能被熄灭；关于这些恶习，若天主允许，我们将在适当的地方，分别在各卷书中详论。但现在我们必须处理贪饕，即口腹之欲，这是我们首先要与之战斗的。因此，凡不能节制食欲的人，绝不可能克制炽烈情欲的冲动。内在之人的贞洁是由这一德行的完美所显出的。因为，若你见一个人能在更高层次的战斗中战胜更强的敌人，却曾被较弱的敌人击败，你绝不会相信他能与更强的对手抗衡。因为所有德行的本性都是同一的，尽管它们似乎被分为许多不同的种类和名称；正如金子虽因金匠的技艺而看似分布于许多不同种类的首饰中，但其本质仍是同一的。因此，一个人若在某一部分德行上有所缺失，便证明他并未完全拥有任何德行。因为我们怎能相信，一个人若不能平息仅由内心无节制而爆发的愤怒之刺，却能熄灭那不仅由身体刺激、也由心灵之恶所点燃的贪欲之火呢？或者，我们怎能认为，一个人若不能克服骄傲这一单纯的过失，却能压制肉身与精神的放纵欲望呢？或者，我们怎能相信，一个人若不能弃绝那属于外在且在我们本性之外的金钱之爱，却能践踏那根植于肉身的放纵之情呢？一个人若连沮丧之病都无力治愈，又怎能战胜肉性与精神的争战呢？无论一座城市如何被高墙和紧闭城门的坚固所保护，但只要有一个甚小的后门被放弃，它仍会沦陷。因为危险的敌人是从高墙和宽阔的城门进入城中心，还是从隐秘狭窄的通道进入，又有何区别呢？\par

% structure: p0035 source=h2 target=subsection reason=relative-heading-rank; base=h2

\subsection{第十二章}

\Emph{在我们的属灵争战中，当从属肉体的争战汲取榜样。}\par

\Quote{凡在竞赛中角逐的，若非按规则较量，就不能得到冠冕。} \BibleRef{弟后 2:5} 凡想熄灭肉体自然欲望的人，应首先赶紧克服那些其位置在我们本性之外的恶习。因为若我们想体验宗徒这话的力量，我们应先了解世俗竞赛的规矩和纪律，以便最终通过与这些的比较，我们能够知道蒙福的宗徒借此比喻愿意教导我们这些在属灵争战中角逐的人什么。因为在那些争战中，正如同一宗徒所说，给胜利者提供了\Quote{可朽坏的冠冕} \BibleRef{格前 9:25}，所遵守的规则是：凡立志为那以豁免特权为装饰的荣耀冠冕而准备自己，并渴望在竞赛中进入最高层次角逐的人，应先在奥林匹克和皮提亚赛会中作为青年展示其能力及其初始的力量；因为在这些赛会中，那些想练习这种训练的青年，要通过赛会主席和全体群众的评判，来测试他们是否配得或应当被接纳。当一个人经过仔细测试，首先证明其生平无任何可耻污点，其次被判定并非因奴隶身份而不配参加这种训练和练习者团体，第三，他充分展示了自己的能力和英勇，通过与青年及同龄人的较量，显出了其作为青年的技巧和勇气，并从儿童比赛逐步晋升，经主席审查获准与成年人和有经验的人混合，且不仅通过与他们不断较量显示出同等的英勇，还多次在他们中间荣获胜利奖，最终他才被允许接近竞赛中最辉煌的角逐。只有胜利者及那些佩带许多冠冕和奖品的人，才被准许参加这种角逐。若我们理解了这来自肉体竞赛的比喻，就应通过与它的比较，也了解我们属灵争战的系统和方式。\par

% structure: p0038 source=h2 target=subsection reason=relative-heading-rank; base=h2

\subsection{第十三章}

\Emph{除非我们已摆脱贪饕之恶，否则不能进入内在之人的争战。}\par

我们也应首先证明自己免受肉体的奴役。因为\Quote{一个人被谁制胜，就是谁的奴隶。} \BibleRef{伯后 2:19} 并且\Quote{凡犯罪的，就是罪的奴隶。} \BibleRef{若 8:34} 当竞赛的主席审查发现我们未沾染可耻情欲的污点，并判定我们不是肉体的奴隶，既不卑贱也不配参加对抗恶习的奥林匹克争战时，我们便能进入与我们的同等者——即肉体的情欲和灵魂的冲动与不安——的较量。因为满腹之人不可能尝试内在之人的争战；而一个在轻微冲突中就能被击败的人，也不配在更艰难的争战中受考验。\par

% structure: p0041 source=h2 target=subsection reason=relative-heading-rank; base=h2

\subsection{第十四章}

\Emph{如何克服贪饕之欲。}\par

首先，我们必须践踏贪食之欲，为此，不仅当以禁食节制心神，更当藉守夜、诵读及频频为心中所知所欺所败之事而痛悔，时而因罪而恐惧叹息，时而因渴慕成全与圣德而炽热，直至心神全然被此等忧思与默想所占据，且视进食非为纵欲之许可，乃为加诸己身之重担，更多是肉体之需，而非灵魂之所欲。藉此心志之热忱与不断之痛悔，我们必能击溃肉体的放荡（它因食物滋养而愈显骄傲自大）及其危险的煽动，并以丰沛之泪与心中之哀哭，扑灭我们身体中那被巴比伦王点燃的火炉——他不断供给我们犯罪之机，以及如同石脑油与沥青般令我们燃烧更烈的恶习——直至因天主的恩宠，藉祂的圣神如露滴入我们心中，肉欲之火得以全然熄灭。这便是我们的首战，这仿佛是我们奥林匹克竞技中的初试：以对成全的渴望扑灭口腹之欲。为此，我们不独当藉修德之默想践踏一切不必要的食欲，即使取用维系本性所需之物，亦不可无心中之忧惧，视之如与贞洁相悖。如此，我们方能步入生命的征程，以至于除因肉体软弱而不得不降身为照料身体之外，再无时刻感到被召离属灵的学习。当我们受此必需所迫——照料生命所需而非灵魂之欲——时，应速速退出，如同它拦阻我们于真正有益的操练。因我们不可能轻视摆在眼前的食物之乐，除非心神专注于默思天主之事，且更沉醉于对德行之爱及天上事物之喜乐。如此，人将视一切现世之物为转瞬即逝，当他稳固地将心灵之眼注视于那些不变而永恒的事物，并虽在肉身中，已用心神默观未来生命的福乐时。\par

% structure: p0044 source=h2 target=subsection reason=relative-heading-rank; base=h2

\subsection{第十五章}

\Emph{隐修士须时常切望保持心灵纯洁。}\par

这好比有人致力击中高悬的某种伟大德行之奖，仅凭一极小标识；他以最敏锐的目光瞄准飞镖，深知荣耀与奖赏之丰厚在于命中；他转目不视其他，必须定睛于此，视赏报所在；因若目光稍有偏失，必丧失技艺之奖与勇武之赏。\par

% structure: p0047 source=h2 target=subsection reason=relative-heading-rank; base=h2

\subsection{第十六章}

\Emph{隐修士应仿效奥林匹克竞技，唯有战胜肉体，方可企望属灵之战。}\par

因此，当口腹之欲藉这些思虑被克服，我们如同在奥林匹克竞赛中被宣告既非肉体之奴，亦非因罪而蒙羞，则将被判定配得上更高层次的争斗，并离开此类课程，相信我们有资格进入与属灵邪恶的竞技，唯有得胜者及获准参与属灵争战者，才被视为堪当此战。因为可以说，这是所有争斗最坚固的基础：首先摧毁肉体欲望的冲动。因为除非先战胜自己的肉体，无人能合法地争斗。不合法争斗者，当然不能参与竞赛，也不能赢得荣耀之冠与胜利之恩。但若在此战中被击败，被证实为如同肉欲之奴，因而显不出自由与力量的印记，我们将立即被从与属灵军旅的争斗中驱逐，如同不配与奴仆，带着种种羞愧。因为\Quote{凡犯罪的就是罪的奴仆。}（\BibleRef{若 8:34}）蒙福的宗徒将这话连同那受指责有淫乱的人一同对我们说：\Quote{你们所遇见的试探，无非是人所能受的。}（\BibleRef{格前 10:13}）因为我们若不寻求心志的力量，便不配经受更严厉的与高天邪恶之争，若我们连那反抗圣神的软弱肉体都无法制服。有些人误解宗徒这话，将直陈语气读作虚拟语气，即：\Quote{愿你们所遇见的试探，无非是人所能受的。}但显然，他是以宣告或责备之意，而非愿望之意说的。\par

% structure: p0050 source=h2 target=subsection reason=relative-heading-rank; base=h2

\subsection{第十七章}

\Emph{属灵争战的基础与根基应建立在与贪饕的斗争上。}\par

你要听一位真正的基督运动员按照争斗的规则和法律奋力拼搏吗？他说：\BibleQuote{我这样跑，不是如同无定向的；我这样打拳，不是如同打空气的；而是我惩戒我的身体，使它服从，免得我给别人传报，自己反而被弃绝。}\BibleRef{格前 9:26}你看他如何把争斗的主要部分放在自己身上，即放在他的肉体上，如同放在最稳固的基础上，并把战斗的结果完全置于对肉体的惩戒和身体的服从之中。\BibleQuote{我这样跑，不是如同无定向的。}他跑得并非无定向，因为仰望天上的耶路撒冷，他有一个标竿，他的心毫不偏离地迅速向它直去。他跑得并非无定向，因为\BibleQuote{忘记背后，努力面前的，向着标竿直跑，要得天主在基督耶稣内从高天召叫的奖赏，}\BibleRef{斐 3:13}他常把心灵的视线投向那里，并以心灵的全部速度向它奔去，满怀信心地宣告：\BibleQuote{这场好仗我已打完，这场赛跑我已跑到终点，这信仰我已保持了。}\BibleRef{弟后 4:7}因为他知道自己以心中热诚的奉献，\BibleQuote{追逐香膏的芬芳}\BibleRef{歌 1:3}不倦地奔跑，并通过对肉体的惩戒赢得了属灵战斗的胜利，他大胆地得出结论说：\BibleQuote{从此以后，正义的冠冕已为我预备下了，就是主，正义的审判者，在那一天要赏给我的。}而且，为了也给我们开启同样报酬的希望，如果我们愿意效法他赛跑的奋力，他补充说：\BibleQuote{不但赏给我，而且也赏给一切爱慕祂显现的人；}\BibleRef{弟后 4:8}宣告在审判之日，我们将分享他的冠冕，如果我们爱慕基督的来临——不仅指那一次连反对者也不得不显现的来临，也指那每日在圣灵魂中发生的来临——并且通过惩戒身体在战斗中获胜。主在福音中论及这来临说：\BibleQuote{我和我父要到祂那里去，并要在祂那里作我们的住所。}\BibleRef{若 14:23}又说：\BibleQuote{看，我立在门口敲门：谁若听见我的声音而给我开门，我要进到他那里去，同他坐席，他也要同我坐席。}\BibleRef{默 3:20}\par

% structure: p0053 source=h2 target=subsection reason=relative-heading-rank; base=h2

\subsection{第十八章}

\Emph{论荣福宗徒藉以升至至高战斗荣冠的各种冲突与胜利的不同数目。}\par

但他说的\BibleQuote{我这样跑，不是如同无定向的}（这句话更特别地关系到他的心意和精神的热情，他借此以全副热诚跟随基督，同新娘一起呼喊：\BibleQuote{我们要因你香膏的芬芳而奔跑；}\BibleRef{歌 1:3}又说：\BibleQuote{我的灵魂紧贴着你；}）但他也见证自己在另一种争斗中得了胜，说：\BibleQuote{我这样打拳，不是如同打空气的，而是我惩戒我的身体，使它服从。}这正关系到戒律的劳苦、身体的斋戒和肉体的折磨：他以此表明自己是自己肉体的有力打击者，并指出他针对肉体所施的节制打击并非徒然；而是通过制服自己的身体赢得了战斗的胜利；因为当肉体以节制的打击受到惩戒、以斋戒的拳击被击倒时，他为自己的得胜精神取得了不朽的荣冠和不朽坏的奖赏。你见到这争斗的正统方法，并默想属灵战斗的结果：基督的运动员如何战胜了反叛的肉体，几乎将它踩在脚下，高高地凯旋前进。因此，他\BibleQuote{不是如同无定向地跑，}因为他相信自己即将进入圣城、天上的耶路撒冷。他\BibleQuote{这样打拳，}即借着斋戒和肉身的屈辱，\BibleQuote{不是如同打空气的，}即不是借着节制的打击在虚空中挥拳，通过这些打击，他不是打击了空虚的空气，而是借着身体的惩戒打击了那些居住在空气中的邪灵。因为说\BibleQuote{不是如同打空气的}的人，表明他所打击的——不是空虚和空洞的空气，而是空气中的某些存在。因为他已在这种争斗中得胜，满载着许多冠冕的奖赏行进，他便开始与更强大的敌人较量，这是理所当然的；在战胜了先前的对手之后，他大胆地宣告说：\BibleQuote{现在我们战斗不是对抗血和肉，而是对抗率领者，对抗掌权者，对抗这黑暗世界的霸主，对抗天界里邪恶的鬼神。}\BibleRef{弗 6:12}\par

% structure: p0056 source=h2 target=subsection reason=relative-heading-rank; base=h2

\subsection{第十九章}

\Emph{论基督的运动员只要在肉身内，就绝不能没有战斗。}\par

基督的运动员，当他仍在肉身之时，从不缺少在竞赛中可争取的胜利；但他愈是因胜利的成就而成长，就愈有更严酷的斗争面对他。因为当肉身被降服并征服时，他的胜利会激起何等成群的仇敌、何等众多的敌人起来反对这位得胜的基督士兵！惟恐在和平的安逸中，基督的士兵会松懈努力，开始忘记他竞赛的光荣斗争，并因免于危险而产生的怠惰而变得松弛，从而被骗取奖赏和胜利的报酬。因此，如果我们想以不断增长的德行上升到这些胜利的阶段，我们也应当以同样的方式进入战斗的竞技场，开始同宗徒一起说：\Quote{我这样打拳，不是如同打空气的；我痛击我身，使它为奴。} \BibleRef{格前 9:26} 好使这场冲突结束后，我们能再次同他一起说：\Quote{我们战斗不是对抗血和肉，而是对抗率领者，对抗掌权者，对抗这黑暗世界的霸主，对抗天界里邪恶的鬼神。} \BibleRef{弗 6:12} 因为如果我们败在属肉的争斗中，被肚腹的斗争击倒，我们就不可能与他们交战，也无法经历属灵的战斗；宗徒将用责备的话对我们说：\Quote{你们所受的试探，无非是普通人所能受的试探。} \BibleRef{格前 10:13}\par

% structure: p0059 source=h2 target=subsection reason=relative-heading-rank; base=h2

\subsection{第二十章。}

\Emph{隐修士若想进行内心冲突的斗争，应当如何不逾越适当的进食时辰。}\par

因此，隐修士若想进行内心冲突的斗争，应当首先为自己定下这项预防措施：即他在任何情况下都不允许自己被任何美味所胜，在斋戒结束和适当的用膳时辰到来之前，不可在用餐时间之外进食或饮水；并且，用餐结束后，绝不允许自己再吃一口食物，无论多么小；同样，他也应遵守睡眠的常规时间和长度。因为必须像根除不洁之罪一样，戒除放纵。若一个人不能克制食欲上不必要的欲望，他怎能熄灭肉欲之火呢？若一个人不能控制那些明显且微小的情欲，他怎能以节制的明智来对抗那些隐秘的、在无人看见时激动他的情欲呢？因此，心灵的力量是在分别的冲动和任何情欲中受到考验的：如果它在非常微小和明显的欲望上被打败，那么它在真正巨大而有力且隐秘的欲望上又将如何承受，这是每个人的良心必须为自己作证的。\par

% structure: p0062 source=h2 target=subsection reason=relative-heading-rank; base=h2

\subsection{第二十一章。}

\Emph{论隐修士的内在平安和属灵的节制。}\par

因为我们所惧怕的不是外在的敌人。我们的仇敌被关在我们自己里面：一场内在的战争每天由我们进行；如果我们在这场战争中得胜，一切外在的事物都将变得软弱，一切对于基督的士兵都将变得和平且顺服。如果我们内在的自我被克服并顺服于灵，我们将没有外在的敌人可惧。我们不要相信仅仅从可见的食物中外在的禁食，就能足以达到心灵的成全和身体的纯洁，除非与此同时也结合了灵魂的禁食。因为灵魂也有它的食物，是有害的，靠这些食物滋养，即使没有肉食的过剩，它也会陷于放荡的堕落。毁谤是它的食物，而且是一种非常可口的食物。一阵怒气也是它的食物，即使是非常轻微的；然而在同一时刻，它用致命的滋味摧毁它。嫉妒是心灵的食物，用其毒汁腐蚀它，并永不停止地使他人的幸福和成功使它变得悲惨和可怜。\OriginalTerm{虚妄光荣}{Kenodoxia}是它的食物，一时以美味的餐食取悦它；但之后剥去它所有美德，使它赤裸且空虚，并遣散它，使它荒芜且缺乏一切属灵的果实，以致它不仅失掉巨大劳苦的奖赏，而且还招致更重的惩罚。一切贪欲和心灵的飘忽不定都是灵魂的一种食物，以有害的肉滋养它，但之后却不给它分享天上的食粮和真正坚实的食物。因此，如果我们尽一切力量以最神圣的禁食禁绝这些，我们身体的禁食实践就会既有用又有益。因为肉体的劳作，若与心灵的痛悔结合，将产生最中悦天主的祭献，并在纯洁无瑕的心灵内室中成为圣洁的圣所。但是，如果我们在身体方面禁食，却在灵魂上被最危险的恶习所缠绕，那么我们肉身的谦卑将毫无益处，同时我们最宝贵的部分被玷污了；因为我们藉以成为圣神居所的那个实体出了差错。因为成为天主的居所和圣神的宫殿的，并非可朽的肉体，而是洁净的心。因此，每当外在的人禁食时，我们也应当约束内在的人不摄入对他有害的食物；即那位内在的人，蒙福的宗徒尤其敦促我们将他纯洁地呈于天主面前，使他堪当在内里接纳基督为宾客，说：\Quote{使基督因着信德，住在你们心里。} \BibleRef{弗 3:16}\par

% structure: p0065 source=h2 target=subsection reason=relative-heading-rank; base=h2

\subsection{第二十二章。}

\Emph{我们应当为此理由操练身体的斋戒，以便借此达到属灵的斋戒。}\par

因此我们晓得，我们应当留意身体的节制，好能藉此禁食达到心灵的纯洁。否则，我们的劳苦必归于徒然，若我们毫无倦怠地忍耐此苦，只顾注视目标，却无法达到我们为忍受这些试炼所追求的目标；那么，与其从身体上禁食那些无关紧要且无害之物，倒不如禁绝灵魂所禁止的食物；因为对于后者，只是单纯而无害地领受天主的受造物，其本身并无过错；但就前者而言，一开始就有吞食弟兄的危险倾向；对此经上说：\Quote{不要喜爱毁谤，免得你被拔除。}\BibleRef{箴 20:13} 关于愤怒和嫉妒，可敬的约伯说：\Quote{愤怒杀死愚人，嫉妒杀死孩童。}\BibleRef{约 5:2} 同时应注意，愤怒的人被定为愚人，嫉妒的人被定为孩童。前者被称为愚人并非没有理由，因为他出于自愿，被愤怒的刺所驱使，自取灭亡；后者在嫉妒时，证明自己是孩童和幼稚者，因为当他嫉妒别人时，他显示他所恼怒其繁荣的那个人，比他更大。\par

% structure: p0068 source=h2 target=subsection reason=relative-heading-rank; base=h2

\subsection{第二十三章。}

\Emph{隐修士的饮食应有何种特征。}\par

我们应当选择这样的食物：不仅不能助长情欲的灼热，还要避免点燃它；而且应当易于准备、价格低廉，适合弟兄们的生活和共同使用。因为贪饕的本性有三重：第一，强迫我们提前吃饭；第二，喜欢填满肚腹，贪食各种食物；第三，喜好精细和美味的宴席。因此，隐修士应当三重防范：第一，等待适宜的时辰开斋；第二，不贪食无度；第三，满足于任何普通的食物。因为任何超过常规和众人共同使用之外的东西，都被教父们的古老传统视为沾染了虚荣、夸耀和显摆之罪。在我们所见因学识和明辨而名副其实杰出的人中，或基督的恩宠分别出来作为众人效法的明灯者，我们未曾见过有人不吃在他们那里算是便宜且容易获得的饼；我们也未曾见有人拒绝此规则而放弃吃饼，改为以豆类、蔬菜或水果为食，就被列入最受尊敬者之中，甚至获得知识和明辨的恩宠。因为他们不仅规定隐修士不应要求他人不常吃的食物，以免他的生活方式公开暴露给众人而变得虚空、无用，并因虚荣病而毁坏；而且坚持禁食的共同惩戒操练不应轻易向任何人透露，而应尽可能隐藏和保密。但当有弟兄到来时，他们规定我们应当表现出慈善和爱德的美德，而不是遵守严格的禁食和日常规则；我们不应考虑自己的意愿和利益或欲望的热切，而应摆在面前，乐于满足客人需要的任何饮食或他的软弱所要求我们的东西。\par

% structure: p0071 source=h2 target=subsection reason=relative-heading-rank; base=h2

\subsection{第二十四章。}

\Emph{我们如何在埃及见到，在我们到达时，每日禁食毫无顾忌地被打破了。}\par

当我们从叙利亚地区来到埃及省，渴望学习长老们的规则时，我们惊奇于那里接待我们时的热忱之心，以至于没有遵守我们在巴勒斯坦隐修院里所养成的、在禁食规定时辰结束前不得进食的规则；除了在例行日子（星期三和星期五）之外，我们每到一处，每日禁食都被打破。当我们问为什么他们毫无顾忌地忽略每日禁食时，一位长者回答说：\Quote{禁食的机会常在我手中。但我将送你们上路，不能总留着你们。禁食虽有益且明智，却是自由奉献。但命令的急迫要求完成一件爱德工作。因此，我接待在你们内的基督，应当款待祂；等你们离开后，我就能以更严格的禁食来平衡为祂所行的款待。因为\InnerQuote{新郎的伴郎们不能在新郎与他们同在时禁食；}但祂离去后，他们就要合宜地禁食。}\par

% structure: p0074 source=h2 target=subsection reason=relative-heading-rank; base=h2

\subsection{第二十五章。}

\Emph{论一位长者如此俭省地进食六次，以致仍然饥饿。}\par

当一位长者在我进食时催我再多吃一点，我说不能时，他回答说：\Quote{我已经为不同的弟兄们六次铺好了餐桌，每次我都催促他们，并与他们一同进食，如今仍感饥饿；而你，现在第一次进食，却说再也吃不下了？}\par

% structure: p0077 source=h2 target=subsection reason=relative-heading-rank; base=h2

\subsection{第二十六章。}

\Emph{论另一位长者，他从不在他的小室中独自进食。}\par

我们曾见过另一位独居者，他声称自己从未独自享用食物；即使连续五天没有弟兄来到他的小室，他也总是推迟进食，直到周六或周日他去教堂参加礼仪时，找到某位陌生人，立刻将他带回自己的小室，和他一起分享身体的饮食——这样做与其说是出于自身需要，不如说是出于仁慈和为了弟兄的缘故。因此，既然他们知道日常斋戒因弟兄到来而被打破而毫无疑虑，当弟兄离开时，他们便以更大的克苦来补偿因弟兄而享受的饮食，并严厉地以更严格的惩罚来弥补接受哪怕是极少量食物的过失——不仅在面包方面，还减少通常的睡眠。\par

% structure: p0080 source=h2 target=subsection reason=relative-heading-rank; base=h2

\subsection{第二十七章}

\Emph{两位院长帕埃修斯和若望论自己热忱的果实。}\par

当年长的\Person{若望}（他管理一所大型隐修院和众多弟兄）来探望年长的\Person{帕埃修斯}（他住在广漠之中），并作为老朋友被问起，在这与他分离、几乎从未被弟兄打扰的四十年中他做了什么：\Quote{从未，}他说，\Quote{太阳从未看见我进食，}\Quote{也未看见我发怒，}另一位说道。\par

% structure: p0083 source=h2 target=subsection reason=relative-heading-rank; base=h2

\subsection{第二十八章}

\Emph{院长若望临终留给门徒的教训和榜样。}\par

当同一位长者，如同即将归去的人，奄奄一息时，弟兄们围立四周，恳求他留下一些特殊的嘱咐作为遗产，使他们能由此轻易达到完美的顶峰，因嘱咐简短而更易遵守。他叹息着说道：\Quote{我从未做过自己的意愿，也未曾教导任何人我自己未曾首先实行的事。}\par

% structure: p0086 source=h2 target=subsection reason=relative-heading-rank; base=h2

\subsection{第二十九章}

\Emph{论院长玛赫特斯，他在属灵谈话中从不入睡，却在世俗故事中总是睡着。}\par

我们认识一位年长者，名叫\Person{玛赫特斯}，他远离弟兄们的喧嚣居住，并通过每日祈祷从主那里获得这恩宠：每逢举行属灵谈话，无论是白天还是夜晚，他从未被睡眠压倒；但如果有人试图引入一句诋毁或闲谈，他立刻入睡，仿佛毁谤的毒药无法渗透污染他的耳朵。\par

% structure: p0089 source=h2 target=subsection reason=relative-heading-rank; base=h2

\subsection{第三十章}

\Emph{同一位长者关于不可判断任何人的言论。}\par

同一位长者教导我们，不应该判断任何人。他指出有三点他曾责备和警告弟兄们：即有些人允许割去自己的悬雍垂，或在室内保留一件斗篷，或祝福油并给予世上的请求者。他说他自己也做过这一切。因为患上悬雍垂的疾病，他由于它的弱点而消瘦了很久，最后被疼痛和所有长者的劝说所迫，允许割掉它。也因为这场病，他被迫保留一件斗篷。他还被迫祝福油并给予那些祈求的人——这是一件他极其憎恶的事，因为他认为这出自内心的狂妄——突然，许多在世上生活的人围住他，他无法以其他方式逃脱，除非他们以不小的暴力和恳求强求他伸手放在他们献上的器皿上，并划十字圣号；于是他们相信得到了祝福的油，终于放他走了。通过这些事，他清楚地发现，一个隐修士正陷入相同的处境，纠缠于他曾冒险判断别人的同样过失。因此，每个人只应判断自己，并谨慎留意所有事，不判断他人的生活和行为，符合宗徒的告诫：\BibleQuote{但是，你为什么判断你的弟兄？他或站住或跌倒，都是他主人的事。}以及这一点：\BibleQuote{你们不要判断，免得你们受判断；因为你们用什么判断来判断，你们也要受什么判断。}因为除了我们已说的理由，判断别人也是危险的，因为在那些我们被冒犯的事情上——我们不知道他们真正在行为上是正确于天主面前，还是至少情有可原——我们发现我们冒失地判断了他们，并因此犯了不小的罪，因为我们形成了与应当的对弟兄的看法不同的意见。\par

% structure: p0092 source=h2 target=subsection reason=relative-heading-rank; base=h2

\subsection{第三十一章}

\Emph{同一位长者看见弟兄们在属灵谈话中入睡，而在闲谈故事中醒来时的斥责。}\par

同一位长者通过这个证据清楚地表明，是魔鬼鼓励闲谈，并总是显明自己为属灵谈话的敌人。因为当他向一些弟兄谈论必要的事和属灵之事时，看见他们被沉睡压倒，无法驱除眼中的睡意，他突然引入一则闲谈故事。当他看见他们立刻醒来，喜形于色，竖耳倾听，他叹息着说道：\Quote{到现在为止，我们一直谈论天上的事，而你们的眼睛都被沉睡压倒了；但一引入闲谈的故事，我们都醒过来，摆脱了压服我们的睡意。因此，从这可以分辨谁是那属灵谈话的敌人，谁显明自己是那无用而属血肉的谈话的建议者。因为最明显地表明，正是那以恶为乐者，从不停止鼓励后者而反对前者。}\par

% structure: p0095 source=h2 target=subsection reason=relative-heading-rank; base=h2

\subsection{第三十二章}

\Emph{论未读而被焚毁的书信。}\par

我也不认为讲述一位弟兄的行为是次要之事，他追求心灵的纯洁，且极其渴望默观神圣之事。当过了十五年，许多信件从他在\Place{本都}省的父亲、母亲及许多朋友那里送到他手中时，他接过那一大捆信，心中思忖片刻，说道：\Quote{读这些信会给我带来什么思想呢？它们要么激起无谓的快乐，要么引起无用的悲伤！这些信会借着对写信者的回忆，使我的心离开我所定意的默观多少日子呢？由此产生的纷乱心境需要多少时间才能平息，而如果心灵一旦被信件的温情所动，又因回忆起那些被他长久离开之人的言语和面容，开始在思想和精神上重新探望他们、留驻其中、与他们同在，那么要恢复先前那种平安状态又需要付出多大的努力呢？况且，如果在心里开始注视他们，并重新接纳和唤起那在舍弃世俗时每个人都像死去一般放弃的记忆，那么身体上的离开也是徒劳无益的。}他心中盘算着，决定不但不读任何一封信，甚至不打开那捆信，唯恐一看到写信者的名字或想起他们的面貌，内心的意志就会动摇。于是，他将那捆信原封不动地扔进火里烧掉，喊道：\Quote{去吧，我家中的思念，烧掉吧，再也别试图把我召回那些我已逃离的事物。}\par

% structure: p0098 source=h2 target=subsection reason=relative-heading-rank; base=h2

\subsection{第三十三章}

\Emph{论阿博特·\Person{德奥多若}通过祈祷获得对问题的解答。}\par

我们也认识阿博特·\Person{德奥多若}，他是一位极具圣德和全备知识的人，不仅在实践生活上，而且在理解圣经上也是如此；他的知识并非主要靠学习和阅读或世俗教育获得，而单凭心灵的纯洁：因为他只能勉强听懂和说出很少几个希腊词。此人曾寻求一个极难问题的解释，便连续七天七夜不停祈祷，直到借着主的启示找到了所提问题的解答。\par

% structure: p0101 source=h2 target=subsection reason=relative-heading-rank; base=h2

\subsection{第三十四章}

\Emph{论同一位长者的话，藉此教导隐修士能通过何种努力获得圣经的知识。}\par

因此，当一些弟兄惊叹于他知识的光辉，向他请教圣经的某些含义时，他说，一位想要获得圣经知识的隐修士，不应将劳力花费在注释者的著作上，而应将自己心灵的一切努力和内心的意向都专注于从肉体的恶习中洁净自己；因为一旦驱除了这些恶习，心灵的眼睛就如揭开了情欲的幔子，立刻开始仿佛自然地凝视圣经的奥秘；因为圣经并非被圣神的恩宠赐予我们，好让它们保持未知和隐晦；相反，它们是由于我们的过错而变得隐晦，因为我们罪恶的幔子覆盖了心灵的眼睛；当心灵恢复其自然健康状态时，仅仅阅读圣经本身便足以看见真知，并不需要注释者的帮助，正如这些肉身的眼睛，只要没有昏花或黑暗的遮蔽，就不需要人教它如何看。正因为此，注释者之间才产生了如此大的分歧和错误，因为大多数人不注意净化心灵，就贸然投入解释圣经的工作，并按照他们心中稠密或污浊的程度，形成彼此冲突且违背信仰的看法，因此无法接受真理的光照。\par

% structure: p0104 source=h2 target=subsection reason=relative-heading-rank; base=h2

\subsection{第三十五章}

\Emph{同一位长者在半夜来到我的小室时的责备。}\par

同一位 \Person{德奥多若}在深夜出乎意料地来到我的小室，以慈父般的好奇心想看看我这个尚未成形的独修隐士独自在做些什么；当他发现我已做完晚祷，正躺在一张席子上开始休息疲倦的身体时，他发自内心地叹息，叫着我的名字说：\Quote{ \Person{若望}啊，此时此刻有多少人在与天主交谈，拥抱祂，并将祂留在身边，而你却因懒惰的睡眠而软弱，被剥夺了如此伟大的光明！}\par

既然诸圣父的德行及赐予他们的恩宠引诱我们转向这样一个故事，我认为在这卷书中记录一件出自那位最优秀之人\Person{阿尔赫比乌斯}善意的值得注意的爱德事迹是好的，以便嫁接在爱德工作上的节制的纯洁性，因着愉悦的变化而更加闪耀。因为禁食的义务只有在以爱德的果实使之完美时，才被天主所悦纳。\par

% structure: p0108 source=h2 target=subsection reason=relative-heading-rank; base=h2

\subsection{第三十六章}

\Emph{关于\Place{迪奥尔科斯}荒漠的描写，独修隐士们居住在那里。}\par

于是，当我们尚为初学者时，从巴勒斯坦的隐修院来到一座名为\Place{迪奥尔科斯}的埃及城市，看到众多隐修士受\OriginalTerm{集体修院}{Cœnobium}的规诫约束，并在那最早且卓越的隐修院体制中受训，我们也以全心智慧渴望看到另一种更优越的制度，即\OriginalTerm{独修隐修士}{anchorites}的制度，因众人对它的赞美激励了我们。因为这些人先在集体修院中生活了很长时间，勤勉学会了忍耐与明辨的一切规则，获得了谦逊与弃绝的德行，并完全克服了所有缺点，为了与魔鬼进行最可怕的战斗，而深入沙漠最幽深之处。之后我们发现这类人居住在\Place{尼罗河}附近，一个一边被该河环绕、另一边被大海包围的地方，形成了一个岛屿，只有寻求这种隐僻之地的隐修士才能居住，因为土壤的盐碱和沙地的干旱使该地不适于任何耕种——这些隐修士，我们急切赶去，对他们因默观德行与热爱独处而忍受的劳作感到无比惊异。因为他们甚至被水的极度匮乏所困扰，以至于他们分配水的细心和精确胜过任何吝啬鬼保存和囤积最珍贵酒类的方式。因为他们从上述河床运水三里甚至更远，以供一切必需；而这距离本已很大，还隔着沙丘，又因这项任务的极大困难而加倍。\par

% structure: p0111 source=h2 target=subsection reason=relative-heading-rank; base=h2

\subsection{第三十七章。}

\Emph{论院长\Person{阿尔赫比乌斯}将其连同家具的斗室让给我们。}\par

于是我们见到此景，心中燃起效法他们的渴望，前述的\Person{阿尔赫比乌斯}，因慈善之恩在他们中最负盛名，将我们拉进他的斗室，发现我们的渴望后，假装想要离开此地，并将他的斗室让给我们，好像他要离去，宣称即使我们没有来他也会这样做。我们因渴望留在那里，对如此伟人的断言深信不疑，欣然同意，并接管了他的斗室及其所有家具和物品。于是他实现了他的虔诚欺诈，离开几日以筹措建造斗室的材料，之后返回，以最大劳力为自己另建了一座斗室。不久之后，当其他一些弟兄怀着同样的渴望来此居住时，他又用类似的慈善谎言欺骗了他们，并将这座斗室及其一切相关物品让给了他们。但他不倦地坚持他的慈善之举，为自己建造了第三座斗室来居住。\par

% structure: p0114 source=h2 target=subsection reason=relative-heading-rank; base=h2

\subsection{第三十八章。}

\Emph{同一\Person{阿尔赫比乌斯}以亲手劳动偿还了母亲的债务。}\par

我以为值得传述同一人的另一件慈善之举，使我地的隐修士能由同一人的榜样学会不仅保持严格的节制，而且维持最真挚的爱情。因为他出身并非微贱之家，幼年时蔑视世俗与亲属之爱，逃到离上述城镇约四里外的隐修院，在那里度过了他的一生，以致整整五十年从未进入或看见他出身的村庄，甚至没有看过任何妇女的面容，连他自己的母亲也不例外。与此同时，他的父亲去世，留下一百\OriginalTerm{索利都斯}{solidi}的债务。虽然他本人因被剥夺了父亲的一切财产而完全免受一切烦扰，但他发现母亲因债主而极度烦恼。于是他出于孝道的思虑，略微缓和了那福音的严厉——从前在他父母兴旺时，他不承认自己在地上拥有父亲或母亲；现在他承认了自己有一位母亲，并急忙去解救她的困苦，却丝毫不放松自己所定的刻苦。因为他留在隐修院的围廊内，要求将他日常工作的任务增加两倍。在那里，他整年昼夜劳作，用他的辛劳和劳动向债主偿还了应还的债务数额，使母亲免受一切烦恼和焦虑；他以这样的方式解除她的债务负担，而不让自己所定的刻苦因孝道的必要借口而有丝毫减少。如此，他保持了惯常的刻苦，从未拒绝母亲心中孝道所要求的劳动，因为他从前为基督之爱而忽视了她，现在出于孝道的考虑又重新承认了她。\par

% structure: p0117 source=h2 target=subsection reason=relative-heading-rank; base=h2

\subsection{第三十九章。}

\Emph{论某位长者的一计，藉此在院长\Person{西默盎}无事可做时为他找到了一些工作。}\par

有一位我们非常亲爱、名叫西默盎的弟兄，完全不懂希腊文，从意大利地区来到。一位长者急于向他这位外乡人表示爱德，假借互惠之名，问他为何在斗室中无所事事；长者由此推测，他可能无法久留，因为无所事事会产生流浪的思绪，也因为他缺乏生活必需品。长者深知，唯有甘愿亲手劳作以获取食物的人，才能忍受独处的攻击。当那位弟兄回答说，他无法做或处理当地弟兄们通常做的任何事，只会写一手好字，如果埃及有谁需要拉丁文书籍的话，长者终于抓住机会，以互惠为名，实现了久已渴望的爱德工作，说道：\Quote{这个机会来自天主，因为我正找人用拉丁文为我抄写书信；我有一个兄弟，被军役的锁链束缚，精通拉丁文，我想送些圣经上的话给他，让他阅读以得造就。}于是，西默盎感恩地视此为天主赐予的机会，长者亦欣然抓住这借口，借此自由地施行爱德。他不仅立即将一整年所需之物作为生意般带来，还提供了羊皮纸及一切书写所需，随后收到抄本。这抄本在那些地方毫无用处（因为那里人人都完全不懂这种语言），且对任何人都无益，唯有借此计谋和大量花费，使一人得以不羞耻、不困惑地凭劳作与代价获得食物与生计，另一人则仿佛因债务所迫而施行了慈惠与慷慨；他为自己获得了更丰厚的酬报，相称于他为外来的弟兄热心地不仅提供了食物，还提供了书写材料和工作的机会。\par

% structure: p0120 source=h2 target=subsection reason=relative-heading-rank; base=h2

\subsection{第四十章}

\Emph{论一些少年，在把无花果带给一位病人时，未及品尝便饿死在旷野。}\par

但在我们打算谈论严格禁食和守斋的那一节中，似乎夹杂了善行和爱德之举。因此，我们再次回到主题，在这本小书中插入一件值得注意的事迹，涉及一些虽在年龄上少年，但在情感上并非如此的人。因为当有人将一些来自\Place{马雷奥蒂斯利比亚}的无花果带给\Person{若望院长}（他是\Place{斯刻特旷野}的管事），作为那里从未见过的稀罕之物时，\Person{若望}在真福司铎\Person{帕弗努提乌斯}的时代管理教会，后者将此责任托付给他。他立即派两名少年将无花果送到一位卧病在旷野深处的老人那里，那老人住在离教会约十八英里的地方。两名少年接过水果，出发前往上述老人的斗室，却完全迷失了正路——这在旷野中即使年长者也会轻易发生，因为突然起了浓雾。他们在旷野的无路荒原上游荡了一天一夜，始终找不到病人的斗室，最终因旅途疲惫、饥渴交加，屈膝祈祷，就在祈祷中把灵魂交托给天主。后来，人们沿着他们的足迹长久寻找（在这些沙地，足迹像印在雪上一样清晰，直到薄薄一层沙被微风吹动又覆盖），发现他们保存着无花果，完好如初，正如他们刚接受时那样。他们宁愿舍弃性命，也不愿放弃对所托之物的忠诚；宁愿在地上丧生，也不愿违背长上的命令。\par

% structure: p0123 source=h2 target=subsection reason=relative-heading-rank; base=h2

\subsection{第四十一章}

\Emph{论\Person{玛加略院长}关于修士行为的话：既当如同要长久活下去，又当如同每日面临死亡。}\par

还有一条\Person{真福玛加略}的宝贵训诫，我们应当提出，好让如此伟大之人的一句话来结束这本关于守斋和节制的书。他说道：一位修士应当如同要在肉身中停留一百年那样关注守斋；又应当如同每日面临死亡那样克制灵魂的激动、忘记伤害、厌恶忧愁、轻视痛苦和损失。因为在前者中，明智有益且恰当，它使修士始终以平衡的谨慎行走，不因身体衰弱而跌入最危险的深渊；在后者中，高尚的心志极为宝贵，不仅使人能轻视现世表面的顺境，也不被逆境和忧愁压垮，视其为微小琐碎之事，因为他不断将心灵的目光定睛在那处，他相信每日每时自己都将被召唤到那里。\par

\SourceRecord{Translated by C.S. Gibson.；Nicene and Post-Nicene Fathers, Second Series；Vol. 11.；Edited by Philip Schaff and Henry Wace.；Buffalo, NY: Christian Literature Publishing Co.,；1894.；<http://www.newadvent.org/fathers/350705.htm>.}

% structure: p0001 source=h1 target=section reason=page-title

\section{\WorkTitle{Institutes (Book VII)}}

% structure: p0002 source=h2 target=subsection reason=relative-heading-rank; base=h2

\subsection{第一章}

\Emph{我们如何与贪财进行外来的争战，以及这缺点在人身上如何不像其他缺点那样出于本性。}\par

我们的第三次战斗是针对贪财，我们可以将其描述为爱金钱；这是一场外来的战争，在我们本性之外，对隐修士而言，它只源于一种败坏和迟钝的心态，并且常常始于其弃绝的不圆满，以及其爱天主之心在根基上的冷淡。因为其他植根于人性的罪恶诱惑，似乎与我们与生俱来，并且以某种方式深深扎根于我们的肉体，几乎与我们出生同时，先于我们辨别善恶的能力，尽管它们在很早的时候就攻击人，却要经过长期的斗争才能被克服。\par

% structure: p0005 source=h2 target=subsection reason=relative-heading-rank; base=h2

\subsection{第二章}

\Emph{贪财之病何等危险。}\par

但这病后来临到我们，从外进入灵魂，既然越容易防备和抵抗，那么，如果忽略它，一旦让它侵入内心，对每个人就越危险，也更难驱逐。因为它成了\Quote{万恶的根}（\BibleRef{弟前 6:10}），并产生多种罪恶的诱惑。\par

% structure: p0008 source=h2 target=subsection reason=relative-heading-rank; base=h2

\subsection{第三章}

\Emph{那些我们本性上的恶习有什么用处。}\par

例如，我们不是看到那些肉体的自然冲动，不仅存在于那些尚以纯真先于善恶辨别的男孩身上，甚至也存在于幼童和婴儿身上吗？他们虽然内心连一丝情欲的接近都没有，却显示出肉体的冲动存在于他们身上，并且自然地受到激发？我们不是也看到，愤怒的致命刺伤同样已经充分存在于小孩子身上吗？在他们学会忍耐的德行之前，我们看到他们因受到委屈而烦恼，感受到甚至以玩笑方式对他们施加的冒犯；有时，尽管力量不足，但当愤怒激动他们时，报复的欲望却不缺乏。我说这些不是为了指责他们的自然状态，而是为了指出，这些从我们发出的冲动，有些是为了有益的目的植入我们里面，有些则是由于疏忽和邪恶意志的欲望而从外引入的。因为我们上面提到的这些肉欲的冲动，是造物主的上智安排为了有益的目的植入我们身体的，即：为了繁衍种族，为后代生育子女；而不是为了犯奸淫和放荡，这些也为法律的权威所谴责。愤怒的刺伤，我们不是看到它被最明智地赐给我们，以便我们对自己的罪过和错误发怒，从而更致力于德行和属灵的操练，对天主显示全部的爱，对弟兄显示全部的爱和忍耐吗？我们也知道悲伤的用处有多大，它被算在其他的恶习之中，当它被用于相反的目的时。因为一方面，当它合乎天主的敬畏时，是最需要的；另一方面，当它合乎世俗时，是最有害的；正如宗徒教导我们的，他说：\Quote{依照天主的悲怆，能成就不变的悔改，以致得救；但世俗的悲怆，却导致死亡。}（\BibleRef{格后 7:10}）\par

% structure: p0011 source=h2 target=subsection reason=relative-heading-rank; base=h2

\subsection{第四章}

\Emph{我们可以说，我们身上存在一些本性的缺点，而不冒犯造物主。}\par

因此，如果我们说这些冲动是由造物主植入我们里面的，他不会因此显得可指责，如果我们错误地滥用它们，将它们歪曲为有害的目的，并且准备通过世上无用的该隐们而悲伤，而不是通过表现忏悔和改正我们的缺点而悲伤；或者至少，如果我们不是对自己发怒（那本是有益的），而是违抗天主的命令对弟兄发怒。因为就铁而言，它是赐给我们用于良好的、有益的用途的，如果有人把它用于杀害无辜者，人们不会因此责备金属的制造者，因为人把他为幸福生活的良好和有益用途所提供的东西用于伤害他人。\par

% structure: p0014 source=h2 target=subsection reason=relative-heading-rank; base=h2

\subsection{第五章}

\Emph{论那些没有自然冲动而因我们自己的过错所沾染的恶习。}\par

但我们肯定，有些恶习的成长没有任何自然机缘产生它们，而仅仅来自败坏和邪恶意志的自由选择，如嫉妒和贪财这个罪过本身；它们（可以这么说）是从外部感染的，在本能上没有任何源于我们的起源。但是，它们越容易被防范和轻易避免，就越使被它们抓住和占据的心灵悲惨，并且几乎不让它获得治愈它的补救方法：要么是因为那些被他们本可以忽视、或避免、或轻易克服的人所伤害的人，不配通过迅速的疗愈得到医治；要么是因为，由于基础打得很差，他们不配建造德行的建筑并达到完美的顶峰。\par

% structure: p0017 source=h2 target=subsection reason=relative-heading-rank; base=h2

\subsection{第六章}

\Emph{贪财之恶一旦被接纳，多么难以驱除。}\par

因此，不要让任何人认为这个恶是无足轻重或无关紧要的：因为它越容易避免，一旦它抓住任何人，就几乎不让他获得医治它的补救方法。因为它是一个罪恶的巢穴，是\Quote{万恶的根}，并且成为无望的邪恶诱惑，如宗徒所说：\Quote{贪财（即爱金钱）是万恶的根。}（\BibleRef{弟前 6:10}）\par

% structure: p0020 source=h2 target=subsection reason=relative-heading-rank; base=h2

\subsection{第七章}

\Emph{论贪财的根源，以及它本身是哪些邪恶之母。}\par

当这种恶习抓住某位隐修士懈怠而冷淡的灵魂时，它开始以一小额金钱诱惑他，给他极好且近乎合理的借口，说明为何他应当为自己留一些钱。因为他抱怨隐修院所供给的不足，且一个健康强壮的身体几乎无法忍受。他说若疾病来临，而他并没有自己特别的储备可支持他的软弱，他该怎么办？他说隐修院的津贴微薄，且对病人极不关心；若他没有自己的东西以便照顾身体的需要，他将悲惨地死去。发给他的衣服不足，除非他准备了东西来获取另一件。最后，他说他不可能长久地留在同一地方和隐修院，除非他准备了旅费和渡海搬迁的费用，否则他不能随意移动，而被需要的强制所拘留，此后将过着悲惨而令人厌倦的生活，毫无进步；他不能不失尊严地被别人的财物供养，像一个乞丐和缺乏的人。因此，当他用这类念头自欺之后，他绞尽脑汁思考如何至少获得一文钱。然后他焦急地寻找一些特别的工作，可以在院长不知情的情况下做。秘密卖掉它，从而获得那渴望的硬币，他越来越痛苦地思考如何使它加倍：困惑于存放在何处，或托付给谁。然后他被更沉重的忧虑压迫，考虑用它买什么，或通过什么交易使其加倍。当这如愿以偿时，对黄金的更贪婪欲望萌发，并且随着金钱储备越来越多，越来越强烈地受到刺激。因为财富增加时，贪财的狂热也增加。然后他预感到长寿、衰弱的老年和各种长期疾病，这些在老年不可忍受，除非年轻时已积攒大量金钱。因此，可怜的灵魂被搅动，仿佛被蛇的圈套紧紧缠住，它试图通过更不合法的关心来增加那已非法获得的堆，而它自己生出更严重折磨它的祸患；完全沉迷于追求利益，只注意如何得到钱以便尽快逃离隐修院的纪律；只要有一线希望得到钱，它就绝不守信。为此，它不回避撒谎、作伪证、偷窃、背弃承诺、放任伤害性的暴怒之罪。如果那人从获利的希望中稍有失落，他就毫无顾忌地逾越谦逊的界限，并且在整个过程中，黄金和贪利之心成了他的神，如肚腹成了别人的神一样。因此，蒙福的宗徒审视这瘟疫的致命毒素，不仅说它是万恶的根源，还称它为拜偶像，说：\BibleQuote{而贪财（希腊文称为\OriginalTerm{贪财}{\GreekText{φιλαργυρία}}）就是拜偶像。}\BibleRef{哥 3:5} 你看，这疯狂如何一步步导致堕落，以至于借着宗徒的口它实际被宣告为拜偶像和假神，因为它越过天主的肖像和模样（一个虔诚事奉天主的人应当在自己内保持这肖像无玷），却选择爱慕和关心铸在金子上的像，而不是天主。\par

% structure: p0023 source=h2 target=subsection reason=relative-heading-rank; base=h2

\subsection{第八章。}

\Emph{贪财如何阻碍一切德行。}\par

以如此的下坡步伐，他每况愈下，最后他不在乎保留——我不说德行，而是连谦逊、爱德和服从的影子；他对一切不满，对每项工作抱怨和叹息；现在他已抛掉一切敬畏，像一匹坏脾气的马，不顾一切地、无缰地冲出去；对他每日的食物和通常的衣服不满，宣布他不能再忍受。他宣告天主不仅在那里，他的救恩也不限于那个地方；他深表遗憾地说，如果他不很快离开那里，他很快就会死。\par

% structure: p0026 source=h2 target=subsection reason=relative-heading-rank; base=h2

\subsection{第九章。}

\Emph{拥有金钱的隐修士如何不能留在隐修院。}\par

于是，他有了钱来供应自己的游荡，用这些钱他仿佛给自己装上了翅膀，现在已完全准备好搬迁，他无礼地回答所有命令，行为举止像个陌生人和访客；凡是他看到需要改进之处，他都蔑视和轻看。虽然他秘密藏有金钱的供应，却抱怨自己没有鞋也没有衣服，并对发给他这些东西如此缓慢感到愤慨。如果因上级的管理，这些东西先给了一个众所周知一无所有的人，他就更加燃起怒火，认为自己被当作陌生人蔑视。他也不满足于动手做任何工作，而是对隐修院需要做的一切事情吹毛求疵。然后他刻意寻找被冒犯和发怒的机会，免得他离开隐修院纪律显得是出于微不足道的理由。并且不满足于独自离开，免得被认为他的离去是出于自己的过错，他不停地通过秘密谈话败坏尽可能多的人。但如果恶劣天气阻碍了他的旅途和旅行，他始终处在心的悬疑和焦虑中，并且不停地散播和激起不满；因为他认为他只能从隐修院的坏品性和缺陷中找到他离开的安慰和他易变的借口。\par

% structure: p0029 source=h2 target=subsection reason=relative-heading-rank; base=h2

\subsection{第十章。}

\Emph{论隐修院的逃亡者因贪财必须承受的辛劳，尽管他过去连最轻微的工作也要抱怨。}\par

于是他被驱赶着，越来越被他的金钱之爱所燃烧；一旦得到金钱，它绝不允许隐修士留在隐修院内，也不允许他生活在会规的纪律之下。它把他像野兽一样从羊群中分离出来，使他因缺乏同伴而成为适于捕食的动物，并使他容易被吞噬，因为他被剥夺了同室之伴；它迫使他——那个一度认为履行隐修院轻微职责有失身份的人——为了获利的希望，昼夜不停地劳作；它不让他遵守祈祷的日课、斋戒的规矩、守夜的纪律；它不允许他履行合宜的代祷职责，只要他能满足贪婪的疯狂，并供应他的日常所需；它更加燃起贪欲之火，尽管他相信通过获取可以熄灭它。\par

% structure: p0032 source=h2 target=subsection reason=relative-heading-rank; base=h2

\subsection{第十一章。}

\Emph{假借保管钱袋之名，她们竟被请求与隐修士同住。}\par

因此，许多人被引向陡峭的悬崖，藉着不可挽回的堕落，跌入死亡；他们不满足于独自拥有那笔钱——无论那是他们以前从未有过的，还是他们从不良开端中扣留的——反而寻求女人与她们同住，以保存他们不正当地聚敛或保留的东西。他们使自己陷入如此多有害而危险的活动中，以至于甚至被抛入地狱的深渊，因为他们不肯同意宗徒的那句话：\Quote{只要有衣有食，就应当知足}——即隐修院节俭所供应的——反而\Quote{想要发财，就陷入诱惑和魔鬼的罗网，以及许多无益而有害的欲望，这使人沉没在败坏和灭亡中。因为贪爱钱财，}即贪婪，\Quote{乃是一切邪恶的根源；有些人贪恋钱财，就偏离了信德，并用许多痛苦刺透了自己。}\BibleRef{弟前 6:8}\par

% structure: p0035 source=h2 target=subsection reason=relative-heading-rank; base=h2

\subsection{第十二章。}

\Emph{一个冷淡的隐修士陷入贪欲之网的实例。}\par

\Emph{我知道}一个人，他自以为是个隐修士，更糟的是，他还自夸其成全。他已被接纳进入一所隐修院，院长命令他不要再把思想转回他已放弃和弃绝的事物，而要摆脱贪婪——一切邪恶的根源——和世俗的罗网；并且告诉他，如果他希望从前的私欲偏情中被洁净（他看见自己时常被这些偏情严重折磨），他就应该停止挂虑那些甚至以前也不是他自己的东西，因为被这些锁链缠住，他肯定无法在净化自己的过失上取得进步。他带着愤怒的表情，毫不犹豫地回答说：\Quote{如果你有可以用来养活别人的东西，为什么你也禁止我拥有它呢？}\par

% structure: p0038 source=h2 target=subsection reason=relative-heading-rank; base=h2

\subsection{第十三章。}

\Emph{年长者在除去罪恶之事上教导后辈的内容。}\par

但愿没有人认为这是多余或令人反感的。因为若不先解释各种罪恶的种类，并追溯疾病的根源和原因，就无法为病人提供适当的治疗药物，也无法为强壮者确保完全健康的保持。因为这些事情以及许多其他事情，常常由年长者在他们的聚会中提出，以教导年轻的弟兄们，因为他们经历过无数次的跌倒和各种人的沦亡。我们自己在许多事情上认出自己，当年长者解释并指明这些事时，作为也被同样情欲困扰的人，我们无需任何羞愧或困惑就被治愈了；因为我们无需说什么就学到了困扰我们的罪的补救方法和原因，我们忽略并缄默不语，并非出于对弟兄们的畏惧，而是唯恐我们的书偶然落入某些未受此生活方式教导的人手中，并向没有经验的人揭示那些只应让那些为达到成全高峰而辛劳奋斗的人知道的事情。\par

% structure: p0041 source=h2 target=subsection reason=relative-heading-rank; base=h2

\subsection{第十四章。}

\Emph{表明贪欲之病是三重的事例。}\par

因此，这种疾病和不健康状态是三重性的，所有教父都以同样的憎恶加以谴责。第一种特征是：我们上文描述过的玷污，它欺骗可怜的人们，使他们即使从未在世俗生活中拥有过自己的东西，也劝诱他们囤积。第二种特征是：它随后迫使人们重新拾起并再次贪恋那些他们初弃绝世俗时所放弃的东西。第三种特征是：它起源于一种错误而有害的开端，起步不良，不让那些一度被这种心灵冷淡感染的人因害怕贫穷和信心不足而脱去自己的全部世俗财物；那些扣留他们本应放弃和离弃的金钱和财产的人，它从不允许他们达到福音的全德。我们在圣经中看到这三种灾难的事例，它们都受到了不轻的惩罚。因为\Person{革哈齐}想要获得他从未有过的东西时，不仅未能获得本该从他师傅继承的预言恩赐，反而被圣\Person{厄里叟}的诅咒所覆盖，染上了永久的癞病；\Person{犹达斯}想要重新占有他跟随基督时已丢弃的财富，不仅陷入出卖主，失去了宗徒的地位，而且不被允许以众人共同的结局终结生命，却以暴死告终。\Person{阿纳尼雅}和\Person{撒斐辣}扣留了他们以前财产的一部分，在宗徒的话下被处死。\par

% structure: p0044 source=h2 target=subsection reason=relative-heading-rank; base=h2

\subsection{第十五章。}

\Emph{论错误弃绝世界者与根本不弃绝世界者之间的区别。}\par

所以，对那些说自己已弃绝此世，后来却被缺乏信德所胜、害怕失去世俗财物的人，在\WorkTitle{申命纪}中有一条奥秘的诫命：\BibleQuote{若有人畏惧胆怯，就不要上阵打仗：让他回去，回家去，免得他使弟兄们的心像他自己一样畏惧胆怯。}\BibleRef{申 20:8}还有什么比这证言更明白的呢？圣经岂不明显地宁愿他们连这修道的初阶和名号都不要接受，反比他们用自己的说服和坏榜样使别人从福音的成全中退缩，并以他们不信的恐惧削弱他人更好吗？因此，他们被命令退出战争，回到家中，因为一个人不能心怀二意打主的仗。因为\BibleQuote{三心二意的人，在他一切的行径上，易变无定。}\BibleRef{雅 1:8}并且，按照福音中的那个比喻\BibleRef{路 14:31}，他们想，那带着一万兵去对抗带着两万兵来的国王的人，是不可能作战的；他们应该趁他还在远处时就求和。也就是说，他们最好连弃绝的第一步都不要迈出，总比以后冷淡地做下去、把自己卷入更大的危险更好。因为\Quote{许愿而不还愿，倒不如不许愿。}但巧妙的是，一个带着一万兵，另一个带着两万兵。因为攻击我们的罪恶数目远远超过为我们作战的德行数目。但\BibleQuote{没有人能事奉天主而又事奉玛门。}\BibleRef{玛 6:24}并且\BibleQuote{手扶着犁向后看的，不适于天主的国。}\BibleRef{路 9:62}\par

% structure: p0047 source=h2 target=subsection reason=relative-heading-rank; base=h2

\subsection{第十六章}

\Emph{那些反对放弃自己财物的人用以庇护自己的权威。}\par

这些人在原初的贪财上试图用圣经的一些权威来辩护，他们以卑鄙的机巧来解释这些权威，渴望歪曲和曲解宗徒的话，或者更确切地说，是主自己的话，使之符合他们的目的：并且，不是使自己的生活或理解适应圣经的意义，而是使圣经的意义弯曲以适应他们自己私欲的渴望，他们试图使之符合自己的观点，并说经上记载：\BibleQuote{施予比领受更为有福。}\BibleRef{宗 20:35}并且通过对这句话的完全错误的解释，他们认为可以削弱主所说的那句话的力量：\BibleQuote{你若愿意是成全的，去，变卖你所有的，施舍给穷人，你必有宝藏在天上；然后来跟随我。}\BibleRef{玛 19:21}他们认为以此为借口，他们不必剥夺自己的财富：反而宣称，如果他们依靠原本属于他们的财富来支持，并出于自己的富余而施舍给别人，他们更为有福。当他们羞于与宗徒一起拥抱那为基督的缘故弃绝的荣耀境界时，他们也不会满足于隐修院的体力劳动和俭朴的饮食。唯一的事情是，这些人要么必须知道他们在欺骗自己，并没有真正地弃绝世界，因为他们紧紧抓住自己先前的财富；要么，如果他们真正地想要体验隐修生活，他们必须放弃并抛弃所有这些，对所弃绝的毫不保留，并且与宗徒一起夸耀\BibleQuote{在饥饿和口渴、寒冷和赤身露体中。}\BibleRef{格后 11:27}\par

% structure: p0050 source=h2 target=subsection reason=relative-heading-rank; base=h2

\subsection{第十七章}

\Emph{论宗徒和初期教会的弃绝}\par

试想，这人（他自称出生时就拥有罗马公民特权，以此证明他在世俗等级中并非卑微之人）难道不也可以靠先前属于他的财产来维持生活吗？又试想，那些在耶路撒冷拥有田产房屋的人，变卖一切，自己一点也不保留，把价银带来放在宗徒脚前；倘若宗徒们认为这是最佳方案，或他们自己认为这样更好，难道他们不也可以从自己的财产中供应身体所需吗？但他们立时放弃全部财产，宁愿靠自己的劳作和向外邦人募捐来生活——圣宗徒在致罗马人的书信中论及外邦人的捐献，表明他自己在这件事上的职责，也激励他们同样进行捐献：\BibleQuote{但我现在往耶路撒冷去，为圣人们服务。因为马其顿和阿哈雅人乐意凑出捐项，给耶路撒冷的贫困圣人；他们实在乐意，且是他们的债户。因为外邦人既分享了他们的神性事物，也该在肉体事物上服事他们。} \BibleRef{罗 15:25} 他对格林多人也同样关心此事，更催促他们在自己到达之前积极准备捐款，他打算派人送去供给他们的需要。\BibleQuote{关于为圣人的捐款，我怎样吩咐了迦拉达众教会，你们也要照样做。每遇一周的第一日，你们每人要各自积蓄，留出他所乐意捐献的，免得我来时才现凑。但当我来到时，你们所认可的人，我就派他们带着你们的恩宠送到耶路撒冷。} 他又鼓励他们多捐，便加上：\BibleQuote{若我也应当去，他们就与我同去。} \BibleRef{格前 16:1} 意思是，如果你们的奉献值得由我来送。对迦拉达人，他也作证说，当他和宗徒们分配宣讲职务时，他与雅各伯、伯多禄、若望如此安排：由他向犹太人传福音，但他绝不推诿对耶路撒冷穷人的关怀和挂虑——这些人为了基督放弃一切，甘愿过贫穷生活。\BibleQuote{他们一看见，}他说，\BibleQuote{那赐给我的恩宠，雅各伯、刻法和若望——那些被视为柱石的——就向我和巴尔纳伯伸出右手，以示共融，好叫我们向外邦人传福音，而他们向受割礼的人；只要求我们记念穷人。}他又说他特别留意这事，说：\BibleQuote{我自己也竭力这样做了。} \BibleRef{迦 2:9} 那么，谁更有福呢？是那些刚刚从外邦人中被召选出来，未能登上福音完美的顶峰，固守自己的财产，宗徒认为他们至少能不敬拜偶像、不犯奸淫、不吃勒死牲畜和血（\BibleRef{宗 15:20}）并连一切财产接受了基督信仰的人呢？还是那些按福音的要求生活，每天背负主的十字架，不愿为自己的使用保留任何财产的人呢？如果蒙福的宗徒自己，被锁链和脚镣束缚，或因旅途困难所阻，不能像往常一样亲手劳动供给饮食，而声明他从马其顿来的弟兄们供应了他的所需——\BibleQuote{我缺乏的，}他说，\BibleQuote{马其顿来的弟兄们补足了：} \BibleRef{格后 11:9} ——他又对斐理伯人说：\BibleQuote{你们斐理伯人也知道，在福音初期，当我从马其顿出发时，除了你们以外，没有别的教会在施与受的事上与我相通；因为甚至在得撒洛尼，你们也一再派人供应我的需用。} \BibleRef{斐 4:15} （如果是这样）那么，按照这些人心中的冷淡想法，那些人是否真的比宗徒更有福，因为发现他们曾以财物服事了他？但没有人敢这样断言，无论他多么愚蠢。\par

% structure: p0053 source=h2 target=subsection reason=relative-heading-rank; base=h2

\subsection{第十八章}

\Emph{如果我们想要效法宗徒，就不应按自己的规定生活，而应跟随他们的榜样。}\par

\Emph{因此}如果我们愿意服从福音的诫命，并显明自己是宗徒和整个初期教会，或那些在我们时代继承了他们德行和完美的教父的追随者，我们就不应安于自己的规定，许诺自己从这可怜不冷不热的状态中获得成全；而应追随他们的脚步，绝不以谋求自己的利益为目标，而应寻求隐修院的纪律和制度，以便真正弃绝这个世界；不保留因信德不足的诱惑而蔑视的任何事物；并应指望以我们自己的劳作来获取日用的食物，而非从自己的金钱储备中支取。\par

% structure: p0056 source=h2 target=subsection reason=relative-heading-rank; base=h2

\subsection{第十九章}

\Emph{凯撒勒雅主教圣巴西略针对辛克雷提乌斯的一句名言}\par

流传着凯撒勒雅主教圣巴西略针对一位叫辛克雷提乌斯的人的一句名言；此人正因我们所说的那种不冷不热而变得漠不关心；他虽然声称弃绝了世界，却为自己保留了一些财产，不喜欢靠双手劳作来维持生活，也不愿通过剥夺自身、艰苦劳作和服从隐修院来获得真正的谦卑。他说：\Quote{你毁掉了辛克雷提乌斯，而没有造就一位隐修士。}\par

% structure: p0059 source=h2 target=subsection reason=relative-heading-rank; base=h2

\subsection{第二十章}

\Emph{被贪婪战胜是多么可耻}\par

因此，若我们想在属灵争战中合法地奋斗，就也要将这危险的仇敌从我们心中驱逐出去。因为战胜他并不多么显示伟大的德行，而被其击败则是可耻和丢脸的。当你被一个强壮的人压倒时，虽因被打倒而有悲伤，因失去胜利而有痛苦，但被击败者仍可从对手的力量中得到些许安慰。但如果敌人是个可怜的东西，争斗也微弱无力，那么除了失败后的悲伤，还有更可耻的混乱，以及比损失更糟糕的羞耻。\par

% structure: p0062 source=h2 target=subsection reason=relative-heading-rank; base=h2

\subsection{第二十一章。}

\Emph{如何战胜贪婪。}\par

在这方面，最大的胜利和持久的凯旋将是：如所说的那样，隐修士的良心不被最小的硬币所有物玷污。因为一个在小所有物上失败、已经将邪恶欲望之根引入心中的人，不可能不立即被更大的欲望之火点燃。因为基督的士兵将是胜利而安全的，摆脱欲望的一切攻击，只要这最恶毒的灵没有在他心中植入这种欲望的种子。因此，虽然在所有种类的罪中，我们通常应注视蛇的头 ，\BibleRef{创 3:15}，但在这方面我们更应加倍警惕。因为一旦被接纳，它就会自行滋长，并为自己点燃更烈的火。因此，我们不仅要提防\Emph{拥有}钱财，还必须从我们的灵魂中驱逐对钱财的\Emph{欲望}。因为我们不应只避免贪婪的结果，而应从根上斩断对它的倾向。因为若我们心中存在获取钱财的欲望，则没有钱财也是无益的。\par

% structure: p0065 source=h2 target=subsection reason=relative-heading-rank; base=h2

\subsection{第二十二章。}

\Emph{实际上没有钱也可能被视为贪婪者。}\par

因为一个没有钱的人也可能完全不摆脱贪婪的病症，而贫穷的恩赐对他毫无益处，因为他未能根除贪财之罪：他喜悦于贫困的益处，而非德行的功绩；满足于必需的负担，而非缺乏内心的热忱。正如福音之言论及那些身体未受玷污的人，说他们在心中是犯奸淫者，\BibleRef{玛 5:28}；同样，那些从未受钱财重量压迫的人，也可能在心态和意图上与贪婪者一同被定罪。因为在他们身上缺乏的只是拥有的机会，而非意愿；而天主常赏报意愿，而非勉强。因此，我们必须竭尽全力，免得我们劳苦的果实白白被毁。因为忍受贫困和匮乏的效果，却因破碎意志的过失而失去其果实，是件可悲的事。\par

% structure: p0068 source=h2 target=subsection reason=relative-heading-rank; base=h2

\subsection{第二十三章。}

\Emph{从犹达斯的案例中引用的例子。}\par

你想知道那刺激（若未精心根除）会如何危险而有害地为它主人的毁灭发芽，并生出各种不同罪的枝子吗？请看犹达斯，他原被列入宗徒之中，看看因他未压碎这蛇的致命头颅，它如何以毒液毁灭了他，以及当他被贪欲的网罗捉住时，如何驱使他犯罪并急速堕落，以致他受劝说以三十块银钱出卖世界的救赎者和人类救恩的创始者。他若未被贪婪之罪污染，决不可能被驱动犯下这卖主的重罪；若非他先习惯于偷窃所托付给他的钱囊，他也不会使自己邪恶地犯下出卖主的罪。\par

% structure: p0071 source=h2 target=subsection reason=relative-heading-rank; base=h2

\subsection{第二十四章。}

\Emph{除非完全舍弃自己的一切，否则无法战胜贪婪。}\par

这是一个相当可怕而明显的暴政实例：一旦心灵被它俘获，它就不允许心灵遵守任何诚实的规则，也不满足于任何所得的增加。因为我们寻求结束这疯狂，不是靠财富，而是靠舍弃它们。最后，当犹达斯接受了为分给穷人所设置并托付给他保管的钱囊时，他本可以至少以充足的金钱满足自己，并限制他的贪欲，但丰裕的供应只是爆发成更贪婪的欲望刺激，以致他不再准备偷偷偷窃钱囊，而是实际上出卖主自己。因为这种贪婪的疯狂不会因任何财富的充足而满足。\par

% structure: p0074 source=h2 target=subsection reason=relative-heading-rank; base=h2

\subsection{第二十五章。}

\Emph{论阿纳尼雅和撒斐辣以及犹达斯的死亡，他们是因贪婪的冲动而丧亡。}\par

最后，众宗徒之首，由这些实例所教导，深知凡有贪财之心者无法自制，且无论钱财多寡均不能遏止此心，唯有靠弃绝一切之德方能止息，遂以死刑惩罚了前述的阿纳尼雅和撒斐辣，因为他们从自己的财产中扣留了一部分，为使犹达斯因出卖主之罪而自缢的死亡，他们也因贪婪说谎而遭受。\EditorialNote{宗 5} 二者的罪恶与惩罚何其相似！在前者，不忠是贪财的结果；在后者，谎言是贪财的结果。在前者，真理被出卖；在后者，犯下说谎之罪。虽然他们行为的结局看似不同，但都指向同一目的。因为前者为逃避贫穷，想要收回他已弃绝之物；后者为避免贫穷，试图从自己的财产中扣留一些，而这些财产本应真诚地献于宗徒，或全然施予弟兄。因此，二者均招致死亡的判决，因为每一罪恶皆从贪财之根生出。既然如此，对于那些不觊觎他人财物，而只是试图节省自己所有，且无意\Emph{获取}、只愿\Emph{保留}的人，尚且发出如此严厉的判决，那么对于那些从未拥有过自己财物却渴望积聚财富，在人前装作贫穷，而在天主面前因贪欲之情而被判为富有的人，又该如何看待呢？\par

% structure: p0077 source=h2 target=subsection reason=relative-heading-rank; base=h2

\subsection{第二十六章}

\Emph{贪财使灵魂患上属灵的癞病}\par

这样的人在灵性和心灵上都被视为癞病人，仿效革哈齐的样式；革哈齐因贪恋此世不确定的财富，被可憎的癞病沾染，由此为我们留下一个明确的例证：每一个被贪欲之污秽玷污的灵魂，都被罪恶的属灵癞病所覆盖，在天主面前被算为不洁，永受诅咒。\par

% structure: p0080 source=h2 target=subsection reason=relative-heading-rank; base=h2

\subsection{第二十七章}

\Emph{圣经证明：追求成全者受教不可再取回已放弃和弃绝之物。}\par

倘若你因追求成全之心，已弃绝一切并跟随那对你说\Quote{你去变卖你所有的一切，施舍给穷人，你必有宝藏在天上；然后来跟随我}的基督，\BibleRef{玛 19:21} 那么，你既已把手扶在犁上，为何向后看，以致被同一主的声音宣告不配进入天国？\BibleRef{路 9:62} 当你在福音的屋顶上安全时，为何下来从屋里取走某物，即那些你从前轻视之物？当你出到田野、操练德行时，为何跑回去，试图重新穿上你弃绝时所脱去属于此世的东西？\BibleRef{路 17:31} 但若你因贫穷所阻，无可放弃之物，你更不应积聚你从未拥有之物。因为主之恩宠正是为此预备了你，使你毫无财富之羁绊，更轻快地奔向祂。但欠缺此德者不必灰心，因为没有人无可放弃之物。凡彻底根除占有欲望的人，就是弃绝了此世一切所有。\par

% structure: p0083 source=h2 target=subsection reason=relative-heading-rank; base=h2

\subsection{第二十八章}

\Emph{惟有赤身裸体、一无所剩，方能战胜贪财。}\par

这就是对贪财的圆满胜利：不让最微小的贪念余烬存于心中，因为我们知道，若在心中哪怕保留一丁点火星，就将无力扑灭它。\par

% structure: p0086 source=h2 target=subsection reason=relative-heading-rank; base=h2

\subsection{第二十九章}

\Emph{隐修士如何保持贫穷。}\par

惟有当我们留在隐修院内，并如宗徒所说，\Quote{只要有衣有食，就当知足}\BibleRef{弟前 6:8}，才能保全此德不受损害。\par

% structure: p0089 source=h2 target=subsection reason=relative-heading-rank; base=h2

\subsection{第三十章}

\Emph{贪财之病的救治方法。}\par

因此，我们当记念阿纳尼雅和撒斐辣的审判，畏惧扣留我们已放弃并誓愿全然弃绝之物。我们也当畏惧革哈齐的榜样，他为贪财之罪受到永久的癞病惩罚。从此，我们当警惕获取我们从前从未拥有过的财富。此外，我们也当畏惧犹达斯的过犯和死亡，尽全力避免取回我们曾抛开的任何财富。尤其要思想我们软弱善变的本性，提防主的日子像夜间的贼一样临到我们，\BibleRef{得前 5:4} 使我们的良心因一文钱而被玷污；因为这将使我们弃绝世界的所有果实化为乌有，并使主在福音中对那富人说的话也转向我们：\Quote{糊涂人哪！今夜就要索回你的灵魂，你所备置的一切，将归谁呢？}\BibleRef{路 12:20} 并且，不为明天忧虑，我们绝不容许自己被诱离隐修院的规条。\par

% structure: p0092 source=h2 target=subsection reason=relative-heading-rank; base=h2

\subsection{第三十一章}

\Emph{除非留在隐修院内，无人能战胜贪财；以及如何留在那里。}\par

但我们必然不被容许这样做，甚至不能留在任何规矩之下，除非那只能源于谦卑的忍耐之德，首先在我们内稳固确立。因为前者教训我们不要烦扰他人；后者教训我们以宽宏之心忍受他人加给我们的侮辱。\par

\SourceRecord{Translated by C.S. Gibson.；Nicene and Post-Nicene Fathers, Second Series；Vol. 11.；Edited by Philip Schaff and Henry Wace.；Buffalo, NY: Christian Literature Publishing Co.,；1894.；<http://www.newadvent.org/fathers/350707.htm>.}

% structure: p0001 source=h1 target=section reason=page-title

\section{《论纲（第八卷）》}

% structure: p0002 source=h2 target=subsection reason=relative-heading-rank; base=h2

\subsection{第一章}

\Emph{我们第四场战斗是如何对抗愤怒之罪，以及这种情欲会产生多少祸患。}\par

在我们的第四场战斗中，必须将愤怒的致命毒害从我们灵魂的最深处彻底根除。因为只要这毒害存留于我们心中，并以它有害的黑暗蒙蔽灵魂的眼睛，我们就既不能获得正确的判断与明辨，也不能获得由正直目光所生的洞察力或多谋善断的成熟；我们不能分享生命，也不能持守正义，甚至不能拥有属灵与真光的能力：\Quote{因为，}有人说，\Quote{我的眼睛因愤怒而昏花。}我们也无法获得智慧，即使举世公认我们为智者，因为\Quote{愤怒安卧于愚人怀中。}我们甚至不能得到永生，即使众人都认为我们谨慎，因为\Quote{愤怒甚至摧毁谨慎的人。}我们也无法以纯洁的心灵判断，获得正义的主权，即使所有人都认为我们完美圣洁，因为\Quote{人的愤怒并不成就天主的正义。}\BibleRef{雅 1:20}我们也绝不可能获得那种甚至常现于世人的尊荣，即使我们因出身高贵而被视为高贵尊荣，因为\Quote{愤怒的人受羞辱。}我们也不能获得任何成熟的计谋，即使我们显得庄重且具备最高学识，因为\Quote{愤怒的人行事不凭计谋。}我们也无法免于危险的纷扰，也不能无罪，即使无人给我们带来任何纷扰，因为\Quote{暴躁的人挑起争执，愤怒的人挖掘罪过。}\par

% structure: p0005 source=h2 target=subsection reason=relative-heading-rank; base=h2

\subsection{第二章}

\Emph{论那些说愤怒并无害处的人，如果我们对做错事的人发怒，因为天主自己也被说成是愤怒的。}\par

我们曾听见有些人试图为灵魂这种最有害的疾病辩护，以一种颇为骇人的解经方式竭力减轻其严重性：他们说，如果对我们做错事的弟兄发怒，并无害处，因为天主自己也被说成向那些要么不愿认识祂，要么认识祂却蔑视祂的人发怒且满溢义怒，正如这里所说：\Quote{上主的怒气向祂的子民发作；}或者先知祈祷说：\Quote{上主，求你不要在你的愤怒中责罚我，也不要在你的怒气中惩戒我。}他们不明白，当他们想要为一种最恶毒的罪向人打开借口时，他们是将人性的情欲污点归给了神圣的无限者和一切纯洁的泉源。\par

% structure: p0008 source=h2 target=subsection reason=relative-heading-rank; base=h2

\subsection{第三章}

\Emph{论那些以拟人方式论及天主的事。}\par

因为如果当这些事被论及天主时，我们按照粗俗的字面意义来理解，那么祂也\Emph{睡觉}，如所说的：\Quote{上主，起来，你为什么睡眠？}尽管关于祂在其他地方却说：\Quote{看，那护卫以色列的，不打盹，也不睡眠。}祂也\Emph{站立}和\Emph{坐着}，因为祂说：\Quote{天是我的宝座，地是我的脚凳：}\BibleRef{依 66:1}尽管\Quote{祂用手量了天，用拳握了地：}\BibleRef{依 40:12}祂也\Quote{醉酒}，如所说的：\Quote{上主像睡眠者醒来，像一位因酒而醉的勇士；}祂\Quote{那独享不死不灭，住在不可接近的光中者：}\BibleRef{弟前 6:16}更不必说圣经中常提及的\Quote{无知}与\Quote{遗忘}；最后还有祂肢体的轮廓，被描述为像人一样排列有序，例如毛发、头、鼻孔、眼、脸、手、臂、指、腹、足。如果我们愿意按纯粹的字面意义来理解这一切，我们就必须认为天主有肢体轮廓和形体样式，这甚至是令人震惊的言词，必须远离我们的思想。\par

% structure: p0011 source=h2 target=subsection reason=relative-heading-rank; base=h2

\subsection{第四章}

\Emph{我们应当以何种意义理解归诸不变无形之天主的情欲和人类技艺。}\par

因此，若没有可怕的亵渎，这些事不能按字面理解那一位——祂被圣经的权威宣告为不可见、不可言说、不可理解、不可估量、纯一、非复合的；同样，愤怒和震怒的激情也不能归于那不变的本性而不致于可怕的亵渎。因为我们必须看到，肢体象征天主的德能和无限作为，这些只能以我们所熟悉的肢体表述来呈现：藉着口，我们应理解祂的言语，即祂的仁慈不断倾注于灵魂隐秘感官的言语，或祂在我们祖先和先知中所说的；藉着眼睛，我们能理解祂视力的无限性，祂以此看见并洞察一切，因此我们所做、所能做、甚至所想的一切，没有一样能向祂隐藏。藉着\Quote{手}，我们理解祂的眷顾和作为，藉此祂是万物的创造者和作者；手臂是祂权能和治理的象征，祂以此托住、统治和控制万物。且不说其他，祂头上的白发除了象征神性的永恒和永存，藉此祂没有任何开端、先于一切时间、超越所有受造物之外，还能是什么？那么同样，当我们读到主的愤怒或震怒时，不应\OriginalTerm{拟人化地}{\GreekText{ἀνθρωποπαθῶς}}理解，即按人类激情的不当意义，而应按合乎天主的方式理解，祂超越一切激情；这样我们应明白祂是这世上所有不义之事的审判者和报复者；并由于这些术语及其含义，我们应敬畏祂作为我们行为的可畏赏报者，并恐惧不敢行任何违背祂旨意的事。因为人性通常畏惧那些知道会发怒的人，并害怕冒犯他们；如同在最公正的审判者面前，报复的忿怒常被那些良心受某种指控折磨的人所畏惧；这并非说这种激情存在于那些以完全公正施行审判之人的心中，而是当他们如此畏惧时，审判者对他们的态度是公正不偏执法的前兆。而这种态度，无论多么仁慈温和，在那些将受公正惩罚的人看来，却被视为最野蛮的忿怒和激烈的愤怒。若要解释圣经中以人的形像隐喻谈论天主的全部事物，那将冗长且超出本书范围。对我们目前的宗旨——针对愤怒之罪——而言，说了这些就够了，免得有人因无知而从那些本应寻求圣德和永生作为带来生命和救恩之药方的圣经中，为自己招来这罪恶和永死的原因。\par

% structure: p0014 source=h2 target=subsection reason=relative-heading-rank; base=h2

\subsection{第五章。}

隐修士应当如何平静。\par

因此，一个追求成全、渴望在属灵战斗中合法争战的隐修士，应摆脱一切愤怒和震怒之罪，并听从那\Quote{蒙选器皿}对他的吩咐。他说：\BibleQuote{一切愤怒、震怒、喧嚷和毁谤，连同一切恶毒，都该从你们中间除掉。}\BibleRef{弗 4:31}当他说\BibleQuote{一切愤怒都该从你们中间除掉}时，他不排除任何愤怒，无论必要或有益于我们。若有需要，他应立即如此对待一个犯错的弟兄，以致当他设法给可能染了小热病的人施以治疗时，自己不致因愤怒而陷入更危险的盲目之病。因为想要医治他人伤口的人，自己必先健康，摆脱一切软弱的情绪，免得福音那句话用在他身上：\BibleQuote{医生，先治好你自己吧；}\BibleRef{路 4:23}免得他看见弟兄眼中的木屑，却看不见自己眼中的大梁，因为他自己眼中有愤怒的大梁，又怎能看见拔出弟兄眼中的木屑呢？\BibleRef{玛 7:3}\par

% structure: p0017 source=h2 target=subsection reason=relative-heading-rank; base=h2

\subsection{第六章。}

论义怒与不义之怒。\par

几乎从每一种原因，愤怒的情绪都会沸腾，并蒙蔽灵魂的眼睛，将一种更严重疾病的致死大梁放在我们视力的锐利之上，阻止我们看见正义的太阳。无论我们眼皮上盖的是金板、铅板还是任何你喜欢的金属，金属的价值并不影响我们的盲目。\par

% structure: p0020 source=h2 target=subsection reason=relative-heading-rank; base=h2

\subsection{第七章。}

论愤怒唯一有益于我们的情况。\par

必须承认，愤怒被极好地植入我们内，我们接纳它仅在一个情况下有益且有利，即当我们对我们心中情欲的冲动感到愤慨和震怒时，并因那些我们在人前羞于做或说的事却在我们心灵的隐秘处兴起而烦恼，因为我们战栗于天使和贯穿万有的天主本身的临在，并以极度的恐惧畏惧祂的眼目——我们的心中秘密不可能向祂隐藏。\par

% structure: p0023 source=h2 target=subsection reason=relative-heading-rank; base=h2

\subsection{第八章。}

从真福达味的生活中看出愤怒被正确感受的例子。\par

而且无论如何（事实如此），当我们对这股怒气本身感到激动，因为它悄然潜入我们心中，针对我们的弟兄，我们便在愤怒中驱逐其致命的煽动，不容许它在心底有危险的藏身之处。那位先知甚至教导我们以这种方式发怒，他如此彻底地从情感中驱逐了怒气，甚至不愿报复他的敌人和那些由天主交到他手中的人：他说\BibleQuote{你们发怒，却不可犯罪}。因为当他渴望白冷井的水，他的勇士们穿过敌营取来给他时，他立刻将水倒在地上；就这样，他在愤怒中熄灭了愿望的快感，将水奠于上主前，没有满足他所表达的渴望，说：\BibleQuote{我决不能这样做！我岂可喝那些冒着生命危险去的人的血呢？}\BibleRef{撒下 23:17}当史米当着众人面前向达味王扔石头并咒骂他，而阿彼瑟，责鲁雅的儿子，军长，想要砍下史米的头，为国王的侮辱报仇时，有福的达味以虔诚的愤怒反对这一可怕的建议，保持着谦卑的适当尺度和严格的忍耐，以沉静温和说：\BibleQuote{责鲁雅的儿子们，我与你们有什么相干？让他咒骂吧！因为上主吩咐他咒骂达味。谁敢说：你为什么这样做？看，我的亲生子还要我的命，更何况这本雅明人呢？由他吧，让他咒骂，因为上主命令了他。也许上主会垂顾我的苦难，今日的咒骂反会使我得到好处。}\BibleRef{撒下 16:10}\par

% structure: p0026 source=h2 target=subsection reason=relative-heading-rank; base=h2

\subsection{第九章}

\Emph{论应针对自己的愤怒。}\par

有些人奉命以一种有益的方式\BibleQuote{发怒}，但却是针对我们自己以及生起的恶念，\BibleQuote{不可犯罪}，即不使它们导致恶果。最后，下一节更清楚地解释了这层含义：\BibleQuote{你们在床上要心中思念，并要悔改。}也就是说，无论你们心里想到什么，当突然而紧张的激动涌来时，要以有益的悔改纠正和改善，仿佛躺在安息的床上，以劝言的缓和力量消除所有恼怒和激动。最后，蒙福的宗徒在引用这节经文时说：\BibleQuote{你们发怒，却不可犯罪；不可让太阳在你们的愤怒中落下，也不可给魔鬼留地步。}\BibleRef{弗 4:26}如果在我们的愤怒上让正义的太阳落下是危险的，并且当我们发怒时，立刻就为魔鬼在心里留地步，那么他为何又吩咐我们发怒，说：\BibleQuote{发怒，却不可犯罪}呢？他岂不是明说：要对你们的过犯和脾气发怒，免得你们若顺从它们，正义的太阳基督会因你们的愤怒开始在你们昏暗的理智中落下，他一离去，你们便在心里给魔鬼留了地步吗？\par

% structure: p0029 source=h2 target=subsection reason=relative-heading-rank; base=h2

\subsection{第十章}

\Emph{论太阳，经上说它不应在你们的愤怒中落下。}\par

天主通过先知明确提到了这太阳，他说：\BibleQuote{但为敬畏我名的人，正义的太阳要升起，以它的翅膀带来痊愈。}\BibleRef{拉 4:2}而这太阳又被说成在正午\BibleQuote{落下}，临到罪人、假先知和发怒的人，正如先知所说：\BibleQuote{他们的太阳在正午落下。}\BibleRef{亚 8:9}而且无论如何，从比喻上说，心灵，即\OriginalTerm{心灵}{\GreekText{νοῦς}}或理性，它被合理地称为太阳，因为它照看所有思想的思虑和分辨，不应因愤怒之罪而熄灭；免得当它\BibleQuote{落下}时，扰乱的黑影连同其作者魔鬼充满我们内心的所有情感，并被愤怒的黑影淹没，如同在昏暗的夜晚，我们不知道该如何行。正是因为这个意义，我们引用宗徒的这段经文，它是藉着长上的教导传给我们的，因为有必要——即便冒着长篇大论的风险——表明他们对愤怒的看法，他们根本不允许愤怒甚至片刻进入我们的心；他们极其谨慎地遵守福音的这句话：\BibleQuote{凡向弟兄发怒的，应受裁判。}\BibleRef{玛 5:22}但若许我们发怒直到日落，那么在这太阳西沉之前，我们愤怒的饱足和报仇就能充分纵容情欲和危险的激动。\par

% structure: p0032 source=h2 target=subsection reason=relative-heading-rank; base=h2

\subsection{第十一章}

\Emph{论那些甚至连日落也不能为其愤怒设定限制的人。}\par

但我该如何说那些人呢（我这样说时自己也感到羞愧），连日落都不能限制他们的不宽容：相反，他们将愤怒延续数日，对激怒他们的人怀怨在心，口头上说他们不生气，但在行动上却表现出极为不安？因为他们不与这些人愉快地交谈，也不以通常的礼貌对待他们，并且认为这样做没有错，因为他们并未寻求为所受的困扰报仇。但由于他们不敢，或至少不能公开表露愤怒并放纵它，他们反而将愤怒之毒吸进体内，对己有害，并在心中秘密怀藏，暗自品味；不用心智的努力摆脱愠怒的性情，而是随着日子流逝将其消化，并在一段时间后稍加缓和。\par

% structure: p0035 source=h2 target=subsection reason=relative-heading-rank; base=h2

\subsection{第十二章}

\Emph{论脾气和愤怒如何以尽其所能付诸行动而告终。}\par

但看来连这一点对每人来说也并非复仇的终结；有些人只有尽其所能地发泄愤怒的冲动，才能完全满足他们的怒气或愠怒；我们知道，那些抑制自己情感的人，并非出于平息情感之愿，而仅仅是缺少报复的机会，便是如此。因为他们对那些激怒他们的人所能做的，无非就是说话时不带通常的礼貌：或者，看来愤怒只能在行为中被节制，而不能从其隐藏在我们心中的巢穴里被彻底根除；以致被其阴影笼罩，我们不仅无法接纳有益忠告和知识的光芒，而且也无法成为圣神的殿宇，只要愤怒的灵居住在我们内。因为心中滋养的恼怒，即便不伤害身边之人，却与公开表露的恼怒一样，排斥圣神光辉的照耀。\par

% structure: p0038 source=h2 target=subsection reason=relative-heading-rank; base=h2

\subsection{第十三章}

\Emph{我们连一刻也不应保留愤怒。}\par

\Quote{若你在祭台前献礼物时，想起你的兄弟有什么怨你的事，就把礼物留在祭台前，先去与你的兄弟和好，然后再来献你的礼物。} \BibleRef{玛 5:23} 那么，我们怎能对兄弟怀怨不放呢？我不是说几天，而是甚至到日落，倘若他还怨我们，我们就不可以向上主献上祈祷？然而宗徒命令我们：\Quote{不断祈祷；} \BibleRef{得前 5:17} 并\Quote{在每一处举起圣洁的手，不发怒，不争论。} \BibleRef{弟前 2:8} 所以，要么我们心中保留这种毒汁，永不祈祷，从而在宗徒或福音的诫命上成为有罪的（这诫命命令我们随处不断地祈祷），要么我们自欺欺人，违背祂的命令冒然献上祈祷，我们必须知道，我们向天主献上的不是祈祷，而是倔强悖逆的态度。\par

% structure: p0041 source=h2 target=subsection reason=relative-heading-rank; base=h2

\subsection{第十四章}

\Emph{论与兄弟和好。}\par

因为常常我们蔑视那些受伤忧愁的兄弟，鄙视他们，并说他们并非因我们的过错而受伤，灵魂的医治者知道一切秘密，愿意彻底根除心中一切愤怒的起因，不仅命令我们若受了亏待要宽恕，要与兄弟和好，并不要记住他们造成的错失或伤害，而且祂也给了类似的指令：倘若我们知道他们对我们有什么怨尤，无论有理无理，我们应放下礼物，即推迟祈祷，先去向他们赔补，如此在兄弟的医治先完成之后，我们才能毫无玷污地献上祈祷的礼物。因为众人的共同主并不那么看重我们的敬拜，以至于允许在我们心中愤怒作王，从而在一人身上失却从另一人身上得到的。因为任何一人的损失都使祂遭受损失，祂愿意用同样的方式寻求所有祂的仆人的救恩。因此，如果我们的兄弟对我们有什么怨尤，我们的祈祷将失去效果，就像我们自己以膨胀愤怒的灵怀藏着对他苦毒的情绪一样。\par

% structure: p0044 source=h2 target=subsection reason=relative-heading-rank; base=h2

\subsection{第十五章}

\Emph{旧律如何不仅从行为也从未念根除愤怒。}\par

\Quote{你不可在心里恨你的兄弟；} 又说：\Quote{不可记念你人民的亏负；} \BibleRef{肋 19:17} 又说：\Quote{那些保存对错事的记念的道路是通向死亡的。} 你看，那里邪恶不仅被限制在行为上，也被限制在秘密的思念中，因为命令要从心里完全根除仇恨，不仅要禁止报复，还要禁止对所受的伤害的记念。\par

% structure: p0047 source=h2 target=subsection reason=relative-heading-rank; base=h2

\subsection{第十六章}

\Emph{那些不放弃恶习的人退隐是何等无用。}\par

\Emph{有时}当我们被骄傲或急躁胜过，想要改善粗暴野蛮的举止时，我们抱怨需要独处，好像在那里无人激怒我们就能找到忍耐之德；我们为自己的疏忽道歉，并说我们不安的原因并非来自自己的急躁，而是来自弟兄的过失。当我们把过失的责任推给别人时，我们就永远无法达到忍耐和成全的目标。\par

% structure: p0050 source=h2 target=subsection reason=relative-heading-rank; base=h2

\subsection{第十七章}

\Emph{我们心灵的平安不取决于别人的意愿，而是在于我们自己的控制。}\par

那么我们修德和心灵平安的主要部分不应取决于别人的意愿（它不可能受我们权威管辖），而更应在于我们自己的控制。因此，我们不发怒不应出于别人的成全，而应出于我们自己的德行；这德行不是靠别人的忍耐获得，而是靠我们自己的坚忍。\par

% structure: p0053 source=h2 target=subsection reason=relative-heading-rank; base=h2

\subsection{第十八章}

\Emph{论我们应如何热切追求旷野，以及在那里我们进步的事物。}\par

此外，惟有那些达致成全、从一切过失中净化的人，才应当寻求旷野；他们在弟兄团体中彻底根除了所有过失之后，进入旷野，不是为了怯懦的逃避，而是为了神圣的默观，并怀着对天上事物更深领悟的渴望——这只能在独居中为成全者所获得。因为我们带入旷野而未治愈的任何过失，都会发现它们隐藏在我们内，无法去除。正如当品格得到改善时，独处能向它开启最纯洁的默观，并让灵性奥秘的知识在其清澈的注视中显现；同样地，对于未经纠正的人，独处通常不仅保存而且加剧他们的过失。因为一个人自以为忍耐和谦逊，只在他与人交往中未遇到任何人时；但只要有机会出现任何形式的激情，他就会立刻恢复本性：我是指那些隐藏的过失会立即浮出表面，像长期闲散期间精心喂养的脱缰之马，更加急切而猛烈地冲出栅栏，导致驾驭者的毁灭。因为当与人实践这些过失的机会被移除时，我们的过失在我们内会更加增长，除非我们先从它们中被净化。而那忍耐的影子——当我们与弟兄们交往时，出于对他们的尊重和公开场合，我们似乎幻想地拥有它——因懒惰和疏忽而完全失去。\par

% structure: p0056 source=h2 target=subsection reason=relative-heading-rank; base=h2

\subsection{第十九章}

\Emph{一个帮助判断那些只在未受任何人试探时才忍耐之人的比喻。}\par

但这就像所有有毒的蛇类或野兽，当它们留在僻静处和自己的巢穴中时，仍然并非无害；因为它们实际上并未伤害任何人，所以不能真正被称为无害。这是因为它们的状态并非出于任何良善的感觉，而是出于独处的需要；当它们获得伤害某人的机会时，立刻就会产生储存于内的毒素，并显示其本性的凶残。同样，在追求成全的人身上，仅仅不对\Emph{人}发怒是不够的。因为我们记得，当我们生活在独处时，一种恼怒的情绪会悄悄侵入我们，对象是我们的笔——因为它太大或太小；或我们的削笔刀——因为它切割不锋利，用钝刃割我们想割的东西；或一块燧石——如果碰巧我们稍迟、匆忙赶去阅读时，它迸出火星，以致我们无法除去和摆脱内心的烦乱，除非咒骂那无生命的物质，或至少咒骂魔鬼。因此，为了成全之道，缺乏可供我们发泄怒气的人是没有用的：因为如果忍耐尚未获得，仍居留于我们心中的激情情绪同样可以在无生命之物和琐碎对象上发泄，不允许我们获得持续的平安状态，或摆脱我们残留的过失——除非我们或许认为，从无生命且无言之物无法回应我们的咒骂和愤怒，也不能将我们的失控脾气激化成更严重的疯狂激情这一事实中，可以为我们的激情带来某种益处和治愈。\par

% structure: p0059 source=h2 target=subsection reason=relative-heading-rank; base=h2

\subsection{第二十章}

\Emph{论按福音如何驱除愤怒。}\par

因此，若我们渴望获得那神圣赏报的本质——经上说：\BibleQuote{心里洁净的人是有福的，因为他们要看见天主。}\BibleRef{玛 5:8}——我们不仅应从行为上驱除愤怒，而应从内心深处完全根除它。因为仅在言语上抑制愤怒、不在行动上表现愤怒，是无益的，如果天主——人心中的隐秘瞒不过祂——看见愤怒仍存留于我们胸中的隐秘角落。因为福音的话命令我们根除过失的根源，而非果实；当所有刺激被移除后，这些根源肯定不会再发芽；当这愤怒不仅从行为和外在表现上被移除，而且从最内在的思想中被移除时，心灵就能在一切忍耐和圣洁中持续。因此，为避免杀人罪，愤怒和仇恨被切断，没有它们杀人罪不可能犯下。因为\Quote{凡向弟兄发怒的，就要受审判；}以及\BibleQuote{凡恨自己弟兄的，便是杀人犯；}\BibleRef{若一 3:15}即，因为他在心中想要杀他——我们确知他并未亲手或持武器在人前流他的血；然而，因他发怒，天主宣告他是杀人犯；天主赐予每个人赏报或惩罚，不仅按行为的结果，也按意图、欲望和意愿；正如祂自己藉先知所说：\BibleQuote{我要来聚集万民及各种语言的人；}\BibleRef{依 66:18}又说：\BibleQuote{他们的思想互相控告，或也为他们辩护；在天主审判人隐秘事的日子。}\BibleRef{罗 2:15}\par

% structure: p0062 source=h2 target=subsection reason=relative-heading-rank; base=h2

\subsection{第二十一章}

\Emph{福音所记\Quote{凡向弟兄发怒的}等处，我们是否应加入\Quote{无故}一词？}\par

但是你应该知道，在许多抄本中出现的这句话：\BibleQuote{凡向弟兄发怒而无故的，应受审判}，其中\Quote{无故}一词是多余的，是那些认为不应禁止因正当原因发怒的人所添加的。因为当然，无论一个人多么不合理地激动，也不会说他无故发怒。因此，似乎是那些不明白圣经主旨的人添加了它，圣经的意图是完全消除动怒的冲动，不保留任何发怒的机会；免得当我们被命令有原因地发怒时，我们就有了无故发怒的机会。因为忍耐的终点与目的，不在于有正当理由发怒，而在于根本不动怒。虽然我知道有些人认为\Quote{无故}这个说法本身是指一个人发怒时不被允许寻求报复。但更好的理解是，正如我们在许多现代抄本和所有古代抄本中所见的那样。\par

% structure: p0065 source=h2 target=subsection reason=relative-heading-rank; base=h2

\subsection{第二十二章}

\Emph{我们能从心中根除愤怒的补救方法。}\par

\Emph{因此}，基督的斗士，凡按规矩竞赛的，应当彻底根除愤怒的情感。对于这种疾病，确切的补救方法是：首先，我们下定决心绝不动怒，无论理由好坏；因为一旦我们心灵的主要光明被愤怒的阴影所蒙蔽，我们就会立即失去明辨之光、稳当的忠告、我们自己的正直以及正义的节制特性。其次，我们的灵魂纯洁将立即被蒙蔽，当愤怒之灵栖居在我们内时，它不可能成为圣神的殿宇。最后，我们应当考虑，我们绝不可在愤怒时祈祷，也不可向天主倾吐我们的祈祷。最重要的是，我们要常将人类的不确定状态摆在眼前，每日意识到我们即将离开肉身，而我们的节制与贞洁、舍弃一切所有、轻视财富、斋戒与守夜的辛劳，将毫无益处——如果仅因愤怒与仇恨，世界的审判者就判给我们永罚。\par

\SourceRecord{Translated by C.S. Gibson.；Nicene and Post-Nicene Fathers, Second Series；Vol. 11.；Edited by Philip Schaff and Henry Wace.；Buffalo, NY: Christian Literature Publishing Co.,；1894.；<http://www.newadvent.org/fathers/350708.htm>.}

% structure: p0001 source=h1 target=section reason=page-title

\section{\WorkTitle{教程（第九卷）}}

% structure: p0002 source=h2 target=subsection reason=relative-heading-rank; base=h2

\subsection{第一章}

\Emph{我们第五场战斗是如何对抗沮丧之灵的，以及它对灵魂造成的伤害。}\par

在第五场战斗中，我们必须抵抗啃噬人心的沮丧之苦：因为如果它通过随意分散的攻击和偶然的变化，获得了占据我们心灵的机会，它就会时时阻止我们洞悉神圣的默观，并使那已从完全纯洁状态失落的心灵彻底毁灭和消沉。它不允许心灵以惯常的喜乐祈祷，也不许它倚靠阅读圣经的安慰，更不许它对弟兄们安静温和；它使心灵在一切工作和敬拜的责任中变得急躁和粗暴：而且，当一切有益的建议丧失，心灵的坚毅被摧毁时，它使情感几乎疯狂和沉醉，并用惩罚性的绝望将其压碎和淹没。\par

% structure: p0005 source=h2 target=subsection reason=relative-heading-rank; base=h2

\subsection{第二章}

\Emph{应当如何小心医治沮丧之病。}\par

\Emph{因此}，如果我们渴望在属灵战斗的奋斗中合法地努力，就应当以同样谨慎的态度来医治这种病。因为\Quote{就如蛾子损坏衣服，虫子蛀坏木头，沮丧也如此腐蚀人的心灵。}圣神已经足够清晰和恰当地表达了这种危险而最有害的毛病的威力。\par

% structure: p0008 source=h2 target=subsection reason=relative-heading-rank; base=h2

\subsection{第三章}

\Emph{遭受沮丧攻击的灵魂可比作什么。}\par

因为被虫蛀的衣服不再有任何商业价值或良好用途；同样，被虫蛀的木头也不再值得用来装饰普通的建筑，只配丢在火里焚烧。因此，遭受啃噬性的沮丧攻击的灵魂，对那司铎的衣服也是无用的——按照圣达味的预言，从天而降的圣神的油膏首先降到亚郎的胡须上，再降到他的衣边上，然后才穿上这衣服。正如经上说：\BibleQuote{这好比珍贵的油倒在亚郎的头上，流到胡须上，又流到他的衣领上。}这灵魂也与那座属灵圣殿的建筑或装饰无关，那圣殿是保禄作为明智的建筑师奠立根基的，他说：\BibleQuote{你们是天主的圣殿，天主圣神住在你们内。}新娘在\BibleBook{}{雅歌}中告诉我们这圣殿的栋梁是怎样的：\BibleQuote{我们的栋梁是香柏木，我们的屋椽是扁柏木。}因此，为天主的圣殿所挑选的木材，是那些芬芳、不易腐烂、不受岁月朽坏、也不被虫蛀的。\par

% structure: p0011 source=h2 target=subsection reason=relative-heading-rank; base=h2

\subsection{第四章}

\Emph{沮丧从何而来，以何种方式产生。}\par

但有时我们发现它是由先前的愤怒之过所致，或出自某种未能实现的利益之欲，当一个人发现他未能获得所计划的东西时，便产生了沮丧。但有时，并无明显的原因驱使我们陷入这种不幸，我们却在狡猾仇敌的煽动下突然陷入如此深沉的阴郁，以至于我们不能以普通的礼貌接待亲近我们的人的来访；他们提出的任何话题，我们都觉得不合时宜、不合场合；我们无法给他们礼貌的回答，因为苦胆占据了我们内心的每一个角落。\par

% structure: p0014 source=h2 target=subsection reason=relative-heading-rank; base=h2

\subsection{第五章}

\Emph{扰乱不是由别人的过错，而是由我们自己的过错引起的。}\par

由此清楚证明，扰乱的痛苦并不总是由别人的过错在我们内引起，而更是由我们自己的过错，因为我们自己储存了得罪的原因和过错的种子，一旦诱惑之雨浇灌我们的灵魂，它们就会立即发芽、结果。\par

% structure: p0017 source=h2 target=subsection reason=relative-heading-rank; base=h2

\subsection{第六章}

\Emph{没有人是因突然跌倒而遭殃，而是因长期不慎而毁灭。}\par

因为除非他心里储存着邪恶的燃料，否则没有人会因别人的冒犯而被驱动犯罪。我们也不应想象一个人在看见女人时突然受骗，落入可耻情欲的深渊；而应认为，由于看见她的机会，那些隐藏在他灵魂深处的疾病症状暴露了出来。\par

% structure: p0020 source=h2 target=subsection reason=relative-heading-rank; base=h2

\subsection{第七章}

\Emph{我们不应为了寻求成全而放弃与弟兄的来往，而应不断修习忍耐之德。}\par

因此，万有的创造者天主，首先关注祂自己作品的改善，并且因为我们跌倒的根源和原因不在别人，而在我们自己，祂命令我们不应放弃与弟兄的来往，也不应躲避那些我们认为被我们伤害过或冒犯过我们的人，而应劝和他们，深知心灵的成全不是靠与人隔绝，而是靠忍耐之德获得的。当这德性被稳固持有时，它能使我们即使与那些仇恨和平的人也和好；但如果它没有被获得，它会使我们与那些比我们更成全和更好的人永远不合。因为只要我们生活在人群中，导致扰乱的机缘就不会缺乏——正是为了逃避这些，我们才急于离开与我们有联系的人；因此，我们无法完全逃脱，而只是改变了我们与先前朋友分离的沮丧原因。\par

% structure: p0023 source=h2 target=subsection reason=relative-heading-rank; base=h2

\subsection{第八章}

\Emph{如果我们改善了品行，就有可能和所有人相处。}\par

我们随后必须尽心竭力，设法改正我们的过错，修正我们的品行。如果我们成功改正了它们，我们必将获得平安——我不说与人，甚至与走兽和一切受造之物，正如\BibleBook{约伯}{传}中所说的：\BibleQuote{田野的走兽也必与你平安相处；} \BibleRef{约 5:23} 因为我们若不从内心接受并植入这些跌倒的根苗，就不会惧怕来自外界的冒犯，也不会有任何来自外部的跌倒机会搅扰我们：因为\BibleQuote{爱慕你法律的，必得大平安；他们绝不会跌倒。}\par

% structure: p0026 source=h2 target=subsection reason=relative-heading-rank; base=h2

\subsection{第九章。}

\Emph{论另一种产生救恩绝望的沮丧。}\par

还有一种更令人反感的沮丧，它使有罪的灵魂不能改善生活或纠正过错，反而产生最具毁灭性的绝望：这种沮丧没有使加音在杀害兄弟后悔改，也没有使犹达斯在背叛后急忙以补赎来解救自己，反而驱使他绝望地上吊自尽。\par

% structure: p0029 source=h2 target=subsection reason=relative-heading-rank; base=h2

\subsection{第十章。}

\Emph{论沮丧对我们唯一有用之处。}\par

因此，我们必须明白，沮丧只有在一种情况下对我们有益，即当我们因悔罪，或因被成全的渴望所激发，或因默观未来的真福而向它屈服时。关于这一点，圣宗徒说：\BibleQuote{按照天主而来的忧苦，能产生稳固的悔改，以致于救恩；但世俗的忧苦却产生死亡。}\BibleRef{格后 7:10}\par

% structure: p0032 source=h2 target=subsection reason=relative-heading-rank; base=h2

\subsection{第十一章。}

\Emph{我们如何分辨什么是有用的、按照天主的忧苦，什么是有魔性的、致死的。}\par

但是，那\BibleQuote{产生稳固悔改以致于救恩}的沮丧和忧苦是顺从的、温文的、谦卑的、和蔼的、温柔的、耐心的，因为它来自天主的爱，并且从成全的渴望不断扩展到每一种身体的痛苦和精神的忧伤；它不知怎地因希望自己的益处而欢喜并以此为食，保存了文雅和宽容的一切温和，因为它本身就拥有圣神的一切果实，同一宗徒列举了它们：\BibleQuote{圣神的果实是爱、喜乐、平安、忍耐、良善、温和、忠信、柔和、节制。}\BibleRef{迦 5:22} 但另一种则是粗暴的、不耐烦的、冷酷的、充满怨恨和无用的忧伤以及惩罚性的绝望，它摧垮它所依附的人，阻止他有精力和健康的忧伤，因为它是非理性的，不仅妨碍祈祷的效能，而且实际摧毁了我们所说的圣神的所有果实，那是另一种忧苦知道如何产生的。\par

% structure: p0035 source=h2 target=subsection reason=relative-heading-rank; base=h2

\subsection{第十二章。}

\Emph{除那以三种方式出现的健康忧苦外，所有忧伤和沮丧都应作为有害的加以抵挡。}\par

\Emph{为此}，除了那为悔罪、为成全或为未来渴望而忍受的忧苦外，所有忧伤和沮丧都必须同等抵挡，因为它们属于这个世界，是那\BibleQuote{产生死亡}的，必须像对淫乱、贪婪和愤怒的精神一样，完全从我们心中驱逐出去。\par

% structure: p0038 source=h2 target=subsection reason=relative-heading-rank; base=h2

\subsection{第十三章。}

\Emph{我们可以将沮丧从心中根除的方法。}\par

这样，我们应当能够将这最有害的情欲从心中驱逐出去，好借着属灵的默想，使我们的心思常专注于未来的希望和所应许幸福的默观。因为如此，我们就能胜过所有那些种类的沮丧——无论是源于先前的愤怒，还是源于得利的失望或某种损失，或是源于所受的冤屈，或是源于不合理的心理扰乱，或是那些给我们带来致命绝望的——如果我们常因洞悉永恒和未来之事而喜乐，并坚定不移，不为眼前事故所沮丧，也不因顺遂而过度得意，而是将每种境遇看作不确定且不久即逝的。\par

\SourceRecord{Translated by C.S. Gibson.；Nicene and Post-Nicene Fathers, Second Series；Vol. 11.；Edited by Philip Schaff and Henry Wace.；Buffalo, NY: Christian Literature Publishing Co.,；1894.；<http://www.newadvent.org/fathers/350709.htm>.}

% structure: p0001 source=h1 target=section reason=page-title

\section{\WorkTitle{训诫集（第十卷）}}

% structure: p0002 source=h2 target=subsection reason=relative-heading-rank; base=h2

\subsection{第一章}

\Emph{论我们第六场战斗如何对抗惰怠之灵，及其特性。}\par

我们第六场战斗是对抗希腊人所谓的\OriginalTerm{倦怠}{\GreekText{ἀκηδία}}，我们可以称之为心灵困乏或沮丧。这与沮丧类似，尤其困扰独居者，是沙漠隐修士危险而常见的敌人；它特别在第六时辰左右扰乱隐修士，如同一场定时发作的热病，在惯常的时刻袭击病人。最后，有些长老声称这就是第九十篇中提到的\Quote{午间恶魔}。\par

% structure: p0005 source=h2 target=subsection reason=relative-heading-rank; base=h2

\subsection{第二章}

\Emph{论惰怠的描述，以及它如何潜入隐修士的心，并对灵魂造成伤害。}\par

当这惰怠占据某个不幸的灵魂时，它便产生对所在地的厌恶、对隐修院小室的憎恶，以及对同住或稍远弟兄的轻蔑和鄙视，仿佛他们是疏忽大意或不热心的人。它也使人在围墙内必须做的各种工作中变得懒惰和迟钝。它不容许人留在小室里，也不容许人费心阅读；他常常叹息，因为留在那里一无所成；他抱怨并叹息，只要还与那个团体联结，就不能结出属灵的果实；他抱怨自己与属灵的收益隔绝，在所在之处毫无用处，仿佛他是一个本来可以管理他人、造福众人的人，却未能用其教导和道理来建造或益任何人。他赞扬远处的隐修院和那些遥远的地方，称那些地方更有益处、更适合获得救恩；此外，他描绘与那里弟兄的交往是甜蜜而充满属灵生命的。另一方面，他说身边一切都是粗糙的，不仅那里的弟兄毫无造就，而且连身体所需的食物也难以获得。最后，他幻想除非离开他的小室（他确信若再留下去必死无疑），并尽快离开此地，否则他永远不会健康。到了第五或第六时辰，他感到身体如此疲乏和渴望食物，以至于他觉得自己仿佛刚完成了一次长途旅行或非常繁重的工作，又好像已经两三天没有进食。此外，他焦急地东张西望，叹息没有弟兄来看望他；他常常进出小室，频频仰望太阳，仿佛它落得太慢；于是某种不合理的混乱思绪如幽冥的黑暗般占据了他，使他对于一切属灵工作都变得懒散无用，以致他认为除了拜访某位弟兄或睡眠的安慰之外，再也找不到任何方法对付如此可怕的攻击。然后，这疾病暗示他应当向弟兄们表示殷勤友好的款待，探望远近的病人。他还谈论一些孝爱和宗教的义务：应当询问亲戚的情况，更频繁地去看望他们；时常去看望那位虔诚的、献身于天主的妇人，她已失去所有亲人的支持，这会是真正的虔诚工作；最棒的事情是去获取那位被自己亲人忽视和鄙视的妇人所需要的东西；他应当虔诚地把时间花在这些事上，而不是无用且无益地待在小室里。\par

% structure: p0008 source=h2 target=subsection reason=relative-heading-rank; base=h2

\subsection{第三章}

\Emph{论惰怠克服隐修士的不同方式。}\par

因此，受困于敌人如此诡计的可怜灵魂便不安起来，直到被惰怠之灵如同沉重的破城锤般耗尽心力，它要么陷入沉睡，要么被逐出隐修院的围墙，习惯于在这些攻击下寻求安慰，去拜访某位弟兄，却只因为这暂时的补救而变得更加软弱。因为敌人会越来越频繁、越来越猛烈地攻击那刚一交战就转身逃跑的人，那人既不靠胜利也不靠战斗，而只靠逃跑来寻求安全；直到他渐渐被引出隐修院，开始忘记自己修会的目的——那无非是默想和瞻仰超越万有的神圣纯洁，而这只能通过静默、恒常留在隐修院和默想来获得——于是基督的士兵便从他服役中逃跑，成为逃兵，并\Quote{卷入世俗事务}，全然不讨他所投效之主的喜悦。\BibleRef{弟后 2:4}\par

% structure: p0011 source=h2 target=subsection reason=relative-heading-rank; base=h2

\subsection{第四章}

\Emph{论惰怠如何阻碍心灵对一切德行的默想。}\par

大卫在一句诗中精彩地表达了这种疾病的所有不适，他说：\Quote{我的灵魂因困乏而沉睡}，即因惰怠而沉睡。他说灵魂沉睡而非身体沉睡，是极为恰当的。因为确实，被这情欲之箭所伤的灵魂，在一切德行的默想和属灵感官的洞察方面都睡着了。\par

% structure: p0014 source=h2 target=subsection reason=relative-heading-rank; base=h2

\subsection{第五章}

\Emph{论惰怠的攻击是双重的。}\par

因此，真正的基督徒运动员，若想在完美的赛场上合法地竞赛，就当赶紧将这疾病也从灵魂深处驱逐出去；他应当从两方面对抗这最邪恶的惰怠之灵，使自己既不被睡眠之箭射倒，也不被赶出隐修院的围栏，哪怕是以某种虔诚的借口或托词，作为逃兵离去。\par

% structure: p0017 source=h2 target=subsection reason=relative-heading-rank; base=h2

\subsection{第六章}

\Emph{论惰怠的危害何等严重。}\par

每当它开始多少战胜某人时，它或使他留在他的小室里闲散怠惰，不作任何神修进步，或由那里将他驱赶出去，使他不得安宁、到处游荡，并且在各项工作上懈怠，又使他不断绕行弟兄们的小室和隐修院，眼目所图不外乎一事，即：他现在在哪里或藉什么借口可以立刻获得一些食物。因为闲散者的心中所思无非是食物和肚腹，直至他寻得某个同样冷淡漠然的男人或女人的陪伴，便沉溺于他们的琐事和俗务中，就这样渐渐被危险的事务缠住，以致好像被蛇的盘绕束缚，再也不能解脱自己，回到起初修道身份的全德。\par

% structure: p0020 source=h2 target=subsection reason=relative-heading-rank; base=h2

\subsection{第七章}

\Emph{关于\OriginalTerm{懈怠}{accidie}之灵来自宗徒的见证。}\par

蒙福的宗徒，如同一位真正的属灵医生，或已看见这由懈怠精神而生的疾病正在蔓延，或藉圣神的启示预见到它将在隐修士中出现，便迅速用他指示的医治之药加以预防。因为当他写信给得撒洛尼人时，起初像一位巧妙卓越的医生，用他安抚和温和的话语之药治疗患者的软弱，并以爱德开始，在那方面称赞他们，好使这致命的创伤经较温和的治疗后，失去其愤怒的溃烂，并能更容易承受更严厉的治疗。他说：\BibleQuote{至于兄弟之爱德，你们无需我写信给你们，因为你们自己蒙天主教导要彼此相爱。你们对全马其顿的众弟兄也确实这样行了。}\BibleRef{得前 4:9} 他首先以称赞的安抚开始，使他们的耳朵顺从并准备好接受医治话语的药方。然后他继续：\Quote{但我们求你们，弟兄们，更该充盈。} 至此他用亲切温和的话语安抚他们；惟恐发现他们尚未准备好接受完全的痊愈。为什么，宗徒啊，你请求他们更该充盈于爱德，而你以上却说\Quote{至于兄弟之爱德，我们无需写信给你们}？为什么你必须对他们说：\Quote{但我们求你们更加充盈}，而他们在此事上根本无需被写信？尤其你添加了他们无需的理由，说：\Quote{因为你们自己蒙天主教导要彼此相爱。} 你又添加了第三件更重要的：他们不仅蒙天主教导，而且实际上履行了所学的。他说：\Quote{因为你们这样行了}，不是向一两人，而是\Quote{向众弟兄}；不仅向自己的同胞和朋友，而是\Quote{在全马其顿}。那么，请告诉我们，为什么你特别以此开始。他又继续：\Quote{但我们求你们，弟兄们，更该充盈。} 最后他艰难地吐露出他之前所推动的：\Quote{并努力安分守己。} 他给出了第一个目标。然后他添加第二个：\Quote{并做自己的事；} 还有第三个：\Quote{用你们自己的手工作，如同我们吩咐过你们的；} 第四个：\Quote{在外人面前品行端正；} 第五个：\Quote{不贪图别人的财物。} 看，我们可以透过那犹豫看出，使他用这些序言推迟说出心中所充满的：\Quote{并努力安分守己；} 即，你们应留在自己的小室中，不被谣言扰乱，这些谣言通常出自游手好闲者的愿望和闲谈，而你们自己也因此扰乱他人。还有，\Quote{做自己的事}，你们不应好奇地追问世间的事，或审视他人的生活，而将精力花在不是改善自己和追求德行，而是贬低你们的弟兄上。\Quote{并用你们自己的手工作，如同我们吩咐过你们的；} 为确保他以上警告你们不可做的事，即他们不应为别人的事烦躁焦虑，也不应在外人面前行事不端，也不贪图别人的财物，他现在添加说：\Quote{并用你们自己的手工作，如同我们吩咐过你们的。} 因为他已清楚表明，闲散是上述他所责备之事的缘由。因为没有人会为别人的事烦躁焦虑，除非他不满足于专心从事自己手中的工作。他又添加了第四个恶，也出自这闲散，即他们不应行事不端：他说：\Quote{并在外人面前品行端正。} 一个不满足于坚守小室的隐修生活和自己手中工作的人，甚至在这世上的人中间也不可能品行端正；他必定不诚实，因为他寻求必要的食物；并努力谄媚，追随新闻和流言，寻找闲谈和故事的机会，以便能立稳脚跟并进入他人的房屋。\Quote{并且你们不应贪图别人的财物。} 一个不关心用自己双手的忠诚和平的劳动来获取每日食物的人，必定会用嫉妒的眼光看待别人的恩赐和恩惠。你们看，何等严重和可耻的状况，唯独出自闲散之病。最后，那些人在他的第一封书信中他曾用话语的温柔治疗对待，在第二封书信中，他试图以更严厉更苛刻的医治来治疗他们，作为未从更温和的治疗中获益的人；他不再使用温和话语的治疗，没有温和仁慈的表达，如\Quote{但我们求你们，弟兄们}，而是\BibleQuote{我们在主耶稣基督的名下吩咐你们，弟兄们，要远离每个行为不检的弟兄。}\BibleRef{得后 3:6} 那里是请求；这里是命令。那里是劝诱者的仁慈；这里是抗议和威胁者的严厉。\Quote{我们在主耶稣基督的名下吩咐你们，弟兄们：} 因为当我们初次请求你们时，你们不屑听从；如今至少服从我们的威胁。并且他使这命令变得可怕，不是单靠话语，而是藉着呼求主耶稣基督的名：惟恐他们再次轻视，以为只是人的话，认为无关紧要。随即，如同一位技术精湛的医生，面对溃烂的肢体，无法施用温和治疗，便试图用属灵的刀切开医治，说：\Quote{要远离每个行为不检、不按我们从你们领受的传统行事的弟兄。} 因此，他命令他们远离那些不肯花时间工作的人，将他们像被闲散溃烂疮疖污染的肢体一样切除：以免懈怠之病像某种致命传染病，通过感染的逐步蔓延，甚至感染他们肢体的健康部分。当他将要谈到那些不肯用自己双手工作、安然吃自己面包的人，并催促他们远离时，请听他以何种谴责一开始便给他们打上烙印。首先他称他们为\Quote{行为不检} 和\Quote{不按传统行事}。换言之，他谴责他们顽固，因为他们不肯按他的规定行事；以及\Quote{不端正}，即不遵守外出、拜访和交谈的合宜时间。因为一个行为不检的人必定会有所有这些缺点。\Quote{并且不按从我们领受的传统。} 由此他给他们打上某种反叛和藐视者的烙印，他们不屑遵守从他领受的传统，不肯追随他们不仅记得师傅在言语上教导过、而且知道他在行为上实行过的事。\Quote{因为你们自己知道你们应如何效法我们。} 他堆积了巨大的指责，当他断言他们没有遵守那仍留在他们记忆中的事，那不仅藉口头教训学习过，而且藉他在工作中的榜样激励而领受的。\par

% structure: p0023 source=h2 target=subsection reason=relative-heading-rank; base=h2

\subsection{第八章。}

\Emph{凡不以亲自劳动为满足者，必不得安宁。}\par

\Quote{因为我们在你们中间并不闲散。}当他想要通过劳动的事实来证明自己并不闲散时，他充分显示了那些不肯劳动的人由于懒惰的过失而常常闲散。\Quote{我们也未曾白吃过任何人的食物。}外邦人的导师用每句话来加强责备。福音的宣讲者说他未曾白吃任何人的食物，因为他知道主曾吩咐说：\Quote{传福音的人应靠福音而生活}\BibleRef{格前 9:14}；又说：\Quote{工人应得他的食物}\BibleRef{玛 10:10}。所以，如果那位宣讲福音、执行如此崇高和属灵工作的人，不敢依靠主的命令白吃食物，那么我们这些不仅没有被托付宣讲圣言、甚至连自己的灵魂也得不着的人，将如何呢？当那位\Quote{蒙选的器皿}因对福音的焦虑和宣讲的工作而不得不亲手劳动时，我们有什么信心敢闲散地白吃食物呢？他说：\Quote{而是带着辛劳和困倦，昼夜工作，免得加重你们任何人的负担}\BibleRef{得后 3:8}。至此，他扩展并加深了他的责备。因为他没有简单地说：\Quote{我们从你们任何人那里未曾白吃过食物}，就停止了。因为有人可能以为他是靠自己的私产、积蓄的金钱或他人的资助（虽然不是靠他们的捐输和礼物）来维持的。然而他说：\Quote{而是带着辛劳和困倦，昼夜工作}，也就是说，特别靠自己的劳动来维持。他还说，这样做并非出于自己的意愿或喜好，如休息和身体锻炼所建议的，而是由于需要和食物的缺乏所迫，并且不是没有巨大的身体困倦。因为不仅在全天，而且在夜晚（看来是为了身体休息而赐予的），我也为了食物而不断从事手工劳动。\par

% structure: p0026 source=h2 target=subsection reason=relative-heading-rank; base=h2

\subsection{第九章。}

\Emph{不仅宗徒，而且与他同在的那两个人也亲手劳动。}\par

他证明不仅他独自这样生活在他们中间，免得这种方法只依靠他的榜样就显得不重要或普遍。但他宣告所有与他一同被派去传福音的人，即与他一同写这封信的\Person{西拉}和\Person{弟茂德}，也同样工作。因为他说\Quote{免得我们加重你们任何人的负担}，使他们深感羞愧。因为如果那位宣讲福音并用奇迹和异能来证明福音的人，尚且不敢白吃食物以免加重他人负担，那么那些每日闲散无事的人如何能不认为自己是在加重他人负担呢？\par

% structure: p0029 source=h2 target=subsection reason=relative-heading-rank; base=h2

\subsection{第十章。}

\Emph{正因为如此，宗徒亲手劳动，为给我们留下劳动的榜样。}\par

\Quote{并非我们没有权力，而是为给你们留下榜样，好使你们效法我们。}他揭露了为何将如此劳苦加诸自己：他说：\Quote{为给你们留下榜样，好使你们效法我们}，这样，如果你们偶然忘记了我们经常从你们耳边经过的话语教训，至少你们可以透过亲眼看见的榜样，记住我的生活方式。这里也有对他们的轻微责备，他说他昼夜从事这样劳苦，没有别的原因，只是为树立榜样，然而他们仍然不受教导，而为了他们的缘故，他虽没有义务，却仍将这样的辛劳加诸己身。他说：\Quote{事实上，我们虽有权力，也有机会使用你们所有的财物和产业，而且我知道主允许我使用它们：但我并没有使用这权力，恐怕我的合法行为会为他人树立危险懒惰的榜样。因此，当我宣讲福音时，我更愿意靠自己的双手和劳动来维持，好为你们这些愿意行走在德性道路上的人开辟成全之路，并借我的劳动树立良好生活的榜样。}\par

% structure: p0032 source=h2 target=subsection reason=relative-heading-rank; base=h2

\subsection{第十一章。}

\Emph{他不仅用榜样，也用言语宣讲并教导人劳动。}\par

他接着说：\Quote{因为当我们与你们同在时，就已向你们吩咐过：若有人不肯工作，就不该吃饭。}他使他们的懒惰更加明显：因为他们知道，他如同一位好老师，为了教导的缘故亲手劳动，为的是训导他们，但他们却羞于效法他；他还强调我们的勤勉和关心，说他不仅亲在时给了他们这个榜样，而且也持续用言语宣讲，说若有人不肯工作，就不该吃饭。\par

% structure: p0035 source=h2 target=subsection reason=relative-heading-rank; base=h2

\subsection{第十二章。}

\Emph{论他所说的：\Quote{若有人不肯工作，就不该吃饭。}}\par

如今他不再以导师或医师的身份向他们提出劝告，而是以司法判决的严厉态度行事，恢复他的宗徒权威，如同从审判席上对轻蔑者宣判：我指的是那种权威——当他以威胁写信给格林多人时，他宣称这权威是主所赐的，当时他命令那些陷于罪恶中的人，要在他到来之前赶紧改正他们的生活；他这样命令说：\Quote{我恳求你们，免得我当面对某些人，使用那赐给我的权柄来对付你们。} 又说：\Quote{即使我稍微夸耀那主赐给我、为成就而非为摧残你们的权柄，我也不至于羞愧。} 我指的是，他宣称：\Quote{若有人不愿工作，就不该吃饭。} 他不是用血肉的剑惩罚他们，而是藉圣神的能力禁止他们享受今生的财物，使那些或许因很少想到未来死亡的惩罚、因贪图安逸而仍顽固的人，最终被自然的需要和对即刻死亡的恐惧所驱使，被迫服从他这有益的诫命。\par

% structure: p0038 source=h2 target=subsection reason=relative-heading-rank; base=h2

\subsection{第十三章}

\Emph{关于他所说的话：\Quote{我们听说你们中间有人行事不规矩。}}\par

然后，在这一切福音严厉的严苛之后，他现在揭示他提出所有这些事情的原因。\Quote{因为我们听说你们中间有人行事不规矩，什么工也不做，反倒好管闲事。} 他绝不满足于把那些不肯投身工作的人说成仿佛只患了一种病。因为在他的第一封书信中，他称他们为\Quote{不规矩的}，并说他们不按从他那里领受的传统行事；他还断言他们不安分，白吃别人的饭。他在这里又说：\Quote{我们听说你们中间有人行事不规矩。} 随即他又附上第二种软弱，这是不安的根源，说：\Quote{什么工也不做；} 他还加上第三种毛病，它像嫩枝一样从最后一种生发出来：\Quote{反倒好管闲事。}\par

% structure: p0041 source=h2 target=subsection reason=relative-heading-rank; base=h2

\subsection{第十四章}

\Emph{体力劳动如何防止许多过犯。}\par

因此，他毫不拖延地立刻将适当的治疗应用于这么多过犯的诱因，放下他刚才使用过的宗徒权威，再次采纳一位好父亲或仁慈医师的温柔性格，如同对待他的孩子或病人一样，用他治愈性的劝告施以疗法来治愈他们，说：\Quote{现在我们吩咐这样的人，并藉主耶稣恳求他们，安静作工，吃自己的饭。} 所有这些溃疡的起因，都是出于懒惰的根源，他像一位技艺精湛的医师一样，藉一条有益的工作命令来医治；因为他知道，所有其他坏症状——它们如同从同一丛生发出来——一旦主要病症的起因被移除，就会立刻消失。\par

% structure: p0044 source=h2 target=subsection reason=relative-heading-rank; base=h2

\subsection{第十五章}

\Emph{如何甚至向懒惰和粗心的人显示仁慈。}\par

\Emph{然而}，像一位有远见且谨慎的医师，他不仅急于医治病人的创伤，也给予全体适当的指导，使他们的健康得以持续保持，并说：\Quote{但你们行善不可厌倦。} 你们这些跟随我们，即我们的道路的人，藉着在劳动上效法我们而复制了给予你们的榜样，不要跟随他们的怠惰和懒散：\Quote{行善不可厌倦；} 也就是说，你们也要向他们显示仁慈，如果碰巧他们未能遵守我们所说的。正如他对那些软弱的人严厉，恐怕他们因懒惰而变得软弱，屈服于不安分和好奇，同样他劝诫那些健康的人，既不要限制那主的诫命吩咐我们对善人和恶人都要显示的爱德（\BibleRef{玛 5:43}），即使有些恶人不会转向健康的教义；也不要停止行善，并用安慰的话和责备，以及普通的仁慈和礼貌来鼓励他们。\par

% structure: p0047 source=h2 target=subsection reason=relative-heading-rank; base=h2

\subsection{第十六章}

\Emph{我们应如何劝戒那些走错路的人：不是出于仇恨，而是出于爱。}\par

但是，唯恐有些人可能因这种温和而受到鼓励，并藐视遵守他的命令，他便以宗徒的严厉态度继续：\Quote{若有人不听从我们这书信上的话，要记下那人，不要和他交往，叫他自觉羞愧。} 在警告他们应当出于对他的尊重和对众人的益处而遵守什么，以及他们应如何谨慎地持守宗徒的命令时，他立刻将一位最慈祥父亲的仁慈与警告结合起来；并且教导他们，如同教导自己的孩子一样，他们应对上述那些人怀有怎样的兄弟情怀，出于爱。\Quote{然而不要以他为仇人，但要劝戒他如弟兄。} 他将审判官的严厉和父亲的慈爱结合起来，用仁慈和温和缓和了以宗徒严厉宣判的判决。因为他命令他们记下那藐视服从他命令的人，不要和他交往；但他并不是吩咐他们出于一种错误的厌恶感这样做，而是出于兄弟间的爱和为他们悔改的考虑。\Quote{不要与他交往，}他说，\Quote{使他自觉羞愧；}这样，即使我的温和警告没有使他变得更好，至少他可能因被公开与你们所有人隔离而感到羞愧，从而有一天开始回归得救之路。\par

% structure: p0050 source=h2 target=subsection reason=relative-heading-rank; base=h2

\subsection{第十七章}

\Emph{宗徒声明我们应当劳动的不同段落，或表明他亲自劳动的不同段落。}\par

在\BibleBook{厄弗所}{书}中，他也同样就工作一事如此训导说：\BibleQuote{那偷窃的，不要再偷窃，却更要劳苦，亲手做正当的事，好能周济有需要的人。} \BibleRef{弗 4:28} 在\BibleBook{宗徒}{大事录}中，我们也发现他不仅教导这事，而且自己实际实践了。因为他来到格林多后，除了阿桂拉和普黎史拉那里，他不住别处，因为他们与他有相同的技艺，那是他惯常做的。因为这样我们读到：\BibleQuote{此后，保禄离开雅典，来到格林多。遇见一个犹太人，名叫阿桂拉，生在本都，同他的妻子普黎史拉，就投奔他们，因为他们与他同技艺，便住在他们那里工作。他们原是制帐棚为业的。} \BibleRef{宗 18:1}\par

% structure: p0053 source=h2 target=subsection reason=relative-heading-rank; base=h2

\subsection{第十八章。}

\Emph{宗徒的劳作足以供给他自己及与他同在的人。}\par

然后他前往米肋托，从那里派人到厄弗所，召请厄弗所教会的长老来，嘱咐他们在他不在时应当如何管理天主的教会，并说：\BibleQuote{我从未贪图过任何人的金银或衣服。你们自己知道：这双手供应了我和同我一起的人的需要。在各方面我都给你们立了榜样，告诉你们必须这样劳动，扶助软弱者，并且要记住主耶稣的话，他说过：\Quote{施予比领受更为有福。}} \BibleRef{宗 20:33} 他用自己的生活方式给我们留下了重要的榜样，因为他证明自己不仅劳作供应自己肉身的需求，也足够供应那些与他在一起的人的需要；我是指那些被必要的职责所占据，没有机会亲手为自己挣取食物的人。并且，正如他告诉得撒洛尼人，他工作是为了给他们树立榜样，使他们可以效法他，同样，在这里他也暗示了同样的事，当他说：\Quote{在各方面我都给你们立了榜样，告诉你们必须这样劳动，扶助软弱者，} ——无论是在精神上还是身体上；也就是说，我们应该勤勉地供应他们的需要，不是从我们丰裕的储备或积攒的钱财中，也不是从他人的慷慨和财物中，而是通过我们自己的劳动和辛苦来赚取所需的金额。\par

% structure: p0056 source=h2 target=subsection reason=relative-heading-rank; base=h2

\subsection{第十九章。}

\Emph{如何理解这句话：\BibleQuote{施予比领受更为有福。}}\par

他说这是主的命令：\Quote{因为祂自己，}即主耶稣，他说道，\Quote{说道：施予比领受更为有福。} 也就是说，施予者的慷慨比接受者的需要更为有福，当这礼物不是来自因不信或失信而扣留的钱财，也不是来自贪吝所积累的宝藏，而是来自我们自己的劳动和诚实辛劳的果实之时。所以\Quote{施予比领受更为有福，}因为施予者虽然分担接受者的贫穷，但他仍然通过自己的辛劳以虔诚的关怀勤勉地提供，不仅足够自己的需要，也足够给予有需要的人；因此他配得到双重的恩宠，因为通过放弃他所有的一切，他确保了基督完全的舍弃，同时通过他的劳动和心思展示了富人的慷慨；这样，他以诚实的劳动光荣天主，并为祂采摘义德的果实，而另一个人，因懒惰和无所事事的怠惰而衰弱，借着宗徒的话证明自己配不上食物，因为他违背宗徒的命令，在闲散中领受，并非没有罪和顽固的罪疚。\par

% structure: p0059 source=h2 target=subsection reason=relative-heading-rank; base=h2

\subsection{第二十章。}

\Emph{论一位懒惰的兄弟试图说服他人离开隐修院。}\par

我们知道一位兄弟，如果说出他的姓名有好处的话，我们就会说出来。他虽然在隐修院中，并被迫每天向总管交出固定的工作成果，但因害怕自己可能被引向更大份额的工作，或因另一位更热心劳作的榜样而感到羞愧，当他看到某位兄弟被接纳进隐修院，凭着他信德的热忱想要完成更大规模的工作时，如果他发现不能通过私下劝说阻止他实现目的，就会用坏建议和耳语说服他离开那里。为了更容易摆脱他，他会假装自己也已经因多种原因感到不满，并想要离开，只要他能找到旅伴和路途支持。当他通过暗中贬损隐修院诱惑他同意，并与他约定离开隐修院的时间、他应该先去的地方以及在哪里等自己时，他自己假装随后就来，却停留在原地。当对方因羞愧于自己的逃跑而不敢再回到他所逃离的隐修院时，这可怜的逃跑制造者却留在隐修院中。举这一个这类人的例子就足以警告初学者，并清楚表明，如圣经所说，闲散能在隐修士心中产生何等邪恶，以及\BibleQuote{交恶友会败坏善行。} \BibleRef{格前 15:33}\par

% structure: p0062 source=h2 target=subsection reason=relative-heading-rank; base=h2

\subsection{第二十一章。}

\Emph{来自撒罗满著作中反对\OriginalTerm{懈怠}{accidie}的不同段落。}\par

最智慧的人撒罗满在多处清楚指出怠惰的过错，他说：\Quote{追求怠惰的，必饱尝贫困；}\BibleRef{箴 28:19}——即可见的或不可见的贫困，怠惰者和陷于各种过失的人必定被牵连其中，且将永远与天主的默观和属灵的财富无缘。蒙福的宗徒说：\Quote{因为你们在一切事上，在他内都成了富足的，包括各种言论和各种知识。}\BibleRef{格前 1:5}至于怠惰者的这种贫困，他在别处也这样写道：\Quote{每个懒惰人都将穿着破烂的衣裳和破布。}\BibleRef{箴 23:21}因为，凡被懒散或怠惰压倒的人，宁愿不凭自己的辛劳和勤勉，而用从圣经整块布匹上撕下的破布来遮蔽自己，为自己披上不是光荣和尊贵的衣服，而是可耻的遮盖和借口，他当然不配装饰那不朽的衣服（宗徒说：\Quote{穿上主耶稣基督；}\BibleRef{罗 13:14}又说：\Quote{穿上公义和爱德的胸甲；}\BibleRef{得前 5:8}关于这衣服，主自己也藉先知对耶路撒冷说：\Quote{兴起，兴起，耶路撒冷，穿上你光荣的衣服）}\BibleRef{依 52:1}因为那些受这种懒惰影响、不愿像宗徒常做并吩咐我们那样亲手劳作自养的人，惯于使用某些圣经证明，试图以此掩盖他们的怠惰，说经上记载：\Quote{不要为那可朽坏的食粮劳碌，而要为那存留到永生的食粮劳碌；}\BibleRef{若 6:27}以及\Quote{我的食物就是承行我父的旨意。}\BibleRef{若 4:34}但这些证明（就像破布）是从福音的整块布匹上撕下来的，其目的是掩盖我们怠惰和羞耻的耻辱，而不是为我们保暖，并用那智慧妇人（在箴言中，她穿着力量和美丽）为自己或丈夫所制的珍贵华服来装饰我们；随后说：\Quote{力量和美丽是她的衣裳，她为日后而喜乐。}\BibleRef{箴 31:25}关于怠惰之恶，撒罗满又这样提到：\Quote{懒惰人的路铺满了荆棘；}即这些以及类似的过失，宗徒以上曾说它们源于怠惰。又说：\Quote{每个懒惰人总是有所缺乏。}宗徒提到这些时说：\Quote{并且你们不要缺少任何人的东西。}\BibleRef{得前 4:11}最后：\Quote{因为怠惰已成了许多邪恶的老师；}\BibleRef{德 33:29}宗徒在他以上论述的段落中清楚地列举了这些：\Quote{什么也不做，反而多管闲事。}他又将另一项与这过错相连：\Quote{并且你们要安心从事；}然后，\Quote{你们应做自己的事，在外人面前要诚实行事，并且不要缺少任何人的东西。}他也藉着这些指责那些认真的人要远离那些不像样的人：\Quote{你们要远离，}他说，\Quote{每一个行事不像样、不按我们从他们那里所领受的传统而行的弟兄。}\par

% structure: p0065 source=h2 target=subsection reason=relative-heading-rank; base=h2

\subsection{第二十二章}

\Emph{埃及的弟兄们如何亲手劳作，不仅供应自己的需要，也供应那些在监狱中的人。}\par

受这些榜样的教导，埃及的教父们绝不允许隐修士、尤其是年轻的隐修士闲散，他们透过勤奋劳作来估量他们内心的意向以及在忍耐和谦逊上的长进；他们不仅不允许从他人那里接受任何东西来供应自己的需要，反而不仅藉着他们的劳动款待来访的朝圣者和弟兄们，而且还收集大量存粮和食物，将它们分发给利比亚遭受饥荒和贫瘠的地区，以及在城市中在监牢污秽中憔悴的人；因为他们相信，藉着这样奉献他们双手的果实，他们向主献上了合理而真实的祭献。\par

% structure: p0068 source=h2 target=subsection reason=relative-heading-rank; base=h2

\subsection{第二十三章}

\Emph{怠惰是西方没有隐修士隐修院的原因。}\par

因此，在这些地区，我们看不到有如此多弟兄的隐修院：因为他们不是靠自己的劳动资源来维持，以致能继续留在其中；即使通过他人的慷慨有时有足够的供给，但贪图安逸和内心的不安也不允许他们在那里久留。因此，在埃及的老教父们中流传着这样一句名言：工作的隐修士只受一个魔鬼攻击，而闲散的隐修士却被无数恶灵折磨。\par

% structure: p0071 source=h2 target=subsection reason=relative-heading-rank; base=h2

\subsection{第二十四章}

\Emph{保禄院长每年用火焚烧他亲手所做的一切工作。}\par

最后，教父中最伟大者之一的院长保禄，当他居住在被称为波尔菲里旷野的广袤荒漠中时，因椰枣树和一个小园子而无忧无虑，生活自足，食物充裕，他找不到任何其他可以养活自己的工作，因为他的住所距城镇和有人居住的地区有七天的路程甚至更远，穿越荒漠，且运输货物的费用会超过工作的价值；他便收集棕榈叶，并定期为自己规定每日的劳作，仿佛要靠此维生。当他的洞穴堆满了整年的工作时，他每年都会用火焚烧那些他如此勤勉劳作的成果：从而证明了没有手工劳作，隐修士无法停留在一个地方，也无法登上完美的顶峰；因此，尽管食物的需要并不要求这样做，但他仍仅仅为了净化心灵、坚固思想、坚持留在自己的小室中，并战胜和驱除\OriginalTerm{灵性怠惰}{accidie}而行此事。\par

% structure: p0074 source=h2 target=subsection reason=relative-heading-rank; base=h2

\subsection{第二十五章}

\Emph{院长梅瑟关于治愈\OriginalTerm{灵性怠惰}{accidie}而向我说的话。}\par

当我开始在荒漠中居住时，我曾对诸圣之首的院长梅瑟说，昨日我被一次\OriginalTerm{灵性怠惰}{accidie}的袭击所困扰，只有立刻跑到院长保禄那里才能摆脱它。他说：\Quote{你并没有使自己摆脱它，反而把自己交付给它，成为它的奴隶和臣属。因为敌人今后将更猛烈地攻击你，视你为逃兵和逃亡者，因为它已看到你在战斗中被击败便立即逃跑：除非下一次你与它交战时，决心不通过舍弃你的小室或通过睡眠的懈怠来驱散它的攻击和热度，而是学习通过忍耐和战斗来战胜它。}因此，经验证明，一次灵性怠惰的发作不应通过逃跑来躲避，而应通过抵抗来战胜。\par

\SourceRecord{Translated by C.S. Gibson.；Nicene and Post-Nicene Fathers, Second Series；Vol. 11.；Edited by Philip Schaff and Henry Wace.；Buffalo, NY: Christian Literature Publishing Co.,；1894.；<http://www.newadvent.org/fathers/350710.htm>.}

% structure: p0001 source=h1 target=section reason=page-title

\section{隐修规章（第十一卷）}

% structure: p0002 source=h2 target=subsection reason=relative-heading-rank; base=h2

\subsection{第一章}

\Emph{论我们第七场战斗是对抗虚荣之灵，以及它的本性。}\par

我们的第七场战斗是对抗\OriginalTerm{虚荣}{\GreekText{κενοδοξία}}之灵，我们可称之为虚空或徒劳的荣耀：这是一个多形、易变且诡诈的灵，以至于最敏锐的眼睛也难以——我不说提防——看透并发现它。\par

% structure: p0005 source=h2 target=subsection reason=relative-heading-rank; base=h2

\subsection{第二章}

\Emph{虚荣如何不仅在肉身方面，也在灵性方面攻击隐修士。}\par

因为此恶习不仅像其他过犯一样攻击隐修士的肉身方面，也攻击其灵性方面，它狡猾而诡诈地潜入心灵：以致那些不能被肉身之欲所欺骗的人，反而因灵性上的进步受到更严重的伤害；要对抗它更为糟糕，因为它更难于防范。因为其他一切恶习的攻击更为公开和直接，在每种情况下，当激起它们的人遇到坚决的拒绝时，他会因此变得软弱，而被击败的对手下次攻击受害人时也会减弱力量。但这病症一旦借着肉身的骄傲攻击了心灵，且被反驳的盾牌所击退，就又像某种多形的邪恶，改变其先前的伪装和特征，在德行的外表下试图击倒并毁灭它的征服者。\par

% structure: p0008 source=h2 target=subsection reason=relative-heading-rank; base=h2

\subsection{第三章}

\Emph{虚荣有多少种形式和形状。}\par

因为我们的其他过犯和情欲可以说较为简单，只有一种形式；但此恶习有多种形式和形状，并且改变自身，从四面攻击那反对它的人，从各方冲击它的征服者。因为它试图伤害基督的士兵——在他的穿着、举止、行走、声音、工作、守夜、禁食、祈祷、退隐、阅读、知识、沉默、服从、谦卑、忍耐中；而且像某些最危险的礁石被汹涌波涛隐藏，它给那些顺风顺水航行的人带来意想不到的悲惨沉船，而他们却没有警惕或防范它。\par

% structure: p0011 source=h2 target=subsection reason=relative-heading-rank; base=h2

\subsection{第四章}

\Emph{虚荣如何在右左两边攻击隐修士。}\par

凡愿意借着\Quote{左右两手的正义武器}行走在王道上的人，应当根据宗徒的教导，经历\Quote{荣誉和羞辱，恶名和美名，}\BibleRef{格后 6:7}，并以这样的谨慎，以明辨为舵，并借主的圣神在我们身上吹拂，在试探的汹涌波涛中引导他的德行之程，因为我们知道，如果我们稍向左或右偏离，将会立即撞上最危险的崖石。因此，最明智的撒罗满警告我们说：\Quote{不可偏向左或右；}即，不要因右手的德行而沾沾自喜，因灵性成就而骄傲；也不要偏向左手的恶习之路，从它们那里为自己寻求（用宗徒的话）\Quote{以你的羞耻为荣耀}\BibleRef{斐 3:19}。因为魔鬼若不能借着合身整洁的衣着在人身上制造虚荣，就试图用肮脏、廉价、不加修饰的样式引入虚荣。若不能以荣誉使人跌倒，就以谦卑推翻他。若不能以知识与口才的恩宠使他骄傲，就以沉默的重量拉倒他。若有人公开禁食，就被虚妄的骄傲攻击。若他为了鄙视虚荣而隐藏禁食，又被同样的骄傲之罪袭击。为了避免被虚荣的污点玷污，他避免在弟兄们面前作长祈祷；然而，因为他私下祈祷且无人知晓，他并没有逃脱虚妄的骄傲。\par

% structure: p0014 source=h2 target=subsection reason=relative-heading-rank; base=h2

\subsection{第五章}

\Emph{一个表明虚荣本性的比喻。}\par

我们的长老们巧妙地描述这疾病的本质如同洋葱和某些鳞茎，当你剥去一层表皮，却发现它被另一层包裹着；你剥多少次，就发现它仍然被保护着。\par

% structure: p0017 source=h2 target=subsection reason=relative-heading-rank; base=h2

\subsection{第六章}

\Emph{虚荣并不能完全借独处的益处被消除。}\par

在独处中，它也不停止追逐那为了荣耀而逃避与所有人交往的人。一个人越是彻底地避开整个世界，它就越热烈地追逐他。它试图使一个人因对工作和劳苦的伟大忍耐而骄傲，另一个人因极其乐于服从，另一个人因在谦卑上超越他人。一个人因知识的渊博而受诱惑，另一个人因阅读的广泛，另一个人因守夜的长久。这疾病除了通过德行之外，并不试图伤害一个人；它借着那些寻求生命供应的事物，引入致死的阻碍。因为当人们渴望行走在圣洁和成全的道路上时，仇敌并不在他们行走的道路之外设置陷阱来欺骗他们，正如有福的达味所说的：\Quote{在我所行走的路上，他们为我设下了网罗；}意思是，就在我们行走的这条德行之路上，当我们向着\Quote{来自高天召叫的奖赏，}\BibleRef{斐 3:14}奋勇前进时，我们可能会因成功而得意，从而沉沦跌倒，灵魂的脚被虚荣的网罗缠住。结果，那些本可在与仇敌的争战中不败的人，竟被自己胜利的伟大所击败；或者（这是另一种欺骗），我们过度紧张那可能达到的自律限度，因身体的软弱而未能坚持我们的路程。\par

% structure: p0020 source=h2 target=subsection reason=relative-heading-rank; base=h2

\subsection{第七章}

\Emph{虚荣被克服后，如何比以往更激烈地重新起来战斗。}\par

一切恶习在被打败时都会变得软弱，每日愈加衰弱，随着时间的推移和地点的变换而减弱、平息；或者至少，因为它们与相反的德行不同，更容易被避开和逃避。但这一恶习在被击倒时，却比以往更激烈地重新崛起，投入战斗；当我们以为它已被消灭时，它却因死亡而变得更强。其他种类的恶习通常只攻击那些在战斗中被它们战胜的人；但这一恶习却只更加热切地追击它的征服者；它被抵抗得越彻底，就越猛烈地攻击那因战胜它而得意洋洋的人。我们仇敌的狡诈诡计正在于此：当他不能用敌人的武器击败基督的士兵时，他就用他自己的矛将其击倒。\par

% structure: p0023 source=h2 target=subsection reason=relative-heading-rank; base=h2

\subsection{第八章}

\Emph{如何不因旷野或岁月而平息虚荣。}\par

正如我们所说，其他恶习有时会因处境的有利而平息，当罪的材料、机会和条件消失时，它们就会松懈并减少；但这一恶习却伴随着逃避它的人进入旷野，无法在任何地方被排除，外在的材料被移除时它也不会消失。因为它仅仅是受到它所攻击之人的德行的成就所鼓励。因为所有其他恶习，如上所述，有时会因时间的流逝而减弱并消失：对于这一恶习，除非有熟练的勤奋和谨慎的判断力支持，否则长寿并非障碍，反而为虚荣提供了新的燃料。\par

% structure: p0026 source=h2 target=subsection reason=relative-heading-rank; base=h2

\subsection{第九章}

\Emph{虚荣因与德行交织而更为危险。}\par

最后，其他完全不同于其相反德行的情感，公开地、仿佛在光天化日之下攻击我们，更容易被克服和防备；但这一恶习与我们的德行交织在一起，纠缠于战斗中，仿佛在黑夜的掩护下作战，更加危险地欺骗那些毫无防备、不警惕的人。\par

% structure: p0029 source=h2 target=subsection reason=relative-heading-rank; base=h2

\subsection{第十章}

\Emph{例子：希则克雅王如何被虚荣之箭击败。}\par

因为我们读到，犹大王希则克雅，在一切事上最完美的义人，且为圣经的见证所认可的人，在无数的德行赞词之后，被一支虚荣之箭击倒。他凭一次祈祷就能使亚述军队的十八万五千人丧命，天使在一夜间将其消灭，但这位君王却被骄傲和虚荣所战胜。关于他——略过那长长的德行清单，若要详述将费时甚久——我只说这一件事。他是一位在其生命终结已定、死亡之日已由主的判决确定之后，凭一次祈祷得以将他生命的期限延长十五年的人，太阳后退了十步，而它原本已在其轨道上向西方运行了，其后退驱散了那些已随着其轨迹的影子所标记的线条，由此，整个大地在同一日经历了两个白天，这是一个违反固定不变的自然规律的惊人奇迹。\EditorialNote{列王纪下第二十章}然而，在如此伟大而不可思议的奇迹之后，在如此众多的善良证据之后，请听圣经如何讲述他是如何被自己的成功所摧毁。\Quote{在那些日子，}我们被告知，\Quote{希则克雅病危，他向主祈祷，主垂听了并给了他一个征兆，}即我们在列王纪下所读到的，由依撒意亚先知通过太阳的后退所给的标记。\Quote{但是，}它说，\Quote{他没有按所得的恩惠回报，因为他的心高傲了；因此愤怒降在他身上，也降在犹大和耶路撒冷；后来他谦卑了，因为他的心高傲过，他和耶路撒冷的居民都谦卑了，因此主的愤怒没有在希则克雅的日子临到他们。}\BibleRef{编下 32:24}虚荣的疾病是何等危险、何等可怕！如此多的善行、如此多的德行、信德和虔敬，足以改变自然本身和整个世界的法则，却因一次骄傲的行为而消亡！以至于他所有的善行都被遗忘，仿佛从未存在过，若不是他通过恢复谦卑而平息了主的愤怒，他立刻就会受到主的愤怒之罚：这样，那个因骄傲的暗示而从如此崇高的卓越高度堕落的人，只能通过同样的谦卑阶梯重新登上他失去的高度。你想看到另一个类似败落的例子吗？\par

% structure: p0032 source=h2 target=subsection reason=relative-heading-rank; base=h2

\subsection{第十一章}

\Emph{乌齐雅王同样被同一疾病感染而败亡的例子。}\par

关于我们刚才所说的这位君王的祖先乌齐雅，他本身也受到圣经见证的完全称赞，在对其德行的极大赞誉之后，在他因虔敬和信德的功劳而获得的无数胜利之后，请听他是如何被虚荣的骄傲所推翻的。\Quote{并且，}我们被告知，\Quote{乌齐雅的名声远扬，因为主帮助他，增强了他。但他刚强起来，心就高傲，以致败坏，他背弃了主他的天主。}\BibleRef{编下 26:15}你看到另一个极其可怕的败落的例子，看到两位如此正直杰出的人是如何因他们的胜利和凯旋而堕落的。由此你可见，顺利的成功通常是多么危险，以至于那些不能被逆境伤害的人，若不谨慎，就会因顺境而毁灭；那些在战斗冲突中免于死亡危险的人，却在他们的战利品和凯旋面前倒下。\par

% structure: p0035 source=h2 target=subsection reason=relative-heading-rank; base=h2

\subsection{第十二章}

\Emph{反对虚荣的几处证言。}\par

宗徒也警告我们：\BibleQuote{不可贪图虚荣。}\BibleRef{迦 5:26} 主责备法利塞人时也说：\BibleQuote{你们彼此领受荣耀，却不寻求从唯一天主而来的荣耀，怎能信呢？}\BibleRef{若 5:44} 有福的达味也以威胁的口吻论及这些人：\Quote{因为天主已经驱散了那些取悦人者的骨头。}\par

% structure: p0038 source=h2 target=subsection reason=relative-heading-rank; base=h2

\subsection{第十三章}

\Emph{论虚荣攻击隐修士的方式}\par

对于初学者和那些在智力或知识上进展甚微的人，虚荣也通常使他们的心骄傲，或是因为他们嗓音优美，唱歌动听，或是因为身体消瘦，或是因为身材匀称，或是因为有富有而高贵的亲属，或是因为他们蔑视了军旅生活和荣誉。有时它也说服一个人，如果他留在世俗中，本可轻易获得荣誉和财富——而这些或许根本不可能得到——并以对不确定之事的虚幻希望使他自大；对于那些他从未拥有过的事物，却以骄傲和虚荣使他膨胀，仿佛他是一位已然蔑视它们的人。\par

% structure: p0041 source=h2 target=subsection reason=relative-heading-rank; base=h2

\subsection{第十四章}

\Emph{论它如何诱导人寻求领受圣秩}\par

但有时它使人渴望领受圣秩，渴求司铎职或执事职。它表示，如果一个人即使违背自己的意愿接受了这一职务，他也将如此圣洁和严谨地履行它，以至于能甚至在其它司铎中树立圣德的榜样；而且他不仅以自己的生活方式，还以他的教导和讲道赢得许多人。它甚至使一个人独自坐在斗室中时，也在心智和想象中周游到他人的住所和隐修院，并在虚幻狂喜的诱因下促使许多人悔改。\par

% structure: p0044 source=h2 target=subsection reason=relative-heading-rank; base=h2

\subsection{第十五章}

\Emph{论虚荣如何使心智沉醉}\par

可怜的灵魂受到这种虚荣的影响——仿佛被深深的睡眠所迷惑——以至于常常被这种思想的快感所牵引，并被这样的想象所充满，以至于甚至不能看见眼前的事物或弟兄，却喜欢沉浸在这些事情中，以游荡的心思做着白日梦，仿佛它们是真的。\par

% structure: p0047 source=h2 target=subsection reason=relative-heading-rank; base=h2

\subsection{第十六章}

\Emph{论一位长上突然到访，发现某人在斗室中被虚幻的虚荣所迷惑}\par

我记得一位长老，当我住在斯刻特旷野时，他曾去拜访某位弟兄的斗室。当他到达门口时，听见里面有人在喃喃自语，便稍作停留，想知道他正在诵读圣经或背诵什么（这是惯常的做法）而同时他正在工作。这位最出色的窃听者勤勉地侧耳倾听，带着一些好奇仔细聆听，发现那人被这种精神的攻击所诱导，幻想自己正向人群发表激昂的讲道。当长老站着，听他结束讲道，又回到他的职务，像执事一样宣告慕道者离席时，他终于敲了门。那人出来，以惯常的敬意迎接长老，领他进去，因知道长老已察觉他的心思而不安，问他何时到达，唯恐他在门口站得太久而受到伤害。老人愉快地开玩笑说：\Quote{我正好在你宣告慕道者离席时才到。}\par

% structure: p0050 source=h2 target=subsection reason=relative-heading-rank; base=h2

\subsection{第十七章}

\Emph{论除非发现过错的根源和原因，否则过错无法治愈}\par

\Emph{我}以为，将这些东西插入我这小小的作品中是好的，好让我们不仅藉理性，也藉实例学习关于诱惑的力量和伤害可怜灵魂的罪的次序，从而更谨慎地避免仇敌的陷阱和多种欺骗。因为埃及的教父们不加区分地提出这些东西，以那些仍在承受它们之人的口吻讲述，以揭露和展示他们实际遭受的以及与一切恶习的斗争，以及年轻者必定会遭受的；这样，当他们解释源于所有情欲的幻觉时，那些初学且心志热忱的人可以知道他们斗争的奥秘，并如同在镜中看见它们，从而既学习困扰他们的罪的根源，也学习对治的方法，并预先获知未来争斗的来临，得以教导他们应如何防范，或如何迎接它们并与之战斗。如同聪明的医生不仅习惯于治愈已有的疾病，也以明智的技巧力求防范未来的疾病，并以他们的处方和治疗药剂预防它们；同样，这些真正的灵魂医生，藉着灵性交谈，如同某种天上的解药，预先摧毁将要出现的灵魂疾病，不容许它们在新进者的心中立足，因为他们向他们揭示了威胁他们的情欲的根源和医治它们的补救方法。\par

% structure: p0053 source=h2 target=subsection reason=relative-heading-rank; base=h2

\subsection{第十八章}

\Emph{论隐修士应如何避远妇女和主教}\par

\Emph{因此}，这是教父们的一句古老格言，至今仍在流传——虽然我本人不能毫不羞愧地提出它，因为我未能避开自己的妹妹，也未能逃脱主教的手——即，隐修士务必要逃避妇女和主教。因为两者都不会允许一个曾经与他们亲近往来的人再关心斗室的宁静，或通过洞察神圣事物而继续以纯洁的双眼进行神圣默观。\par

% structure: p0056 source=h2 target=subsection reason=relative-heading-rank; base=h2

\subsection{第十九章}

\Emph{论我们可以克服虚荣的补救方法}\par

因此，那位渴望在真正属灵战斗中合规争战的\Person{基督}的运动员，应尽力制服这个变化多端的畸形怪物，它如同多种邪恶从各方面攻击我们，我们能通过这样的补救措施逃脱，即：思虑\Person{达味}那句话：\Quote{主分散了那些取悦于人者的骨头。}首先，我们绝不可允许自己在虚荣的暗示下做任何事，或为了获取虚浮的荣耀而做任何事。其次，当我们善始善终地开始一件事后，应当以同样的谨慎努力维持它，以防之后虚荣的疾患潜入，使我们一切劳作的果实化为乌有。凡在弟兄们的共同生活中用处或价值甚微的事，我们都应避免，因为会引向夸耀；凡会使我们在他人中显得突出，并因此博得人间的称赞，好像只有我们才能做到的事，都该避之。因为这些迹象表明，致命的虚荣污秽紧附我们；我们极易摆脱它，只需思虑：我们不仅会丧失那些在虚荣暗示下所行劳作的果实，更要犯下重大的罪，并如不敬之人受永罚，因为我们本该为天主而行的事，却为了取悦人而做，从而得罪了天主；且被那洞察一切秘密的天主定罪，我们竟将人置于天主之上，将世俗的赞美置于主的赞美之上。\par

\SourceRecord{Translated by C.S. Gibson.；Nicene and Post-Nicene Fathers, Second Series；Vol. 11.；Edited by Philip Schaff and Henry Wace.；Buffalo, NY: Christian Literature Publishing Co.,；1894.；<http://www.newadvent.org/fathers/350711.htm>.}

% structure: p0001 source=h1 target=section reason=page-title

\section{《隐修院规章》（第十二卷）}

% structure: p0002 source=h2 target=subsection reason=relative-heading-rank; base=h2

\subsection{第一章}

\Emph{我们第八场战斗如何是对抗骄傲之灵，以及它的特性。}\par

我们第八场且是最后一场战斗，是对抗骄傲之灵。这恶虽在我们与过错的斗争中最晚出现，在清单上也位列最后，但在起源与时间顺序上却是最先的：这是一头最凶残、比先前一切更可怕得多的恶兽，主要试探那些成全的人，并以可怖的咬噬吞吃那些几乎已达到德行圆满的人。\par

% structure: p0005 source=h2 target=subsection reason=relative-heading-rank; base=h2

\subsection{第二章}

\Emph{骄傲有两种。}\par

这骄傲有两种：一种，我们说过是困扰最优秀之人与属灵之人的；另一种，则连初学者与属血气之人也会攻击。虽然每种骄傲皆因危险的得意而涉及天主及人，但前一种更特别与天主有关；后一种则尤其指向人。关于后一种的起源与疗方，我们将靠天主的助佑，尽可能在本卷后半部分讨论。现在我们打算就前一种说几句话，如前所述，它特别试探那些成全的人。\par

% structure: p0008 source=h2 target=subsection reason=relative-heading-rank; base=h2

\subsection{第三章}

\Emph{骄傲如何同样毁灭一切德行。}\par

因此，没有别的过犯能像骄傲这恶那样毁灭一切德行，剥夺并劫掠一个人所有的义德与圣洁。它如同某种瘟疫般攻击整个人，不满足于只损伤一部分或一个肢体，而以致命的毒害败坏整个身体，并力图将那些已登临诸德之树顶端的人，以最致命的坠落推倒并毁灭。因为每一种别的过犯都满足于自身的界限与范围，虽然也遮蔽其他德行，但主要只针对一种，并特别攻击它。所以（为使我的意思更清楚），饕餮，即肚腹的贪欲与口舌的享乐，毁灭严格的节德；淫欲玷污贞洁；忿怒摧毁忍耐：以至于有时一个被某罪捆绑的人，并非完全缺乏其他德行；只是单纯被那在斗争中屈服于其对立恶习的德行所剥夺，却能在某种程度上保全其他德行。但这骄傲一旦占据了某个不幸的灵魂，就像最残暴的暴君，当诸德的高耸堡垒被攻破后，它彻底摧毁并踏平整座城邑，将圣善往昔的高墙夷为平地，与恶习混杂，不许那受其辖制的灵魂此后存有任何自由的痕迹。而且，它起初越富有，如今奴役的轭就越沉重，通过它残酷的蹂躏，它将剥去它所征服的灵魂的一切德能。\par

% structure: p0011 source=h2 target=subsection reason=relative-heading-rank; base=h2

\subsection{第四章}

\Emph{骄傲如何使路济弗尔从总领天使变为魔鬼。}\par

为使我们明白它可怕暴政的力量，我们看到那位天使，因他的光辉与美丽之大而被称作路济弗尔，除这骄傲外没有别的罪，就被逐出天庭，被骄傲之箭刺穿，从崇高显赫的天使高位坠入地狱。既然单是心中的骄傲就足以使如此伟大、满有如此大能之装饰的能力从天坠地，那么这坠落之巨大正告诉我们，我们这些被血肉软弱所包围的人应当怎样谨慎提防。但是，我们可以学习如何避免这恶的最致命毒害，如果我们追溯他坠落的原因与根源。因为软弱永不能被治愈，除非首先通过明智的查究来探明其起源与原因。因为当他（即路济弗尔）被赋予天主的光辉，并因造物主的慷慨而在其他更高权能中闪耀时，他相信自己获得那智慧的光辉与那些能力的美——这些本是造物主恩赐的礼物——是出于自己本性的能力，而非祂慈善的惠赐。因此他被高傲所充满，仿佛他在保持这种纯洁状态上不需要天主助佑，并自视如同天主，好像像天主一样不需要任何人，且信赖自己意志的力量，幻想通过它他能丰裕地自供一切为德行之圆满或完美真福之永续所需之物。仅此思想就是他最初坠落的原因。因此他被天主——他幻想不再需要的那位——所舍弃，突然变得不稳与动摇，发现了自己本性的软弱，并丧失了曾因天主恩赐而享有的福乐。因为他\Quote{喜爱毁灭之言}（他说过\Quote{我要升到天上}），和\Quote{诡诈的舌头}（他对自己说\Quote{我要与至高者同等}，\BibleRef{依 14:13}；又对亚当与厄娃说\Quote{你们将如同天主}），因此\Quote{天主必永远消灭他，拔除他，从他帐幕中把他拔出，从活人之地将他的根除去}。于是\Quote{义人见其毁灭，必害怕，且要嘲笑他说}（这话也最可公正地指向那些自信不靠天主的保护与助佑就能获得至善之人）：\Quote{看哪，那人不以天主为他的帮助者，却依赖财富的丰裕，在自己的虚妄中显赫。}\par

% structure: p0014 source=h2 target=subsection reason=relative-heading-rank; base=h2

\subsection{第五章}

\Emph{一切罪的唆使皆源于骄傲。}\par

这就是最初坠落的原因，以及原始疾病的起点。这疾病又通过那已被它毁灭者渗入第一个人，产生了所有过犯的软弱与质料。因为他相信凭自己自由意志和自身努力就能获得神性的荣耀，结果他恰恰失去了那本已通过造物主白白恩赐而拥有的荣耀。\par

% structure: p0017 source=h2 target=subsection reason=relative-heading-rank; base=h2

\subsection{第六章}

\Emph{骄傲之罪在战斗的实际顺序中居后，但在时间与起源上却居先。}\par

因此，藉由圣经中的实例和见证，已极为清楚地确立：骄傲之祸虽然在后来的战斗顺序中出现，但在起源上却是更早的，且是一切罪恶与过犯的开端；它（不像其他恶习）不仅单纯致命于与之相反的德行（此处即谦逊），同时也摧毁\Emph{一切}德行：它不只诱惑普通人和小人物，更主要的是那些已立于勇德高峰的人。因为先知论及此灵说：\Quote{他的食物是精选的。}因此，真福\Person{达味}虽然以极大的谨慎守护他内心的深处，以致敢于向那其良心秘密无法隐藏的那一位说：\Quote{主，我的心不骄矜，我的眼不高傲；我未曾行伟大超卓的事，也未曾行超过我能力的奇事；若我没有谦逊自持，}又说：\Quote{行事骄傲的人不得在我家中居住，}然而，由于他知道即使对成全者也难以保持这种警醒，他并未倚赖自己的努力，而是向天主祈祷并恳求祂的帮助，以求能不受这仇敌的箭矢所伤，他说：\Quote{不要让骄傲的脚来到我这里，}因为他害怕并畏惧陷入那论及骄傲者的话，即：\BibleQuote{天主拒绝骄傲的人；}\BibleRef{雅 4:6} 又说：\Quote{凡心高气傲的人，在主前是不洁的。}\par

% structure: p0020 source=h2 target=subsection reason=relative-heading-rank; base=h2

\subsection{第七章。}

\Emph{骄傲之恶如此巨大，以至它理应有天主亲自作为它的对手。}\par

骄傲之恶何等巨大，以至它理应有天使或别的德行作其对手，而唯有天主亲自作为它的对手！因为应当注意，对于陷入其他罪过的人，从未说过天主抵抗他们；我的意思是，从未说过天主反对\Quote{贪饕者、淫乱者、易怒者或贪婪者}，而只反对\Quote{骄傲者}。因为那些罪过只影响犯罪者自身，或似乎只对共享者（即他人）犯下；但这骄傲罪更直接地与天主有关，因此它理应有天主作为其对手，这尤为恰当。\par

% structure: p0023 source=h2 target=subsection reason=relative-heading-rank; base=h2

\subsection{第八章。}

\Emph{天主如何以谦逊之德摧毁了魔鬼的骄傲，以及证明此事的几处经文。}\par

因此，天主，万有的创造者与医治者，深知骄傲是邪恶的起因与泉源，便谨慎地以相反者医治相反者，使那些因骄傲而毁坏的事物因谦逊而得以恢复。因为一个说：\BibleQuote{我要升到天上；}\BibleRef{依 14:13} 另一个说：\Quote{我的灵魂被降低到了尘埃。}一个说：\Quote{我要与至高者相同；}另一个说：\BibleQuote{他虽具有天主的形体，却空虚了自己，取了奴仆的形体，并贬抑自己，听命至死。}\BibleRef{斐 2:6} 一个说：\Quote{我要高举我的宝座在天主的星辰之上；}另一个说：\BibleQuote{跟我学吧！因为我是良善心谦的。}\BibleRef{玛 11:29} 一个说：\BibleQuote{我不认识主，也不放以色列人走；}\BibleRef{出 5:2} 另一个说：\BibleQuote{如果我说我不认识他，我就像你们一样是个说谎者；但是我认识他，也遵守他的话。}\BibleRef{若 8:55} 一个说：\Quote{我的河流是我的，是我命它们存在的；}另一个说：\Quote{我什么也不能做，而是住在我内的父，他做这些工作。}一个说：\BibleQuote{这一切王国及其荣华都是我的，我愿意给谁，就给谁；}\BibleRef{路 4:6} 另一个说：\BibleQuote{他本是富有的，为了你们却成了贫困的，好使你们因着他的贫困而成为富有的。}\BibleRef{格后 8:9} 一个说：\BibleQuote{如同人收集遗下的卵，我收集了整个大地；没有一只翅膀扇动，张口或发出最轻微的声音；}\BibleRef{依 10:14} 另一个说：\Quote{我成了荒野中的鹈鹕；我警醒，成了屋顶上孤单的麻雀。}一个说：\BibleQuote{我用脚掌踏干了所有被围困的河流；}\BibleRef{依 37:25} 另一个说：\BibleQuote{我不能求我父，他立刻就会给我十二军以上的天使吗？}\BibleRef{玛 26:53} 如果我们思考我们原始堕落的原因和救恩的根基，并考虑后者由谁并以何种方式奠定，前者如何起源，我们就可以藉着魔鬼的堕落或基督的榜样，学会如何避免如此可怕的骄傲之死。\par

% structure: p0026 source=h2 target=subsection reason=relative-heading-rank; base=h2

\subsection{第九章。}

\Emph{我们如何也能克服骄傲。}\par

因此，我们能逃脱这最邪恶之灵的陷阱，只要我们在所取得进步的每一种德行上，说出宗徒的这话：\Quote{不是我，而是天主的恩宠偕同我；}以及\BibleQuote{因天主的恩宠，我成了今日的我；}\BibleRef{格前 15:10} 以及\BibleQuote{是天主在你们内工作，使你们愿意，并付诸实行，为成就他的善意。}\BibleRef{斐 2:13} 正如我们救恩的创始者亲自所说：\BibleQuote{那住在我内的，我也住在他内的，他就结许多果实；因为离了我，你们什么也不能做。}\BibleRef{若 15:5} 以及\Quote{若不是主建造房屋，建造的人也是徒然劳苦；若不是主看守城池，看守的人也是徒然警醒。}以及\Quote{你们清早起来，也是徒然。}因为\BibleQuote{不在于那愿意的，也不在于那奔跑的，而在于那施行仁慈的天主。}\BibleRef{罗 9:16}\par

% structure: p0029 source=h2 target=subsection reason=relative-heading-rank; base=h2

\subsection{第十章。}

\Emph{无人能单凭自身力量获得完美的德行和所应许的真福。}\par

因为任何人的意愿和追求，无论多么热切和焦虑，当他仍被那与灵魂争战的肉体所困时，都不足以获得如此伟大的完美奖赏和正直纯洁的棕榈枝，除非他得到天主的慈悲保护，从而有幸达到他所切望并奔跑追求的目标。因为\BibleQuote{一切美好的赠与和完美的恩赐都是从上面来的，从光明之父降下来的。}\BibleRef{雅 1:17}\BibleQuote{你有什么不是领受的呢？既然是领受的，为什么你夸耀好像不是领受的呢？}\BibleRef{格前 4:7}\par

% structure: p0032 source=h2 target=subsection reason=relative-heading-rank; base=h2

\subsection{第十一章。}

\Emph{关于盗贼和达味的例子，以及我们的蒙召，以说明天主的恩宠。}\par

因为如果我们回想那个因一次宣认而被接入乐园的盗贼\BibleRef{路 23:40}，我们就会感到他获得如此福乐并非凭生平的功绩，而是靠仁慈天主的恩赐。或者如果我们想到达味王那两个严重可耻的罪，因补赎的一句话而被涂抹了\BibleRef{撒下 12:13}，我们就会看到，这里也不是他工作的功绩足以获得如此大罪的宽恕，而是天主的恩宠格外丰富，当真心补赎的机会被抓住时，祂藉着仅仅一句话的完全告解，除去了所有罪的重担。如果我们再思考人类蒙召与得救的开端，如宗徒所说，我们得救并非由于自己，也不是由于我们的工作，而是由于天主的恩赐和恩宠，我们就能清楚看见，整个完美\BibleQuote{不在于那愿意的，也不在于那奔跑的，而在于那施予慈悲的天主}，祂使我们战胜我们的过犯，而我们方面没有任何工作和生命的功绩足以补偿，也没有任何意志的努力足以攀登完美的高峰，或制服我们不得不使用的肉体：因为此身体的任何痛苦和心灵的任何痛悔，都不足以仅凭人的努力（即没有天主的帮助）获得内心真正贞洁所需的东西，以赢得纯洁这一伟大美德（这唯天使所固有，且如同本地于天上）：因为一切善行的施行都源于祂的恩宠，祂藉着倍增祂的恩赐，赐予我们软弱的意志和短暂渺小的人生历程如此持久的福乐和巨大的荣耀。\par

% structure: p0035 source=h2 target=subsection reason=relative-heading-rank; base=h2

\subsection{第十二章。}

\Emph{没有任何劳苦堪与所许诺的福乐相比。}\par

因为今世所有悠长岁月，如果考虑到将来荣耀的永恒，都消失无踪；我们所有忧愁，在默想那巨大福乐时，都烟消云散，归于无有，如灰烬不再可见。\par

% structure: p0038 source=h2 target=subsection reason=relative-heading-rank; base=h2

\subsection{第十三章。}

\Emph{长老们关于获得纯洁之方法的教导。}\par

\Emph{因此}，现在是用他们流传下来的原话提出教父们的意见的时候了；即那些并非用高调言辞描绘完美之路及其特征，而是实事求是，以自己经验和可靠榜样拥有它并将其传下来的人。他们如此说：没有人能完全洗净肉体的罪，除非他意识到自己一切劳苦和努力都不足以达到如此伟大完美的目标；并且，他并非由传承的系统所教导，而是由自己的感受、德行和经验所教导，认识到惟有靠天主的慈悲和帮助才能获得。因为要获得如此辉煌崇高的纯洁与完美的奖赏，无论加上多少斋戒、守夜、阅读、独居和隐退的努力，都不足以凭实际努力和劳苦的功绩来确保获得。因为人自己的努力和人为的辛劳永远无法弥补天主恩赐的缺乏，除非因天主的慈悲应允他的祈祷而赐予。\par

% structure: p0041 source=h2 target=subsection reason=relative-heading-rank; base=h2

\subsection{第十四章。}

\Emph{天主的助佑赐给那些劳苦的人。}\par

我这样说并不是轻视人的努力，也不是企图劝阻任何人停止工作和尽己所能。但我清楚而极其认真地提出，不是我的意见，而是长老们的意见：没有这些，完美不可能获得，但仅仅依靠这些而没有天主的恩宠，谁也不能达到。因为当我们说人的努力没有天主的帮助不能自行保证完美时，我们是在强调天主的慈悲和恩宠只赐给那些劳苦努力的人，并且（用宗徒的话）赐给那些\BibleQuote{愿意}和\BibleQuote{奔跑}的人，正如\BibleRef{咏 88}中以天主位格所唱的：\BibleQuote{我把救助赐给一位有能者，从我的百姓中提拔了一位被选者。}因为我们按照救主的话说：求的，就给他们；叩门的，就给他们开门；寻找的，就使他们找到\BibleRef{玛 7:7}。但是，我们的求、寻找和叩门是不够的，除非天主的慈悲赐予我们所求的，为我们打开所叩的，使我们能够找到所寻找的。因为祂很乐意赐予所有这些，只要我们以善意给祂机会。因为祂比我们自己更渴望并期待我们的完美和得救。蒙福的达味深知靠自己的努力无法保证工作的增长，所以他以更新的祈祷恳求从主获得他工作的\BibleQuote{引导}，说：\BibleQuote{愿我们手的工作在我们身上得到引导；是的，我们手的工作，求祢引导。}又说：\BibleQuote{天主，求祢坚固在我们身上所成就的。}\par

% structure: p0044 source=h2 target=subsection reason=relative-heading-rank; base=h2

\subsection{第十五章。}

\Emph{我们可以从谁学习完美之路。}\par

因此，如果我们确实渴望借由实际行为与真理获得德行的冠冕，就应当听从那些导师和向导，他们不是用浮夸的演说做梦，而是借由行动和经验学习，也能同样教导我们，指引我们，并为我们展示一条最确实的道路，使我们能抵达那里；他们也证明自己到达那里是凭信德而非自己努力的任何功行。而且，他们所获得的心灵纯洁尤其教导他们这一点：即越来越认识到自己背负着罪（因为他们的痛悔随着灵魂纯洁的进步而与日俱增），并不断从心底叹息，因为他们看到自己无法避免那些根深蒂固的过错的污点和瑕疵，这些过错源于思绪的无数琐碎。因此，他们宣称，他们期待的来生的赏报，并非来自自己功行的功劳，而是来自主的慈悲；他们不因自己在内心上的重大谨慎而自夸胜过别人，因为他们将此归功于天主的恩宠，而非自己的努力；他们也不因那些冷淡、比自己更差的人的疏忽而自满，反而力求持久的谦逊，注视那些他们知道真正无罪、已在天国享受永福的人，借此思虑避免骄傲的堕落，同时总是看到自己的目标与需要哀悼的事：因为他们知道，在肉身的重负下，无法达到所渴望的心灵纯洁。\par

% structure: p0047 source=h2 target=subsection reason=relative-heading-rank; base=h2

\subsection{第十六章}

我们若不蒙天主的慈悲和灵感，连努力追求成全都不可能。\par

因此，我们应当依照他们的教导和训诲，如此努力追求，勤于斋戒、守夜、祈祷，以及心灵和身体的痛悔，以免这一切因这疾病的攻击而变得无用。因为我们不仅应当相信，无法凭自己的努力和刻苦获得这实际的成全，而且也无法为了成全而实践的那些事——即我们的努力、刻苦和渴望——没有天主保护之助和祂的灵感、惩戒与劝勉之恩宠，这些恩宠祂通常通过别人的工具或亲自来探望我们，倾注在我们心中。\par

% structure: p0050 source=h2 target=subsection reason=relative-heading-rank; base=h2

\subsection{第十七章}

各种章节清楚表明，没有天主的助佑，我们无法做出任何与我们的救恩相关的事。\par

最后，我们救恩的作者教导我们，我们不仅应当思考，也应当在一切所行的事上承认：\Quote{我凭自己什么也不能做，但那住在我内的父，祂在做这些工作。}祂是就自己所取的人性说的，说自己凭自己什么也不能做；我们这些灰土，竟以为在关乎我们救恩的事上不需要天主的帮助吗？因此，让我们在一切事上，在感到自己软弱的同时也感到祂的帮助，学习同圣人们一起宣告：\Quote{我几乎跌倒，但主扶持了我。主是我的力量和赞美，祂成了我的救恩。}又说：\Quote{若非主帮助了我，我的灵魂几乎已居于阴府。若我说：我的脚滑了；主啊，你的仁慈扶持了我。我心中愁烦众多，你的安慰却使我的灵魂欢乐。}又因我们的心在主内并在忍耐中得以坚固，让我们说：\Quote{主成了我的保护者，并领我到宽阔之处。}又知道知识因工作的进展而增长，让我们说：\Quote{主啊，你点亮了我的灯；我的天主，你照亮我的黑暗，因为藉着你我将脱离诱惑，藉着我的天主我将跳过城墙。}然后，感觉到我们已为自己求得勇气和忍耐，且更容易、毫不费力地被导向德行的道路，让我们说：\Quote{是天主以力量束我腰，使我的道路完全；祂使我的腿快如母鹿的腿，并使我在高处屹立；祂教导我的手作战。}又既已获得明辨，借以能击倒仇敌，让我们向天主大声呼喊：\Quote{你的管教使我达到终点，你的管教必教导我。你在我脚下开阔我的步伐，我的脚没有软弱。}因为我这样藉着你的知识和力量得以坚强，我将勇敢地接续以下的话语，说：\Quote{我要追赶仇敌，并赶上他们；不将他们消灭，决不回转。我将击碎他们，他们不能站起；他们必倒在我的脚下。}又想到自己的软弱，以及仍背着软弱的肉身，若无祂的帮助，就无法战胜如此凶恶的罪恶仇敌，让我们说：\Quote{藉着你，我们将击散仇敌；凭你的名，我们将践踏那起来攻击我们的人。因为我不倚靠我的弓，我的剑也不能拯救我。因为你救了我们脱离那压迫我们的人，并羞辱了那恨我们的人。}此外：\Quote{你以力量引导我作战，并使那起来攻击我的屈服于我。你又使我的仇敌向我转身逃跑，并毁灭了那恨我的人。}又思量单凭自己的武器不能得胜，让我们说：\Quote{拿起盾牌与甲胄，起来帮助我。抽出剑来，为那追逼我的挡住前路；向我的灵魂说：我是你的救恩。}你使我的膀臂如铜弓。你又以你救恩的盾牌护卫了我，你的右手扶持了我。\Quote{因为我们的祖先并非靠自己的剑取得那地，也不是凭自己的膀臂得救；而是你的右手、你的膀臂和你脸上的光亮，因你喜悦他们。}最后，当我们以忧心忡忡的心感激地省思祂所有的恩惠时，让我们以内心最深处的感情向祂呼喊这一切，因为我们作战，并从那获得了知识和自律与明辨的光照，因为祂用自己的武器装备了我们，并以德行之带加强了我们，因为祂使我们的仇敌转身而逃，并赐予我们能力，将他们如风前的尘土驱散：\Quote{我爱你，主，我的力量；主是我的堡垒、我的避难所、我的拯救者。我的天主是我的帮助者，我必信赖祂。我的保护者，我救恩的角，我的支持。我要赞美呼求主的名，我必从仇敌中得救。}\par

% structure: p0053 source=h2 target=subsection reason=relative-heading-rank; base=h2

\subsection{第十八章}

\Emph{我们如何不但因本性状况，也因祂每日的眷顾而受天主的恩宠保护。}\par

不仅为祂创造我们为理性受造物、赋予我们自由意志的能力、祝福我们领受圣洗的恩宠、赐予我们法律的知识和帮助而感谢祂，也为祂每日的眷顾所赐予我们的这些事物感谢祂：即祂拯救我们脱离仇敌的诡计；祂与我们合作，使我们能战胜肉身的罪恶；甚至在我们不知不觉中，祂庇护我们免于危险；祂保护我们免于堕入罪恶；祂帮助我们并光照我们，使我们能理解并认识祂所给予我们的实际帮助（有些人以为这就是法律的意义）；当我们因祂的影响而暗中为自己的罪过和疏忽感到痛悔时，祂以祂的眷顾垂顾我们，并以有益于我们灵魂的方式惩戒我们；甚至违背我们的意愿，祂有时也引领我们归向救恩；最后，这个更容易倾向罪恶的自由意志本身，被祂转向更好的目的，并在祂的催促和暗示下，转向德行的道路。\par

% structure: p0056 source=h2 target=subsection reason=relative-heading-rank; base=h2

\subsection{第十九章}

\Emph{这有关天主恩宠的信德如何由古代圣祖传授给我们。}\par

这就是对天主的谦卑，这就是古代教父们那真正的信德，这信德在他们继承者中仍完好无损。对这信德，宗徒的诸德——他们常以此彰显——不仅在我们中间，也在教外人和不信者中间，作了不容置疑的见证：因为他们存心纯朴，持守渔夫们纯朴的信德，并未以世俗精神通过辩证三段论或\Person{Cicero}的雄辩来领受它，而是藉纯洁生活的经验、无瑕的行为、纠正过失，而（说得更确切）藉可见的证明，学会了圆满的品格在于那信德；没有它，无法确保对天主的虔敬、罪的除净、生活的改正，或德行的圆满。\par

% structure: p0059 source=h2 target=subsection reason=relative-heading-rank; base=h2

\subsection{第二十章}

\Emph{论一人因亵渎而被交给最不洁之灵}\par

\Emph{我认识}弟兄中的一位，我从心底希望从未认识他；因为后来他允许自己承担了我这修会的责任。他向一位最可敬的长老告解，说自己受到一种可怕肉欲的袭击：他被一种难忍的贪欲所燃烧，带着一种不自然的、宁愿受苦而非行可耻之事的欲望。那位长老如同一名真正的属灵医生，立刻看穿了这邪恶的内在原因和根源。他深深地叹息说：\Quote{主绝不会容许你被交给如此污秽的灵，除非你曾亵渎祂。}这事被识破后，他立刻俯伏在地，在他脚前，极其惊愕，仿佛看见自己内心的秘密被天主揭穿，承认他曾以恶念亵渎了天主之子。由此可见，被骄傲之灵附身的人，或犯过亵渎天主之罪的人——如同对那必须祈求纯洁恩赐者行不义——便被剥夺了他的正直与成全，不配得到贞洁的圣化恩宠。\par

% structure: p0062 source=h2 target=subsection reason=relative-heading-rank; base=h2

\subsection{第二十一章}

\Emph{犹大王约阿施的例子，显示其骄傲的后果。}\par

某些类似之事我们在\BibleBook{}{编年纪}中读到。因为犹大王\Person{约阿施}七岁时被司铎\Person{耶何耶大}召来继承王位，并且圣经见证他的一切行为，只要上述司铎还在世，便都受称赞。但请听圣经在\Person{耶何耶大}死后关于他的记载，以及他如何骄傲自大，被交给最可耻的状态。\Quote{耶何耶大死后，众首领进去朝拜王，王被他们的侍奉所安抚，就听从了他们。他们离弃了主、他们祖先的天主的圣殿，事奉了木偶和偶像，因此大发烈怒临到\Place{犹大}和\Place{耶路撒冷}。}过了一会儿：\Quote{过了一年，\Place{叙利亚}的军兵上来攻击他，到了\Place{犹大}和\Place{耶路撒冷}，杀了民中的众首领，将所掠的财物都送到\Place{大马士革}王那里。虽然来到的\Place{叙利亚}军兵很少，主却将极多的人交在他们手中，因为他们离弃了主、他们祖先的天主；他们在\Person{约阿施}身上执行了可耻的审判，并且离开时使他患了重病。}你看，骄傲的后果就是他交给可憎又污秽的情欲。因为凡骄傲自大，并容许自己被当作天主敬拜的人，正如宗徒所说的，\Quote{交给可耻的情欲和邪僻的心，去做那些不合理的事。}并且，因为圣经说：\Quote{凡心里骄傲的，在天主面前为不洁。}凡心中充满骄傲膨胀的人，就被交给最可耻的混乱，受它迷惑，这样当受羞辱时，他才能知道自己因肉体的不洁和知道不洁欲望而成为不洁——这件事他在骄傲的心中曾拒绝承认；而且，肉体可耻的感染可能揭露他心中隐秘的不洁，这污秽是他因骄傲之罪而染上的，并且通过身体明显的玷污，证明他是不洁的，他先前没有看见自己因灵里的骄傲而成了不洁。\par

% structure: p0065 source=h2 target=subsection reason=relative-heading-rank; base=h2

\subsection{第二十二章}

\Emph{每一个骄傲的灵魂都屈服于属灵的邪恶，被其欺骗。}\par

这清楚表明，凡被骄傲的膨胀所占据的灵魂，都被交给灵魂的叙利亚人，即属灵的邪恶，并陷入肉体的情欲；这样，灵魂最终因世上的过犯而谦卑，并因肉体而受玷污，就能认清自己的不洁——尽管当它心中冷漠地昂然站立时，无法理解因心中的骄傲在天主眼中已成为不洁；而藉此谦卑下来，人就能摆脱先前的冷漠，被降卑、被肉欲的羞耻所困惑，从而从此更急切地奔向灵性的热忱与温暖。\par

% structure: p0068 source=h2 target=subsection reason=relative-heading-rank; base=h2

\subsection{第二十三章}

\Emph{如何惟藉谦卑之德才能达到成全。}\par

因此清楚表明，除非藉着真正的谦卑——人首先向弟兄们表现出来，也在内心深处向天主显示，相信自己若没有天主时刻伸出的保护与帮助，绝不可能获得他所渴望、所迫切奔赴的成全——否则无人能达到成全与纯洁的终极。\par

% structure: p0071 source=h2 target=subsection reason=relative-heading-rank; base=h2

\subsection{第二十四章}

\Emph{谁受属灵骄傲的攻击，谁受属肉体骄傲的攻击。}\par

以上这些，靠着天主的助佑，尽我们微薄的能力，已经足够论述关于属灵的骄傲了，我们说过，这种骄傲攻击进步的基督徒。这种骄傲不为大多数人所熟悉或经历，因为多数人不以追求心灵的完全纯洁为目标，以达到这些争战的阶段；他们也没有从先前我们已在各卷中解释了性质与补救方法的过犯中获得任何净化。但它通常只攻击那些已经克服了先前过犯、并几乎已到达德行之巅的人。因为我们那最狡猾的仇敌不能通过肉体的跌倒摧毁他们，就企图通过属灵的灾难使他们跌倒并推翻他们，试图借此夺走他们以极大辛劳获得的先前赏报的代价。至于我们，仍然纠缠于世俗的情欲，他从不屑于以这种方式试探我们，而是以一种更粗俗的、即我所说的肉性的骄傲来推翻我们。因此，我认为，如我所承诺的，有必要稍微谈谈这种骄傲，即通常影响我们以及我们这类人，尤其危及年轻人和初学者的心灵的骄傲。\par

% structure: p0074 source=h2 target=subsection reason=relative-heading-rank; base=h2

\subsection{第二十五章}

\Emph{描述肉性的骄傲，以及它在隐修士灵魂中所产生的邪恶。}\par

这种肉性的骄傲，我们谈到的，当它进入一个隐修士的心时，这心只是冷淡的，并且在弃绝世界上有了一个坏的开端，它不容许他从先前世俗的高傲状态屈从于基督的真正谦卑，而是首先使他成为不顺从和粗暴的；然后不允许他温和与慈善；也不容许他与弟兄们平等和相似；不允许他按照天主和我们的救主所命令的那样，被剥夺和舍弃他的世俗财物；并且，虽然弃绝世界无非是死灭和十字架的标志，并且不能从其他任何基础开始或升起，除了这些基础：即一个人应认识到他不仅在精神上对这个世界的行为是死的，而且也应意识到他必须在身体上死去——它反而使他期望长寿，在他面前提出许多漫长的疾病，并使他充满羞愧和混乱。如果他在被剥夺一切后开始依靠他人的而非自己的财物生活，它说服他，根据那段经文（如前所述，那些因这种迟钝和心冷而变得愚钝的人不可能理解），由自己的财物而非他人的来提供食物和衣服要好得多：\BibleQuote{施比受更为有福。}\BibleRef{宗 20:35}\par

% structure: p0077 source=h2 target=subsection reason=relative-heading-rank; base=h2

\subsection{第二十六章}

\Emph{一个人若基础不好，就会日益变坏。}\par

那些被这种不信的心所占据，并且因魔鬼自身的不信而从起初皈依时似乎被点燃的信德火花中坠落的人，开始更焦虑地看守他们先前开始施舍的钱财，并以更大的贪婪囤积它，如同那些一旦浪费就无法再收回的人；或者——更糟的是——取回他们先前丢弃的东西；或者（第三种最可恶的罪），收集他们以前从未拥有过的东西，因此被证实他们在弃绝世界上只不过取了隐修士的名字和样式而已。因此，以这个开端，在这个坏而腐朽的基础上，整个罪恶的上层建筑必然层层加增，在这种恶劣的基础上不可能建造任何东西，除非是将可怜的灵魂带入毫无希望的坍塌。\par

% structure: p0080 source=h2 target=subsection reason=relative-heading-rank; base=h2

\subsection{第二十七章}

\Emph{描述从骄傲的邪恶中产生的各种过错。}\par

如此情绪使心灵刚硬，始于这种可悲的冷漠，必然每日每况愈下，并以更可怕的结局终结生命。它沉溺于先前的欲望，并被宗徒所说的不虔敬的贪婪所战胜（如他所说：\Quote{贪财，即是偶像崇拜，或拜偶像}，又说\Quote{贪爱金钱}，他说，\Quote{是万恶之根}），因而永不能将基督真实无伪的谦卑接纳心中。人若自恃出身高贵，或因在世上的地位（身体虽已离弃，心却未离）而自高，或以其财富自傲而自取灭亡，便不再甘愿承受隐修院的轭，或听从任何长老的训导，不仅拒绝遵守任何服从或听命的规则，甚至连关于成全的教导也不愿听。如此，对属灵谈论的厌恶在他心中滋长，以至于若发生这样的谈话，他无法定睛一处，目光茫然地四处游移，眼睛左右转动，如同常人之态。他不作有益的咳嗽，却从干涩的喉咙吐痰；无故故意咳嗽，用手指敲击、捻动，像写字般乱画；四肢不停地躁动，以至于属灵谈话进行时，你会觉得他坐在荆棘上，而且是极尖锐的荆棘，或坐在一堆蠕虫之中。若谈话单纯地转向有益听者的事，他便认为那是针对他而提。在属灵生活的探讨进行时，他被自己的猜疑之心占据，不注意吸取有益之处，反而焦虑地探究为何说这话，或暗自思忖如何反驳，以致完全无法领受那些极好地提出的道理，也不能从中获益。因此，结果便是属灵的谈话不仅对他无益，反而有害，成为他更大犯罪的机会。因为他良心不安，觉得一切都在针对他，便更加固执地硬着心，更激烈地感受到愤怒的刺痛。随后，他声音高亢，言语苛刻，回答尖刻喧嚷，步态倨傲专横；舌头过于轻率，在交谈中冒失，除非心中对弟兄怀有苦毒，否则不爱沉默，而他的沉默并非表示痛悔或谦卑，而是骄傲与愤怒。因此，几乎很难说他身上哪一点更令人反感：是那无节制、喧闹的欢乐，还是这可怖而致命的严肃。因为在前者，我们见到不合时宜的饶舌、轻浮无聊的笑、无节制无纪律的喜乐；在后者，则是充满愤怒与致命的沉默，单纯源于希望尽可能延长对某位弟兄怀有的怨恨，而这怨恨在沉默中滋长，并非为了从中获得谦卑与忍耐的德行。正如那受激情支配的人轻易使他人痛苦，且羞于向被他得罪的弟兄道歉，同样，当那弟兄主动向他道歉时，他轻蔑拒绝。他不仅不为弟兄的谦卑举动所触动或软化，反而更加愤怒，因为弟兄在谦卑上超过了他。那种有益的谦卑与道歉，通常能结束魔鬼的诱惑，对他却成了更严重爆发的机会。\par

% structure: p0083 source=h2 target=subsection reason=relative-heading-rank; base=h2

\subsection{第二十八章。}

\Emph{论某位兄弟的骄傲。}\par

我在这地区曾听闻一事，回想起来仍战栗羞愧：有一位初学生，因院长责备他表现出放弃那在弃绝世俗时短暂体验过的谦卑，并被魔鬼的骄傲所膨胀，竟极其无礼地回答：\Quote{我一时谦卑，难道是为了永远服从吗？}对于他这放肆而邪恶的回答，长老完全惊愕，无言以对，仿佛这回答是来自老路济弗尔本人，而非来自人；以致他无法开口反驳如此无耻之言，只能从心中发出叹息与呻吟，心中默默思量关于我们救主的话：\Quote{祂虽具有天主的形体，却谦卑自抑，听命至死}——并非如那被魔鬼骄傲之灵所附者所说的\Quote{一时}，而是\Quote{至死}。\BibleRef{斐 2:6}\par

% structure: p0086 source=h2 target=subsection reason=relative-heading-rank; base=h2

\subsection{第二十九章。}

\Emph{藉以识别灵魂中肉身骄傲的征兆。}\par

为简要总结这种骄傲的论述，我们尽可能收集其一些外在标记，通过外在行为向渴求成全教导的人传达其特性。我认为最好用几句话来阐述，使我们可以方便地识别这些标记，从而发现并察觉它；这样，当这种情欲的根源被揭露、暴露，并通过亲眼见证被看到和追踪时，便更容易被拔除和避免。因为只有那时，这最致命的恶才能完全逃脱；如果我们不等到它已经控制我们时才开始防备其危险的热度和有害影响，而是在事先（可以这样说）识别其症状，并以明智和谨慎的远见加以预防，我们才能逃脱。因为正如我们之前所说，从外在的举止可以判断一个人的内在状态。通过这些标记，我们先前所讲的肉性骄傲便显露出来。首先，在交谈中，声音响亮；沉默时，带有苦毒；欢笑时，笑声喧闹过度；严肃时，却无理地阴沉；回答中带有怨恨；言语过于随便，说话不加斟酌。他完全缺乏忍耐，没有爱德；敢于侮辱他人，自己却难以忍受冒犯；在服从方面令人厌烦，除非职责与自己的意愿和喜好相符；不接受劝告时毫不宽容；放弃自己意愿时软弱；固执己见，不愿为他人让步；结果，尽管他不能给出明智的建议，但在一切事上总是偏爱自己的意见胜过长者的意见。\par

% structure: p0089 source=h2 target=subsection reason=relative-heading-rank; base=h2

\subsection{第三十章}

\Emph{当人因骄傲而冷淡时，如何想要管辖他人}\par

当被骄傲控制的人通过这些堕落阶段跌落后，他对隐修院的纪律感到恐惧，仿佛弟兄的陪伴妨碍了他的成全，他人的罪恶阻碍了他忍耐与谦卑的进步，于是他渴望独居一室；或者急于建造隐修院，召集他人来教导和训导，似乎他能使更多人受益，使自己从坏门徒变成更坏的师傅。因为当人因心中的骄傲陷入这种最危险、最有害的冷淡时，他既不能成为真正的隐修士，也不能成为世俗之人；更糟的是，他竟许诺通过这种可怜的状态和生活方式获得成全。\par

% structure: p0092 source=h2 target=subsection reason=relative-heading-rank; base=h2

\subsection{第三十一章}

\Emph{如何克服骄傲并获得成全}\par

\Emph{因此}，若我们希望建筑的顶端完美无缺，并悦乐天主，就当努力按照福音严谨的规则奠定基础，而非随从自己情欲的私愿；这基础只能是敬畏天主和谦卑，源自内心的良善与纯朴。然而，谦卑若不舍弃一切就不能获得；只要人对此陌生，就不可能获得服从的德行、忍耐的力量、良善的平安或爱德的圆满；没有这些，我们的心就不可能成为圣神的居所，正如主通过先知所说：\BibleQuote{我的灵将安息在谁身上？岂非在谦卑、安静、听我话的人身上吗？}或按照精确表达希伯来文的版本：\BibleQuote{我要垂顾谁？岂非贫穷、卑微、心灵痛悔、因我言语战栗的人吗？}\par

% structure: p0095 source=h2 target=subsection reason=relative-heading-rank; base=h2

\subsection{第三十二章}

\Emph{如此毁灭一切德行的骄傲，如何能被真正的谦卑摧毁}\par

\Emph{因此}，在属灵争战中按规矩奋斗、并渴望被主加冕的基督运动员，应当尽一切努力摧毁这毁灭一切德行的最凶猛野兽，深知只要它还盘踞在心中，他不仅永远无法摆脱各种邪恶，甚至看似有的善行也会因它的恶性影响而丧失。因为，德行的任何架构（可以这样说）都绝不可能在我们灵魂中建立起来，除非先在心中打下真正谦卑的根基：这根基牢固地奠定后，才能承受圆满与爱德的重负，使得我们，如我们所说，首先从心底向弟兄们展示真正的谦卑，绝不认同使他们忧伤或受伤害；而除非真正的自我舍弃——即剥夺并放弃我们一切所有——因基督的爱植根于我们内，我们根本不可能做到。其次，必须以纯朴之心毫无伪饰地负起服从与顺命的轭，除院长命令外，我们内再无自己的意愿。只有认为自己不仅向世界死去，而且愚拙无知，并毫无异议地执行长上所命的一切，相信是来自天上并受命于天的人，才能确保这一点。\par

% structure: p0098 source=h2 target=subsection reason=relative-heading-rank; base=h2

\subsection{第三十三章}

\Emph{对抗骄傲之恶的补救方法}\par

当人处于这种状态时，毫无疑问，这种安宁而稳妥的谦卑状态将会随之而来，以至于我们视自己低于他人，并以极大的忍耐承受一切加于我们的事，即使是有害的、令人忧伤的、有损害的——如同这些来自比我们优越的人。我们不仅会极其轻松地承受这些事，而且会视它们为微不足道和无关紧要，只要我们不断牢记我们的主和祂所有圣人的受难：考虑到我们所受的伤害远不及他们的，如同我们远落后于他们的功德和生命；同时记住我们不久将离开这个世界，并很快通过我们在此生迅速终结而分享他们的命运。因为这样的考量不仅是骄傲的终结，也是一切罪恶的确定终结。然后，在此之后，我们必须牢牢把握这种对天主的谦卑：我们应当如此确保，不仅承认没有祂的帮助和恩宠，我们不可能完成任何与达到完美德行相关的事，而且真正相信我们能够理解这一点本身就是祂的恩赐。\par

\SourceRecord{Translated by C.S. Gibson.；Nicene and Post-Nicene Fathers, Second Series；Vol. 11.；Edited by Philip Schaff and Henry Wace.；Buffalo, NY: Christian Literature Publishing Co.,；1894.；<http://www.newadvent.org/fathers/350712.htm>.}
